\documentclass[italian]{unipd-notes}

\unipdsetup{
  course = {Analisi Matematica 1},
  short-course = {Analisi Matematica 1},
  course-code = {IN10100190},
  channel = {B},
  professor = {Prof. Alessandro Languasco},
  academic-year = {2026--2027},
  degree-year = {1},
  semester = {1},
  document-type = {Appunti delle lezioni},
  author = {Simone Siega},
  date = {29 settembre 2026},
  version = {0.1.0}
}

\setcounter{tocdepth}{1}

\crefname{chapter}{lezione}{lezioni}
\Crefname{chapter}{Lezione}{Lezioni}
\newcourseterm{naturali}{\ensuremath{\N}}{Insieme dei numeri naturali, con $0\in\N$}
\newcourseterm{interi}{\ensuremath{\Z}}{Insieme dei numeri interi}
\newcourseterm{razionali}{\ensuremath{\Q}}{Insieme dei numeri razionali}
\newcourseterm{reali}{\ensuremath{\R}}{Insieme dei numeri reali}
\newcourseterm{reali-estesi}{\ensuremath{\overline{\R}}}{Insieme dei numeri reali estesi}

\newcourseterm{esiste-unico}{\ensuremath{\exists!}}{Esiste ed è unico}
\newcourseterm{inclusione}{\ensuremath{\subseteq}}{Inclusione tra insiemi, con uguaglianza ammessa}

\newcourseterm{parte-intera-inferiore}{\ensuremath{\lfloor x\rfloor}}{Parte intera inferiore di $x$}
\newcourseterm{parte-intera-superiore}{\ensuremath{\lceil x\rceil}}{Parte intera superiore di $x$}


\begin{document}
\makecoursefrontmatter
\makecoursemainmatter

\printnoidxglossary[title={Convenzioni e notazione},sort=def]

\chapter[Lezione 1 --- 29 settembre 2026]{Lezione 1: numeri naturali, ordinamento e principio di induzione}
\label{chap:lezione-01}

\section{Numeri naturali e le loro operazioni}
\label{sec:lez01-numeri-naturali}

Si considera l'insieme dei numeri naturali
\[
  \N=\{0,1,2,3,\ldots\},
\]
adottando la convenzione \(0\in\N\).

Su \(\N\) si considerano le operazioni di somma e prodotto, entrambe
\emph{operazioni binarie}, ossia applicazioni che associano a ogni
coppia ordinata di numeri naturali un numero naturale:
\[
  +\colon\N\times\N\longrightarrow\N,
  \qquad (a,b)\longmapsto a+b,
\]
\[
  \cdot\colon\N\times\N\longrightarrow\N,
  \qquad (a,b)\longmapsto ab.
\]
La notazione \(\N\times\N\) indica le coppie ordinate
di naturali (\cref{def:lez06-prodotto-cartesiano}).

Le proprietà fondamentali delle due operazioni vengono assunte come
\emph{assiomi}, senza dimostrazione.

La somma è \emph{commutativa} (l'ordine degli addendi è irrilevante)
e \emph{associativa} (lo è il raggruppamento): per ogni
\(a,b,c\in\N\) si ha
\[
  a+b=b+a,
  \qquad
  (a+b)+c=a+(b+c).
\]
Ammette inoltre un unico \emph{elemento neutro} \(e_0\in\N\),
caratterizzato dalla proprietà
\[
  \forall a\in\N,\qquad a+e_0=a.
\]
Tale elemento è \(e_0=0\).

Anche il prodotto è commutativo e associativo: per ogni
\(a,b,c\in\N\) vale
\[
  ab=ba,
  \qquad
  (ab)c=a(bc).
\]
Il suo unico \emph{elemento neutro} \(e_1\in\N\) soddisfa
\[
  \forall a\in\N,\qquad ae_1=a,
\]
ed è \(e_1=1\).

Il prodotto è infine \emph{distributivo} rispetto alla somma:
per ogni \(a,b,c\in\N\),
\[
  \begin{aligned}
    a(b+c) &= ab+ac,\\
    (a+b)c &= ac+bc.
  \end{aligned}
\]

\clearpage

\section{Ordinamento totale dei numeri naturali}
\label{sec:lez01-ordine-naturali}

\begin{definition}[Ordine naturale]
\label{def:ordine-naturali}
Per \(a,b\in\N\), si pone
\[
  a\le b
  \quad\Longleftrightarrow\quad
  \exists c\in\N:\ a+c=b.
\]
\end{definition}

La relazione della \cref{def:ordine-naturali} soddisfa le seguenti
proprietà.

\emph{Riflessività.}
Per ogni \(a\in\N\),
\[
  a\le a.
\]
Infatti, nella \cref{def:ordine-naturali} è sufficiente
scegliere \(c=0\), poiché \(a+0=a\).

\emph{Antisimmetria.}
Per ogni \(a,b\in\N\),
\[
  a\le b
  \quad\text{e}\quad
  b\le a
  \quad\Longrightarrow\quad
  a=b.
\]

\emph{Transitività.}
Per ogni \(a,b,c\in\N\),
\[
  a\le b
  \quad\text{e}\quad
  b\le c
  \quad\Longrightarrow\quad
  a\le c.
\]

\emph{Dicotomia.}
Per ogni \(a,b\in\N\),
\[
  a\le b
  \quad\text{oppure}\quad
  b\le a.
\]
Le due possibilità non sono esclusive quando \(a=b\),
poiché in tal caso valgono entrambe le disuguaglianze.

Queste proprietà caratterizzano \(\le\) come un \emph{ordinamento
totale} su \(\N\). L'ordine è inoltre compatibile con le
operazioni di somma e prodotto.

\begin{proposition}[Compatibilità dell'ordine]
\label{prop:compatibilita-ordine}
Per ogni \(a,b,c\in\N\), se \(a\le b\), allora
\[
  a+c\le b+c.
\]
Inoltre, se \(c>0\), vale
\[
  ac\le bc.
\]
\end{proposition}

Per \(c=0\) vale comunque \(ac=bc=0\).
La dimostrazione della \cref{prop:compatibilita-ordine}
non viene sviluppata in questa sede.

\section{Richiami su insiemi e dimostrazioni}
\label{sec:lez01-insiemi-dimostrazioni}

Siano \(A\) e \(B\) due insiemi. La loro \emph{intersezione}
raccoglie gli elementi comuni, mentre l'\emph{unione} comprende
gli elementi che appartengono ad almeno uno dei due:
\[
  A\cap B=\{x:x\in A\text{ e }x\in B\},
  \qquad
  A\cup B=\{x:x\in A\text{ oppure }x\in B\}.
\]

\clearpage

\begin{figure}[htbp]
  \centering
  \begin{tikzpicture}[unipd/diagram]
  \begin{scope}
    \clip (0,0) circle (1);
    \fill[unipd-accent!25] (1.2,0) circle (1);
  \end{scope}
  \draw[unipd/edge] (0,0) circle (1);
  \draw[unipd/edge] (1.2,0) circle (1);
  \node at (-0.45,0) {$A$};
  \node at (1.65,0) {$B$};
  \node at (0.6,-1.35) {$A\cap B$};
  \begin{scope}[xshift=5cm]
    \fill[unipd-accent!25] (0,0) circle (1);
    \fill[unipd-accent!25] (1.2,0) circle (1);
    \draw[unipd/edge] (0,0) circle (1);
    \draw[unipd/edge] (1.2,0) circle (1);
    \node at (-0.45,0) {$A$};
    \node at (1.65,0) {$B$};
    \node at (0.6,-1.35) {$A\cup B$};
  \end{scope}
\end{tikzpicture}

  \caption{Intersezione (a sinistra) e unione (a destra)
  di due insiemi.}
  \label{fig:venn-operazioni}
\end{figure}

Negli enunciati matematici si distinguono le \emph{ipotesi},
che si assumono valide, e la \emph{tesi}, che si intende dimostrare.
Un teorema può essere formulato nella forma
\[
  \text{IPOTESI}\Longrightarrow\text{TESI}.
\]
Per due proposizioni \(P\) e \(Q\), l'implicazione
\(P\Longrightarrow Q\) afferma che dalla verità di \(P\)
segue quella di \(Q\). Essa è logicamente equivalente alla
propria \emph{contronominale}:
\[
  P\Longrightarrow Q
  \quad\Longleftrightarrow\quad
  \neg Q\Longrightarrow\neg P.
\]

Nella \emph{dimostrazione per assurdo} si assumono le ipotesi e
la negazione della tesi, ricavandone una contraddizione, ossia
una proposizione \(C\) insieme alla sua negazione \(\neg C\).

\section{Principio di induzione}
\label{sec:lez01-induzione}

Il principio di induzione consente di dimostrare una proprietà
per tutti i naturali a partire da un valore iniziale \(p\).

\begin{theorem}[Principio di induzione]
\label{thm:principio-induzione}
Siano \(A\subseteq\N\) e \(p\in\N\). Se
\[
  p\in A
  \qquad\text{e}\qquad
  \forall n\in\N,\quad n\in A\Longrightarrow n+1\in A,
\]
allora
\[
  \{n\in\N:n\ge p\}\subseteq A.
\]
\end{theorem}

Le due ipotesi sono dette rispettivamente \emph{passo base}
e \emph{passo induttivo}.

Per dimostrare che una proprietà \(P(n)\) vale per ogni
\(n\ge p\), si considera l'insieme
\[
  A:=\{n\in\N:n\ge p\text{ e }P(n)\text{ è vera}\}.
\]
L'applicazione del \cref{thm:principio-induzione} richiede
di verificare \(P(p)\) e di dimostrare, per un generico
\(n\ge p\), che
\[
  P(n)\Longrightarrow P(n+1).
\]

Si considerano due applicazioni del principio: la somma dei
primi naturali positivi e la somma geometrica finita.
Nel secondo esempio compare l'insieme dei numeri reali \(\R\),
che sarà introdotto nella \cref{sec:lez03-numeri-reali}.

\begin{example}[Somma dei primi \(n\) naturali positivi]
\label{ex:somma-primi-naturali}
Per ogni \(n\in\N\) con \(n\ge1\), vale
\[
  \sum_{i=1}^{n}i=\frac{n(n+1)}2.
\]
\end{example}

\begin{proof}
Si applica il \cref{thm:principio-induzione} con valore
iniziale \(p=1\), considerando l'insieme
\[
  A:=\left\{n\in\N:n\ge1\text{ e }
    \sum_{i=1}^{n}i=\frac{n(n+1)}2\right\}.
\]

\emph{Passo base.}
Per \(n=1\) si ha
\[
  \sum_{i=1}^{1}i=1=\frac{1(1+1)}2,
\]
quindi \(1\in A\).

\emph{Passo induttivo.}
Sia \(n\ge1\) e si supponga che valga l'ipotesi induttiva
\[
  \sum_{i=1}^{n}i=\frac{n(n+1)}2.
\]
Separando l'ultimo addendo e applicando l'ipotesi induttiva,
si ottiene
\begin{align*}
  \sum_{i=1}^{n+1}i
    &=\sum_{i=1}^{n}i+(n+1)\\
    &=\frac{n(n+1)}2+(n+1)\\
    &=\frac{n(n+1)+2(n+1)}2\\
    &=\frac{(n+1)(n+2)}2.
\end{align*}
Pertanto \(n\in A\Longrightarrow n+1\in A\).

Per il \cref{thm:principio-induzione}, si conclude che
\(A=\{n\in\N:n\ge1\}\), ossia che la formula
vale per ogni \(n\ge1\).
\end{proof}

\begin{example}[Somma geometrica finita]
\label{ex:somma-geometrica-finita}
Sia \(q\in\R\setminus\{1\}\). Per ogni \(n\in\N\) vale
\[
  \sum_{k=0}^{n}q^k=\frac{1-q^{n+1}}{1-q}.
\]
Si adotta la convenzione \(q^0=1\) anche per \(q=0\);
in particolare, \(\sum_{k=0}^n0^k=1\) e la formula resta valida.
Per \(q=1\), escluso perché annulla il denominatore, si ha
\[
  \sum_{k=0}^{n}1^k=n+1.
\]
\end{example}

\begin{proof}
Fissato \(q\in\R\setminus\{1\}\), si applica il
\cref{thm:principio-induzione} con valore iniziale \(p=0\),
considerando l'insieme
\[
  A:=\left\{n\in\N:
    \sum_{k=0}^{n}q^k=\frac{1-q^{n+1}}{1-q}\right\}.
\]

\emph{Passo base.}
Per \(n=0\), poiché \(q\ne1\), si ha
\[
  \sum_{k=0}^{0}q^k=q^0=1
  =\frac{1-q}{1-q},
\]
quindi \(0\in A\).

\emph{Passo induttivo.}
Sia \(n\in\N\) e si supponga che valga l'ipotesi induttiva
\[
  \sum_{k=0}^{n}q^k=\frac{1-q^{n+1}}{1-q}.
\]
Separando l'ultimo addendo e applicando l'ipotesi induttiva,
si ottiene
\begin{align*}
  \sum_{k=0}^{n+1}q^k
    &=\sum_{k=0}^{n}q^k+q^{n+1}\\
    &=\frac{1-q^{n+1}}{1-q}+q^{n+1}\\
    &=\frac{1-q^{n+1}+q^{n+1}(1-q)}{1-q}\\
    &=\frac{1-q^{n+2}}{1-q}.
\end{align*}
Pertanto \(n\in A\Longrightarrow n+1\in A\).

Per il \cref{thm:principio-induzione}, si conclude che
\(A=\N\), ossia che la formula vale per ogni \(n\in\N\).
Poiché \(q\in\R\setminus\{1\}\) era arbitrario,
essa vale per ogni \(q\in\R\setminus\{1\}\).
\end{proof}

\chapter[Lezione 2 --- 30 settembre 2026]{Lezione 2: numeri interi, razionali e fattoriale}
\label{chap:lezione-02}

\section{Numeri interi}
\label{sec:lez02-numeri-interi}

Dato un insieme numerico \(A\), si indica con \(-A\) l'insieme degli
opposti dei suoi elementi:
\[
  -A:=\{-a:a\in A\}.
\]
Poiché \(0\in\N\), l'insieme \(-\N\) contiene lo zero e gli opposti
dei naturali positivi. L'unione di \(\N\) e \(-\N\) permette quindi
di introdurre i numeri interi.

\begin{definition}[Numeri interi]
\label{def:numeri-interi}
L'insieme dei \emph{numeri interi} è
\[
  \Z:=\N\cup(-\N).
\]
\end{definition}

Si assume l'\emph{assioma di esistenza e unicità dell'opposto}: per ogni
\(m\in\Z\) esiste un unico \(n\in\Z\) tale che
\[
  m+n=0.
\]
Tale elemento si indica con \(-m\) ed è detto \emph{opposto} di \(m\).

L'ordinamento dei naturali, introdotto nella
\cref{sec:lez01-ordine-naturali}, si estende agli interi. La
moltiplicazione per un intero positivo conserva il verso delle
disuguaglianze, mentre la moltiplicazione per un intero negativo
lo inverte.

\begin{proposition}[Compatibilità dell'ordine in \(\Z\)]
\label{prop:compatibilita-ordine-interi}
Per ogni \(a,b,c\in\Z\), se \(a\le b\), allora
\[
  c>0\Longrightarrow ac\le bc,
  \qquad
  c<0\Longrightarrow ac\ge bc.
\]
\end{proposition}

Per \(c=0\) si ha invece \(ac=bc=0\). La dimostrazione della
\cref{prop:compatibilita-ordine-interi} non viene sviluppata
in questa sede.

Si consideri ora il problema di determinare, per \(a\in\Z\), un
intero \(b\) tale che
\[
  ab=1.
\]
Quando esiste, tale elemento è detto \emph{inverso moltiplicativo}
di \(a\) e si indica con \(a^{-1}\). Negli interi, tuttavia, non
ogni elemento ammette un inverso moltiplicativo.

\begin{proposition}[Inversi moltiplicativi in \(\Z\)]
\label{prop:inversi-interi}
Gli unici interi che ammettono un inverso moltiplicativo in
\(\Z\) sono \(-1\) e \(1\). In particolare,
\[
  (-1)^{-1}=-1,
  \qquad
  1^{-1}=1.
\]
\end{proposition}

\begin{proof}
Siano \(a,b\in\Z\) tali che \(ab=1\). Il prodotto è positivo,
quindi \(a\) e \(b\) hanno lo stesso segno.

Se entrambi sono positivi, si ha \(a,b\ge1\). Poiché \(ab=1\),
necessariamente \(a=b=1\). Se entrambi sono negativi, gli interi
positivi \(-a\) e \(-b\) soddisfano \((-a)(-b)=1\), da cui
\(a=b=-1\).

Viceversa, \(1\cdot1=1\) e \((-1)(-1)=1\), quindi entrambi gli
interi indicati ammettono un inverso moltiplicativo in \(\Z\).
\end{proof}

\section{Numeri razionali}
\label{sec:lez02-numeri-razionali}

La \cref{prop:inversi-interi} mostra che in \(\Z\) soltanto
\(-1\) e \(1\) ammettono un inverso moltiplicativo. Si introduce
pertanto un insieme numerico più ampio, nel quale sia possibile
rappresentare i rapporti tra interi con denominatore non nullo.

Per due insiemi \(A\) e \(B\), la \emph{differenza insiemistica}
è definita da
\[
  A\setminus B:=\{x\in A:x\notin B\}.
\]
In particolare, \(\N\setminus\{0\}\) è l'insieme dei naturali
strettamente positivi.

\begin{definition}[Numeri razionali]
\label{def:numeri-razionali}
L'insieme dei \emph{numeri razionali} è
\[
  \Q:=\left\{\frac{m}{n}:m\in\Z,\ n\in\N\setminus\{0\}\right\}.
\]
\end{definition}

Si ottiene così la catena di inclusioni
\[
  \N\subseteq\Z\subseteq\Q.
\]

Il denominatore è scelto positivo, riportando l'eventuale segno
negativo al numeratore. Uno stesso razionale può essere rappresentato
da frazioni diverse, per esempio \(\frac12=\frac24\). Riducendo
una frazione ai minimi termini si ottiene una rappresentazione
\[
  \frac{m}{n},\qquad \operatorname{mcd}(m,n)=1,
\]
dove il massimo comune divisore è quello positivo di \(|m|\) e \(n\).
Con \(n>0\), la rappresentazione ridotta è unica; in particolare,
lo zero si rappresenta mediante \(\frac01\).

In \(\Q\) si assume l'\emph{assioma dell'inverso moltiplicativo}:
\[
  \forall a\in\Q\setminus\{0\}\;\exists!\,b\in\Q:
  \quad ab=1.
\]
L'unico elemento \(b\) che soddisfa questa condizione si indica
con \(a^{-1}\) ed è detto anche \emph{reciproco} di \(a\). Per esempio,
\[
  2^{-1}=\frac12.
\]

\section[Irrazionalità della radice quadrata di 2]{Irrazionalità di \(\sqrt{2}\)}
\label{sec:lez02-irrazionalita-radice-due}

L'insieme dei razionali non è sufficiente a risolvere tutte le
equazioni polinomiali. Si consideri, per esempio, il polinomio
\[
  p(x)=x^2-2,
  \qquad p\in\Q[x],
\]
dove \(\Q[x]\) indica l'insieme dei polinomi nella variabile
\(x\) a coefficienti razionali. Si dimostrerà che l'equazione
\(x^2-2=0\) non ammette soluzioni razionali. Il risultato sarà
anche espresso nella forma \(\sqrt{2}\notin\Q\): l'esistenza
della radice quadrata reale non negativa di \(2\) sarà giustificata
nella \cref{sec:lez03-completezza-radice-due}.

Si richiamano anzitutto le nozioni di numero pari e dispari,
necessarie per la dimostrazione.

\clearpage

\begin{definition}[Numeri pari e dispari]
\label{def:lez02-pari-dispari}
Un intero \(m\) si dice \emph{pari} se esiste \(k\in\Z\) tale che
\[
  m=2k,
\]
mentre si dice \emph{dispari} se esiste \(k\in\Z\) tale che
\[
  m=2k+1.
\]
\end{definition}

Ogni intero è pari oppure dispari, e le due possibilità sono
mutuamente esclusive.

\begin{lemma}[Parità del quadrato]
\label{lem:quadrato-pari}
Per ogni \(m\in\Z\), se \(m^2\) è pari, allora \(m\) è pari.
\end{lemma}

\begin{proof}
Si dimostra la contronominale, secondo lo schema introdotto nella
\cref{sec:lez01-insiemi-dimostrazioni}. Sia \(m\) dispari. Per la
\cref{def:lez02-pari-dispari}, esiste \(k\in\Z\) tale che
\[
  m=2k+1.
\]
Pertanto
\begin{align*}
  m^2
    &=(2k+1)^2\\
    &=4k^2+4k+1\\
    &=2\bigl(2k(k+1)\bigr)+1.
\end{align*}
Poiché \(2k(k+1)\in\Z\), il quadrato \(m^2\) è dispari.
La contronominale è verificata e il lemma segue.
\end{proof}

\begin{theorem}[Irrazionalità di \(\sqrt{2}\)]
\label{thm:irrazionalita-radice-due}
L'equazione \(x^2-2=0\) non ammette soluzioni in \(\Q\).
Equivalentemente, con la notazione introdotta sopra,
\[
  \sqrt{2}\notin\Q.
\]
\end{theorem}

\begin{proof}
Si procede per assurdo. Si supponga che esista \(x\in\Q\) tale che
\[
  x^2=2.
\]
Scrivendo \(x\) in forma ridotta, esistono \(m\in\Z\) e
\(n\in\N\setminus\{0\}\) tali che
\[
  x=\frac{m}{n},
  \qquad
  \operatorname{mcd}(m,n)=1.
\]
Sostituendo la rappresentazione di \(x\) nell'uguaglianza
\(x^2=2\), si ottiene
\[
  \left(\frac{m}{n}\right)^2=2.
\]
Moltiplicando per \(n^2>0\), segue
\[
  m^2=2n^2.
\]
Di conseguenza, \(m^2\) è pari e, per il \cref{lem:quadrato-pari},
anche \(m\) è pari. Esiste quindi \(h\in\Z\) tale che \(m=2h\).
Sostituendo nell'uguaglianza precedente,
\[
  4h^2=2n^2
  \quad\Longrightarrow\quad
  n^2=2h^2.
\]
Anche \(n^2\) è pari e, applicando nuovamente il
\cref{lem:quadrato-pari}, si conclude che \(n\) è pari.

Gli interi \(m\) e \(n\) sono dunque entrambi divisibili per \(2\),
in contraddizione con \(\operatorname{mcd}(m,n)=1\).
L'ipotesi iniziale è falsa: l'equazione \(x^2-2=0\) non ammette
soluzioni razionali.
\end{proof}

Si richiama ora la nozione di numero primo per enunciare
una generalizzazione del risultato precedente.

\begin{definition}[Numero primo]
\label{def:numero-primo}
Un numero naturale \(p>1\) si dice \emph{primo} se i suoi unici
divisori positivi sono \(1\) e \(p\).
\end{definition}

Per esempio, \(2\) e \(3\) sono primi, mentre \(4\) non lo è,
poiché è divisibile anche per \(2\). Il numero \(1\) non è primo,
come stabilito dalla definizione.

\begin{proposition}[Irrazionalità della radice di un numero primo]
\label{prop:irrazionalita-radice-primo}
Sia \(p\in\N\) un numero primo. L'equazione
\[
  x^2-p=0
\]
non ammette soluzioni in \(\Q\). Indicando con \(\sqrt{p}\) la
radice quadrata reale non negativa di \(p\), il risultato si esprime
nella forma
\[
  \sqrt{p}\notin\Q.
\]
\end{proposition}

La dimostrazione della \cref{prop:irrazionalita-radice-primo} non viene
sviluppata in questa sede.

\section{Fattoriale}
\label{sec:lez02-fattoriale}

\begin{definition}[Fattoriale]
\label{def:fattoriale}
Il \emph{fattoriale} di un numero naturale è definito ricorsivamente
dalle condizioni
\[
  0!:=1,
  \qquad
  (n+1)!:=(n+1)n!,
  \quad n\in\N.
\]
\end{definition}

Applicando la \cref{def:fattoriale}, si ottiene
\begin{align*}
  0!&=1,\\
  1!&=1\cdot0!=1,\\
  2!&=2\cdot1!=2,\\
  3!&=3\cdot2!=3\cdot2\cdot1.
\end{align*}
Per \(n\ge1\), il fattoriale ammette anche la rappresentazione
\[
  n!=1\cdot2\cdot\ldots\cdot n=\prod_{i=1}^{n}i.
\]

\begin{proposition}[Una disuguaglianza per il fattoriale]
\label{prop:fattoriale-maggiore-n}
Per ogni \(n\in\N\) vale
\[
  n!\ge n.
\]
\end{proposition}

\begin{proof}
Si consideri l'insieme
\[
  A:=\{n\in\N:n!\ge n\}.
\]
Per \(n=0\), dalla \cref{def:fattoriale} si ha
\[
  0!=1\ge0,
\]
quindi \(0\in A\).

Si procede per induzione su \(n\ge1\), applicando il
\cref{thm:principio-induzione} con valore iniziale \(p=1\).

\emph{Passo base.}
Per \(n=1\), si ha
\[
  1!=1\ge1.
\]

\emph{Passo induttivo.}
Sia \(n\ge1\) e si supponga che valga l'ipotesi induttiva
\[
  n!\ge n.
\]
Per la \cref{def:fattoriale} e la compatibilità dell'ordine
rispetto alla moltiplicazione per un numero positivo
(\cref{prop:compatibilita-ordine}), si ottiene
\[
  (n+1)!=(n+1)n!\ge(n+1)n\ge n+1,
\]
poiché \(n+1>0\) e \(n\ge1\). Il passo induttivo è quindi
verificato.

Per il principio di induzione, la disuguaglianza vale per ogni
\(n\ge1\). Poiché è verificata anche per \(n=0\), si conclude
che \(A=\N\), e dunque la tesi.
\end{proof}

\chapter[Lezione 3 --- 1 ottobre 2026]{Lezione 3: coefficienti binomiali e numeri reali}
\label{chap:lezione-03}

\section{Coefficienti binomiali}
\label{sec:lez03-coefficienti-binomiali}

La nozione di fattoriale, introdotta nella \cref{def:fattoriale},
permette di definire i coefficienti binomiali per ogni coppia di
numeri naturali.

\begin{definition}[Coefficiente binomiale]
\label{def:coefficiente-binomiale}
Siano \(n,k\in\N\). Il \emph{coefficiente binomiale} è definito da
\[
  \binom{n}{k}
  :=
  \begin{cases}
    \displaystyle\frac{n!}{k!(n-k)!}, & \text{se }k\le n,\\[8pt]
    0, & \text{se }k>n.
  \end{cases}
\]
\end{definition}

Dalla \cref{def:coefficiente-binomiale} si ricavano alcuni casi
particolari. Per ogni \(n\in\N\), poiché \(0!=1\), si ha
\[
  \binom{n}{0}=\frac{n!}{0!\,n!}=1,
  \qquad
  \binom{n}{n}=\frac{n!}{n!\,0!}=1.
\]
Inoltre, per ogni \(n\in\N\) con \(n\ge1\),
\[
  \binom{n}{n-1}
  =\frac{n!}{(n-1)!\,(n-(n-1))!}
  =\frac{n!}{(n-1)!\,1!}
  =n.
\]

\begin{proposition}[Simmetria dei coefficienti binomiali]
\label{prop:simmetria-coefficienti-binomiali}
Per ogni \(n,k\in\N\) con \(k\le n\) vale
\[
  \binom{n}{k}=\binom{n}{n-k}.
\]
\end{proposition}

\begin{proof}
Per la \cref{def:coefficiente-binomiale} e la commutatività
della moltiplicazione, si ha
\[
  \binom{n}{n-k}
  =\frac{n!}{(n-k)!\,k!}
  =\frac{n!}{k!(n-k)!}
  =\binom{n}{k}.
\]
\end{proof}

Un'altra proprietà fondamentale è l'identità di Pascal, che esprime
un coefficiente binomiale come somma di due coefficienti della riga
precedente.

\begin{proposition}[Identità di Pascal]
\label{prop:identita-pascal}
Per ogni \(n,k\in\N\) con \(n\ge1\) e \(1\le k\le n\) vale
\[
  \binom{n}{k}
  =\binom{n-1}{k-1}+\binom{n-1}{k}.
\]
\end{proposition}

\begin{proof}
Si considera anzitutto il caso \(1\le k\le n-1\).
Per la \cref{def:fattoriale} si ha
\[
  k!=k(k-1)!,
  \qquad (n-k)!=(n-k)(n-k-1)!.
\]
Queste identità permettono di ridurre allo stesso denominatore
i due addendi della \cref{def:coefficiente-binomiale}:
\begin{align*}
  \binom{n-1}{k-1}+\binom{n-1}{k}
    &=\frac{(n-1)!}{(k-1)!(n-k)!}
      +\frac{(n-1)!}{k!(n-k-1)!}\\
    &=\frac{k(n-1)!}{k!(n-k)!}
      +\frac{(n-k)(n-1)!}{k!(n-k)!}\\
    &=\frac{\bigl(k+(n-k)\bigr)(n-1)!}{k!(n-k)!}\\
    &=\frac{n!}{k!(n-k)!}
     =\binom{n}{k}.
\end{align*}

Resta il caso \(k=n\). Per definizione,
\[
  \binom{n-1}{n}=0,
  \qquad
  \binom{n-1}{n-1}=\binom{n}{n}=1.
\]
Pertanto anche in questo caso
\[
  \binom{n}{n}
  =\binom{n-1}{n-1}+\binom{n-1}{n}.
\]
\end{proof}

Sebbene la \cref{def:coefficiente-binomiale} sia espressa mediante
un rapporto di fattoriali, il coefficiente binomiale è sempre un
numero naturale:
\[
  \binom{n}{k}\in\N
  \qquad\forall n,k\in\N.
\]
Questa proprietà viene qui enunciata senza dimostrazione.

\section{Numeri reali}
\label{sec:lez03-numeri-reali}

Gli insiemi numerici introdotti nelle lezioni precedenti soddisfano
le inclusioni
\[
  \N\subseteq\Z\subseteq\Q.
\]
Come mostrato nel \cref{thm:irrazionalita-radice-due},
l'equazione \(x^2=2\) non ammette soluzioni razionali.
Si considera quindi l'insieme dei numeri reali \(\R\),
più ampio di \(\Q\), per il quale vale la catena di inclusioni
\[
  \N\subseteq\Z\subseteq\Q\subseteq\R.
\]
L'esistenza in \(\R\) di un numero non negativo il cui quadrato
sia \(2\) sarà ripresa nella
\cref{sec:lez03-completezza-radice-due}.

Una proprietà fondamentale di \(\R\) è la completezza, che in
questa sede viene formulata mediante l'assioma di separazione.

Siano \(A,B\subseteq\R\) due insiemi non vuoti. Si dice che
\(A\) e \(B\) sono \emph{separati}, con \(A\) a sinistra di
\(B\), se
\[
  a\le b
  \qquad\forall a\in A,\ \forall b\in B.
\]

\phantomsection
\label{ass:assioma-separazione}
\textbf{Assioma di separazione.}
Se \(A,B\subseteq\R\) sono non vuoti e separati, allora esiste
almeno un numero reale \(\lambda\) tale che
\[
  a\le\lambda\le b
  \qquad\forall a\in A,\ \forall b\in B.
\]
Il numero \(\lambda\) si dice \emph{elemento separatore} di
\(A\) e \(B\).

L'assioma esprime una proprietà di completezza dei numeri reali.
Le altre formulazioni equivalenti non vengono considerate in
questa sede.

\section{Proprietà archimedea}
\label{sec:lez03-proprieta-archimedea}

L'\hyperref[ass:assioma-separazione]{assioma di separazione}
permette di dimostrare che l'insieme dei numeri naturali non è
superiormente limitato in \(\R\).

\begin{theorem}[Proprietà archimedea]
\label{thm:proprieta-archimedea}
Per ogni \(x\in\R\) esiste \(n\in\N\) tale che
\[
  n>x.
\]
\end{theorem}

\begin{proof}
Sia \(x\in\R\). Se \(x\le0\), è sufficiente scegliere
\(n=1\). Si consideri quindi \(x>0\) e si ponga
\[
  A:=\{n\in\N:n\le x\}.
\]
Poiché \(0\in A\), l'insieme \(A\) è non vuoto.
Se \(A\ne\N\), esiste \(n\in\N\setminus A\), per il quale
\(n>x\). Resta dunque da escludere il caso \(A=\N\).

Si supponga per assurdo che \(A=\N\) e si ponga
\[
  B:=\{y\in\R:n\le y\ \text{per ogni }n\in\N\}.
\]
Dall'ipotesi \(A=\N\) segue che \(n\le x\) per ogni
\(n\in\N\). Pertanto \(x\in B\), e anche \(B\) è non vuoto.
Inoltre, per definizione di \(B\),
\[
  n\le y
  \qquad\forall n\in A=\N,\ \forall y\in B.
\]
Gli insiemi \(A\) e \(B\) sono quindi separati. Per
l'\hyperref[ass:assioma-separazione]{assioma di separazione}
esiste \(\lambda\in\R\) tale che
\[
  n\le\lambda\le y
  \qquad\forall n\in\N,\ \forall y\in B.
\]

Poiché \(\lambda\) maggiora ogni naturale e
\(n+1\in\N\) per ogni \(n\in\N\), si ha
\[
  n+1\le\lambda
  \quad\Longrightarrow\quad
  n\le\lambda-1
  \qquad\forall n\in\N.
\]
Per definizione di \(B\), ciò implica \(\lambda-1\in B\).
D'altra parte, essendo \(\lambda\) un elemento separatore,
\(\lambda\le y\) per ogni \(y\in B\). Scegliendo
\(y=\lambda-1\), si ottiene
\[
  \lambda\le\lambda-1,
\]
una contraddizione. Pertanto \(A\ne\N\), ed esiste
\(n\in\N\setminus A\) tale che \(n>x\).
\end{proof}

\section{Parte intera, floor e ceiling}
\label{sec:lez03-parte-intera}

Il \cref{thm:proprieta-archimedea} conduce alla seguente
proprietà, che individua la posizione di ogni numero reale
rispetto a due interi consecutivi.

\begin{proposition}[Proprietà della parte intera]
\label{prop:parte-intera}
Per ogni \(x\in\R\) esiste un unico \(n\in\Z\) tale che
\[
  n\le x<n+1.
\]
\end{proposition}

La dimostrazione della \cref{prop:parte-intera} non viene
sviluppata in questa sede.

L'intero \(n\) individuato dalla \cref{prop:parte-intera} è
il massimo intero minore o uguale a \(x\) e prende il nome di
\emph{parte intera} di \(x\).

\begin{definition}[Funzione floor]
\label{def:floor}
Per ogni \(x\in\R\), si definisce
\[
  \lfloor x\rfloor:=n,
\]
dove \(n\in\Z\) è l'unico intero tale che
\[
  n\le x<n+1.
\]
Equivalentemente,
\[
  \lfloor x\rfloor\le x<\lfloor x\rfloor+1.
\]
\end{definition}

La funzione floor associa a ogni numero reale il più grande
intero che non lo supera. Utilizzando la notazione degli
intervalli, formalizzata nella \cref{sec:lez04-intervalli},
per ogni \(n\in\Z\) si ha
\[
  \lfloor x\rfloor=n
  \quad\Longleftrightarrow\quad
  n\le x<n+1
  \quad\Longleftrightarrow\quad
  x\in[n,n+1).
\]
In particolare, \(\lfloor x\rfloor=0\) per ogni
\(x\in[0,1)\).

\begin{figure}[htbp]
  \centering
  \begin{minipage}[c]{0.56\textwidth}
    \centering
    \begin{tikzpicture}[unipd/diagram,scale=0.9]
  \fill[green!12] (0,-0.3) rectangle (1,0.3);
  \draw[unipd-rule,very thin] (-2,-2) grid (2,2);
  \draw[->,unipd-muted] (-2.4,0) -- (2.6,0) node[right] {\(x\)};
  \draw[->,unipd-muted] (0,-2.4) -- (0,2.6) node[above] {\(y\)};

  \foreach \k in {-2,-1,0,1,2} {
    \draw[unipd-muted] (\k,-2.45) -- (\k,-2.55)
      node[below] {\(\k\)};
    \draw[unipd-muted] (-2.45,\k) -- (-2.55,\k)
      node[left] {\(\k\)};
  }

  \foreach \n in {-2,-1,0,1} {
    \draw[very thick,unipd-primary] (\n,\n) -- ({\n+1},\n);
    \fill[unipd-primary] (\n,\n) circle (2pt);
    \draw[unipd-primary,fill=white,thick] ({\n+1},\n) circle (2pt);
  }
  \draw[very thick,unipd-primary] (2,2) -- (2.4,2);
  \fill[unipd-primary] (2,2) circle (2pt);

  \draw[very thick,green!40!black] (0,0) -- (1,0);
  \fill[green!40!black] (0,0) circle (2pt);
  \draw[green!40!black,fill=white,thick] (1,0) circle (2pt);
  \node[fill=green!12,draw=green!40!black,rounded corners,inner sep=3pt]
    at (0.5,-0.55) {\(\lfloor x\rfloor=0\), \(x\in[0,1)\)};
\end{tikzpicture}

  \end{minipage}
  \hfill
  \begin{minipage}[c]{0.38\textwidth}
    \small
    \textbf{Esempi in Python}
    \begin{unipdcode}[language=Python,numbers=none,xleftmargin=0pt,framexleftmargin=0pt]
from math import floor

print(floor(0.5))   # Output: 0
print(floor(-0.5))  # Output: -1
print(floor(1.0))   # Output: 1
\end{unipdcode}

    \vspace{0.5em}
  \end{minipage}
  \caption{Grafico della funzione floor ed esempi d'uso in Python.}
  \label{fig:floor}
\end{figure}

Applicando la \cref{prop:parte-intera} a \(-x\),
per ogni \(x\in\R\) esiste un unico \(n\in\Z\) tale che
\[
  n\le-x<n+1.
\]
Cambiando segno e invertendo i versi delle disuguaglianze,
si ottiene
\[
  -n-1<x\le-n.
\]
Ponendo \(m=-n\), segue \(m-1<x\le m\),
con \(m\in\Z\) unico.

\begin{definition}[Funzione ceiling]
\label{def:ceiling}
Per ogni \(x\in\R\), si definisce
\[
  \lceil x\rceil:=m,
\]
dove \(m\in\Z\) è l'unico intero tale che
\[
  m-1<x\le m.
\]
Equivalentemente,
\[
  \lceil x\rceil-1<x\le\lceil x\rceil.
\]
\end{definition}

La funzione ceiling associa a ogni numero reale il più piccolo
intero maggiore o uguale a esso. Per ogni \(m\in\Z\) si ha
\[
  \lceil x\rceil=m
  \quad\Longleftrightarrow\quad
  m-1<x\le m
  \quad\Longleftrightarrow\quad
  x\in(m-1,m].
\]
In particolare, \(\lceil x\rceil=1\) per ogni
\(x\in(0,1]\).

\begin{figure}[htbp]
  \centering
  \begin{minipage}[c]{0.56\textwidth}
    \centering
    \begin{tikzpicture}[unipd/diagram,scale=0.9]
  \fill[green!12] (0,0.7) rectangle (1,1.3);
  \draw[unipd-rule,very thin] (-2,-2) grid (2,2);
  \draw[->,unipd-muted] (-2.4,0) -- (2.6,0) node[right] {\(x\)};
  \draw[->,unipd-muted] (0,-2.4) -- (0,2.6) node[above] {\(y\)};

  \foreach \k in {-2,-1,0,1,2} {
    \draw[unipd-muted] (\k,-2.45) -- (\k,-2.55)
      node[below] {\(\k\)};
    \draw[unipd-muted] (-2.45,\k) -- (-2.55,\k)
      node[left] {\(\k\)};
  }

  \draw[very thick,unipd-primary] (-2.4,-2) -- (-2,-2);
  \fill[unipd-primary] (-2,-2) circle (2pt);
  \foreach \m in {-1,0,1,2} {
    \draw[very thick,unipd-primary] ({\m-1},\m) -- (\m,\m);
    \draw[unipd-primary,fill=white,thick] ({\m-1},\m) circle (2pt);
    \fill[unipd-primary] (\m,\m) circle (2pt);
  }

  \draw[very thick,green!40!black] (0,1) -- (1,1);
  \draw[green!40!black,fill=white,thick] (0,1) circle (2pt);
  \fill[green!40!black] (1,1) circle (2pt);
  \node[fill=green!12,draw=green!40!black,rounded corners,inner sep=3pt]
    at (0.5,-0.55) {\(\lceil x\rceil=1\), \(x\in(0,1]\)};
\end{tikzpicture}

  \end{minipage}
  \hfill
  \begin{minipage}[c]{0.38\textwidth}
    \small
    \textbf{Esempi in Python}
    \begin{unipdcode}[language=Python,numbers=none,xleftmargin=0pt,framexleftmargin=0pt]
from math import ceil

print(ceil(0.5))   # Output: 1
print(ceil(-0.5))  # Output: 0
print(ceil(1.0))   # Output: 1
\end{unipdcode}

    \vspace{0.5em}
  \end{minipage}
  \caption{Grafico della funzione ceiling ed esempi d'uso in Python.}
  \label{fig:ceiling}
\end{figure}

\section[Cenno sulla radice quadrata di 2 nei reali]
{Cenno su \(\sqrt{2}\in\R\setminus\Q\)}
\label{sec:lez03-completezza-radice-due}

Nel \cref{thm:irrazionalita-radice-due} è stato dimostrato che
l'equazione \(x^2=2\) non ammette soluzioni razionali.
Si illustra ora come l'assioma di separazione permetta di
motivare l'esistenza di una soluzione reale non negativa.

Si considerino gli insiemi
\[
  A:=\{x\in\R:x\ge0,\ x^2<2\},
  \qquad
  B:=\{x\in\R:x\ge0,\ x^2>2\}.
\]
Entrambi sono non vuoti, poiché \(0\in A\) e \(2\in B\).
Inoltre, per ogni \(a\in A\) e \(b\in B\),
\[
  0\le a,\ b,
  \qquad a^2<2<b^2
  \quad\Longrightarrow\quad a<b.
\]
Pertanto \(A\) e \(B\) sono separati. Per
l'\hyperref[ass:assioma-separazione]{assioma di separazione}
esiste \(\lambda\in\R\) tale che
\[
  a\le\lambda\le b
  \qquad\forall a\in A,\ \forall b\in B.
\]
Poiché \(0\in A\), si ha anche \(\lambda\ge0\).

Per concludere occorrerebbe dimostrare che
\[
  \lambda^2=2.
\]
Questa verifica non viene sviluppata in questa sede.
Una volta stabilita l'identità, \(\lambda\) è la radice
quadrata non negativa di \(2\). Insieme al
\cref{thm:irrazionalita-radice-due}, si ottiene quindi
\[
  \sqrt{2}\in\R\setminus\Q.
\]

\chapter[Lezione 4 --- 2 ottobre 2026]{Lezione 4: densità, intervalli ed estremi}
\label{chap:lezione-04}

\section{Densità dei numeri razionali}
\label{sec:lez04-densita-razionali}

La proprietà archimedea e la proprietà della parte intera, introdotte
rispettivamente nel \cref{thm:proprieta-archimedea,prop:parte-intera},
permettono di approssimare ogni numero reale mediante numeri razionali.

\begin{theorem}[Densità di \(\Q\) in \(\R\)]
\label{thm:densita-razionali}
Per ogni \(x\in\R\) e per ogni \(\varepsilon>0\) esiste \(z\in\Q\) tale che
\[
  z\le x<z+\varepsilon.
\]
\end{theorem}

\begin{proof}
Siano \(x\in\R\) e \(\varepsilon>0\). Per il
\cref{thm:proprieta-archimedea} esiste \(q\in\N\) tale che
\[
  q>\frac{1}{\varepsilon}>0.
\]

Applicando la \cref{prop:parte-intera} al numero reale \(qx\),
esiste \(p\in\Z\) tale che
\[
  p\le qx<p+1.
\]
Dividendo per \(q>0\), il verso delle disuguaglianze
rimane invariato. Tenendo conto che \(1/q<\varepsilon\),
si ottiene
\[
  \frac{p}{q}\le x
  <\frac{p+1}{q}
  =\frac{p}{q}+\frac{1}{q}
  <\frac{p}{q}+\varepsilon.
\]
Ponendo \(z=p/q\), si ha \(z\in\Q\) e quindi
\[
  z\le x<z+\varepsilon.
\]
\end{proof}

\begin{figure}[htbp]
  \centering
  \begin{tikzpicture}[unipd/diagram]
  \fill[green!12] (0,-0.18) rectangle (4,0.18);
  \draw[->,unipd-muted] (-0.8,0) -- (5,0) node[right] {\(\R\)};
  \draw[very thick,unipd-primary] (0,0) -- (4,0);
  \fill[unipd-primary] (0,0) circle (2pt);
  \draw[unipd-primary,fill=white,thick] (4,0) circle (2pt);
  \fill[unipd-primary] (2.5,0) circle (2pt);
  \node[below=6pt] at (0,0) {\(z\)};
  \node[below=6pt] at (2.5,0) {\(x\)};
  \node[below=6pt] at (4,0) {\(z+\varepsilon\)};
  \draw[unipd-rule,thin] (0,0.2) -- (0,0.7) (4,0.2) -- (4,0.7);
  \draw[<->,unipd-muted] (0,0.6) -- (4,0.6)
    node[midway,above] {\(\varepsilon\)};
  \node at (2,-0.95) {\(z\le x<z+\varepsilon\)};
\end{tikzpicture}

  \caption{Interpretazione geometrica del \cref{thm:densita-razionali}:
  il numero razionale \(z\) dista da \(x\) meno di \(\varepsilon\).}
  \label{fig:densita-razionali}
\end{figure}

La densità dei razionali, insieme all'esistenza di numeri irrazionali
considerata nella \cref{sec:lez03-completezza-radice-due}, conduce
al seguente risultato.

\begin{theorem}[Razionali e irrazionali tra due reali]
\label{thm:razionali-irrazionali-tra-reali}
Siano \(a,b\in\R\) con \(a<b\). Allora:
\begin{enumerate}
  \item esiste \(r\in\Q\) tale che \(a<r<b\);
  \item esiste \(\alpha\in\R\setminus\Q\) tale che \(a<\alpha<b\).
\end{enumerate}
\end{theorem}

La dimostrazione del \cref{thm:razionali-irrazionali-tra-reali}
non viene sviluppata in questa sede.

\section{Intervalli della retta reale}
\label{sec:lez04-intervalli}

Siano \(a,b\in\R\) con \(a<b\). Si introducono le notazioni
\[
\begin{aligned}
  (a,b)  &=\{x\in\R:a<x<b\},\\
  [a,b]  &=\{x\in\R:a\le x\le b\},\\
  [a,b)  &=\{x\in\R:a\le x<b\},\\
  (a,b]  &=\{x\in\R:a<x\le b\}.
\end{aligned}
\]

\begin{definition}[Intervallo della retta reale]
\label{def:intervallo-retta-reale}
Sia \(A\subseteq\R\) non vuoto. Si dice che \(A\) è un
\emph{intervallo} se, per ogni coppia \(x_1,x_2\in A\) con
\(x_1<x_2\), vale
\[
  \{x\in\R:x_1<x<x_2\}\subseteq A.
\]
\end{definition}

\begin{example}[Un insieme che non è un intervallo]
\label{ex:insieme-non-intervallo}
Si consideri
\[
  A=(0,1)\cup\{2\}.
\]
L'insieme \((0,1)\) è un intervallo, mentre \(A\) non lo è.
Infatti, scegliendo
\[
  x_1=\frac12\in A,
  \qquad x_2=2\in A,
\]
si ha
\[
  \frac32\in\left(\frac12,2\right),
  \qquad \frac32\notin A.
\]
Ciò contraddice la \cref{def:intervallo-retta-reale}.
\end{example}

\begin{figure}[htbp]
  \centering
  \begin{tikzpicture}[unipd/diagram]
  \draw[->,unipd-muted] (-0.7,0) -- (5,0) node[right] {\(\R\)};
  \draw[very thick,unipd-primary] (0,0) -- (2,0);
  \foreach \p in {0,2}
    \draw[unipd-primary,fill=white,thick] (\p,0) circle (2pt);
  \foreach \p in {1,4}
    \fill[unipd-primary] (\p,0) circle (2pt);
  \draw[unipd-warning,fill=white,thick] (3,0) circle (2pt);
  \node[below=6pt] at (0,0) {\(0\)};
  \node[below=6pt] at (1,0) {\(\frac12\)};
  \node[below=6pt] at (2,0) {\(1\)};
  \node[below=6pt] at (3,0) {\(\frac32\)};
  \node[below=6pt] at (4,0) {\(2\)};
  \node[above=8pt] at (1,0) {\(\frac12\in A\)};
  \node[above=8pt,text=unipd-warning] at (3,0) {\(\frac32\notin A\)};
  \node[above=8pt] at (4.3,0) {\(2\in A\)};
  \node at (2,-1) {\(A=(0,1)\cup\{2\}\)};
\end{tikzpicture}

  \caption{L'insieme \(A=(0,1)\cup\{2\}\) non è un intervallo:
  tra \(\frac12\) e \(2\) esistono punti, come \(\frac32\),
  che non appartengono ad \(A\).}
  \label{fig:insieme-non-intervallo}
\end{figure}

\section{Maggioranti, minoranti e insiemi limitati}
\label{sec:lez04-maggioranti-minoranti}

Si consideri un insieme non vuoto \(A\subseteq\R\).

\begin{definition}[Maggiorante]
\label{def:maggiorante}
Un numero reale \(x\in\R\) si dice \emph{maggiorante} di \(A\) se
\[
  a\le x\qquad\forall a\in A.
\]
\end{definition}

\begin{definition}[Minorante]
\label{def:minorante}
Un numero reale \(x\in\R\) si dice \emph{minorante} di \(A\) se
\[
  x\le a\qquad\forall a\in A.
\]
\end{definition}

Un maggiorante o un minorante non deve necessariamente appartenere
all'insieme considerato. Per esempio, per
\[
  A=(0,1)\cup\{2\},
\]
i numeri \(2\) e \(3\) sono maggioranti, mentre \(0\) e
\(-\frac12\) sono minoranti.

\begin{definition}[Limitatezza superiore e inferiore]
\label{def:limitatezza-superiore-inferiore}
Sia \(A\subseteq\R\) non vuoto.
\begin{enumerate}
  \item \(A\) si dice \emph{superiormente limitato} se ammette
  almeno un maggiorante;
  \item \(A\) si dice \emph{inferiormente limitato} se ammette
  almeno un minorante.
\end{enumerate}
\end{definition}

\begin{definition}[Insieme limitato]
\label{def:insieme-limitato}
Sia \(A\subseteq\R\) non vuoto. L'insieme \(A\) si dice
\emph{limitato} se è sia superiormente sia inferiormente limitato.
\end{definition}

\begin{proposition}[Limitatezza dei sottoinsiemi]
\label{prop:limitatezza-sottoinsiemi}
Siano \(A,B\subseteq\R\), con \(B\neq\varnothing\) e \(B\subseteq A\).
\begin{enumerate}
  \item Se \(A\) è superiormente limitato, allora \(B\) è
  superiormente limitato.
  \item Se \(A\) è inferiormente limitato, allora \(B\) è
  inferiormente limitato.
\end{enumerate}
\end{proposition}

\begin{proof}
\emph{Primo punto.}
Se \(A\) è superiormente limitato, per la \cref{def:maggiorante}
esiste \(u\in\R\) tale che
\[
  a\le u\qquad\forall a\in A.
\]
Poiché \(B\subseteq A\), si ha \(b\le u\) per ogni \(b\in B\):
quindi \(B\) è superiormente limitato.

\emph{Secondo punto.}
Se \(A\) è inferiormente limitato, per la \cref{def:minorante}
esiste \(\ell\in\R\) tale che
\[
  \ell\le a\qquad\forall a\in A.
\]
Poiché \(B\subseteq A\), si ha \(\ell\le b\) per ogni \(b\in B\):
quindi \(B\) è inferiormente limitato.
\end{proof}

\subsection{Esempi fondamentali}
\label{subsec:lez04-esempi-limitatezza}

Per ogni \(n\in\N\) si ha \(0\le n\). Pertanto \(0\) è un
minorante di \(\N\), che risulta inferiormente limitato.

L'insieme \(\N\) non è invece superiormente limitato. Infatti,
l'esistenza di un numero reale \(x\) tale che
\[
  n\le x\qquad\forall n\in\N
\]
contraddirebbe il \cref{thm:proprieta-archimedea}.

Anche \(\Z\) non è superiormente limitato. In caso contrario,
poiché \(\N\subseteq\Z\), dalla
\cref{prop:limitatezza-sottoinsiemi} seguirebbe la limitatezza
superiore di \(\N\), già esclusa.

L'insieme \(\Z\) non è neppure inferiormente limitato. Supponendo
per assurdo che esista \(b\in\R\) tale che
\[
  b\le z\qquad\forall z\in\Z,
\]
cambiando segno si ottiene
\[
  -z\le-b\qquad\forall z\in\Z.
\]
Poiché \(-z\) descrive tutti gli interi al variare di \(z\in\Z\),
\(-b\) sarebbe un maggiorante di \(\Z\), in contraddizione
con quanto dimostrato.

Infine, dalle inclusioni
\[
  \Z\subseteq\Q\subseteq\R
\]
e dalla \cref{prop:limitatezza-sottoinsiemi} segue che né \(\Q\)
né \(\R\) sono superiormente o inferiormente limitati.

\section{Massimo e minimo}
\label{sec:lez04-massimo-minimo}

A differenza di un maggiorante o di un minorante, il massimo
e il minimo devono appartenere all'insieme considerato.

\begin{definition}[Massimo e minimo]
\label{def:massimo-minimo}
Sia \(A\subseteq\R\) non vuoto.
\begin{enumerate}
  \item \(A\) ammette \emph{massimo} se esiste \(x_0\in A\) tale che
  \[
    a\le x_0\qquad\forall a\in A.
  \]
  In tal caso si scrive \(x_0=\max A\).

  \item \(A\) ammette \emph{minimo} se esiste \(x_0\in A\) tale che
  \[
    x_0\le a\qquad\forall a\in A.
  \]
  In tal caso si scrive \(x_0=\min A\).
\end{enumerate}
\end{definition}

Per esempio, se
\[
  A=(0,1],
\]
allora \(\max A=1\), mentre \(A\) non ammette minimo.
Infatti, per ogni \(a\in A\), si ha
\[
  0<\frac a2<a\le1,
\]
quindi \(a/2\in A\) e \(a/2<a\). Nessun elemento di
\(A\) può pertanto essere il minimo.

Si introducono gli insiemi dei maggioranti e dei minoranti di \(A\):
\[
  M_A:=\{x\in\R:x\text{ è un maggiorante di }A\},
\]
\[
  m_A:=\{x\in\R:x\text{ è un minorante di }A\}.
\]
Per \(A=(0,1]\), si ha
\[
  M_A=\{x\in\R:x\ge1\},
  \qquad
  m_A=\{x\in\R:x\le0\},
\]
e pertanto
\[
  \min M_A=1,
  \qquad
  \max m_A=0.
\]

\begin{theorem}[Estremi degli insiemi di maggioranti e minoranti]
\label{thm:estremi-maggioranti-minoranti}
Sia \(A\subseteq\R\) non vuoto.
\begin{enumerate}
  \item Se \(A\) è superiormente limitato, allora \(M_A\) ammette minimo.
  \item Se \(A\) è inferiormente limitato, allora \(m_A\) ammette massimo.
\end{enumerate}
\end{theorem}

La dimostrazione del \cref{thm:estremi-maggioranti-minoranti}
non viene sviluppata in questa sede.

\section{Estremo superiore ed estremo inferiore}
\label{sec:lez04-estremi}

Il \cref{thm:estremi-maggioranti-minoranti} consente di definire
gli estremi superiore e inferiore di un insieme non vuoto.

\begin{definition}[Estremo superiore]
\label{def:estremo-superiore}
Sia \(A\subseteq\R\) non vuoto e superiormente limitato.
L'\emph{estremo superiore} di \(A\), indicato con \(\sup A\),
è il minimo dell'insieme dei maggioranti:
\[
  \sup A:=\min M_A.
\]
\end{definition}

\begin{definition}[Estremo inferiore]
\label{def:estremo-inferiore}
Sia \(A\subseteq\R\) non vuoto e inferiormente limitato.
L'\emph{estremo inferiore} di \(A\), indicato con \(\inf A\),
è il massimo dell'insieme dei minoranti:
\[
  \inf A:=\max m_A.
\]
\end{definition}

Per convenzione, se \(A\) non è superiormente limitato si scrive
\[
  \sup A=+\infty,
\]
mentre, se \(A\) non è inferiormente limitato, si scrive
\[
  \inf A=-\infty.
\]
I simboli \(+\infty\) e \(-\infty\) non rappresentano numeri reali;
sono qui utilizzati come convenzioni notazionali e saranno
introdotti formalmente nella \emph{retta reale estesa}
(\cref{def:lez07-retta-estesa}).

\begin{proposition}[Relazione tra estremi, massimo e minimo]
\label{prop:relazione-estremi-massimo-minimo}
Sia \(A\subseteq\R\) non vuoto.
\begin{enumerate}
  \item Se \(\sup A\in A\), allora
  \[
    \sup A=\max A.
  \]
  \item Se \(\inf A\in A\), allora
  \[
    \inf A=\min A.
  \]
\end{enumerate}
\end{proposition}

La dimostrazione della \cref{prop:relazione-estremi-massimo-minimo}
non viene sviluppata in questa sede.

\section{Caratterizzazione di estremo superiore e inferiore}
\label{sec:lez04-caratterizzazione-estremi}

Le \cref{def:estremo-superiore,def:estremo-inferiore}
ammettono le seguenti caratterizzazioni.

\begin{theorem}[Caratterizzazione dell'estremo superiore]
\label{thm:caratterizzazione-sup}
Sia \(A\subseteq\R\) non vuoto e superiormente limitato.
Per \(x\in\R\), l'uguaglianza
\[
  x=\sup A
\]
vale se e solo se sono soddisfatte entrambe le condizioni
\[
  x\ge a\qquad\forall a\in A
\]
e
\[
  \forall\varepsilon>0\;\exists\bar a\in A:
  \quad\bar a>x-\varepsilon.
\]
\end{theorem}

Per la \cref{def:maggiorante}, le due condizioni equivalgono
a richiedere che, per ogni \(\varepsilon>0\), esista
\(\bar a\in A\) tale che
\[
  x-\varepsilon<\bar a\le x.
\]

\begin{figure}[htbp]
  \centering
  \begin{tikzpicture}[unipd/diagram]
  \fill[green!12] (2,-0.18) rectangle (5,0.18);
  \draw[->,unipd-muted] (-0.5,0) -- (6,0) node[right] {\(\R\)};
  \foreach \p in {0,0.8,1.5,2.4,3.5,4.1,4.45}
    \fill[unipd-primary] (\p,0) circle (2pt);
  \node[above=8pt] at (0.8,0) {\(A\)};
  \node[fill=green!12,inner sep=1pt] at (4.7,0) {\(\cdots\)};
  \draw[unipd-muted] (2,-0.08) -- (2,0.08);
  \draw[unipd-muted] (5,-0.08) -- (5,0.08);
  \node[below=6pt] at (2,0) {\(x-\varepsilon\)};
  \node[below=6pt] at (5,0) {\(x=\sup A\)};
  \fill[green!40!black] (3.5,0) circle (2pt);
  \node[above=8pt,text=green!40!black] at (3.5,0) {\(\bar a\in A\)};
  \draw[unipd-rule,thin] (2,0.2) -- (2,0.85) (5,0.2) -- (5,0.85);
  \draw[<->,unipd-muted] (2,0.75) -- (5,0.75)
    node[midway,above] {\(\varepsilon>0\)};
  \node at (2.75,-1) {\(\forall\varepsilon>0\;\exists\bar a\in A:\quad x-\varepsilon<\bar a\le x\)};
\end{tikzpicture}

  \caption{Caratterizzazione di \(x=\sup A\): per ogni
  \(\varepsilon>0\) esiste \(\bar a\in A\) tale che
  \(x-\varepsilon<\bar a\le x\).}
  \label{fig:caratterizzazione-sup}
\end{figure}

\begin{proof}
Si dimostrano separatamente le due implicazioni.

\emph{Prima implicazione \((\Rightarrow)\).}
Si supponga \(x=\sup A\). Per la
\cref{def:estremo-superiore}, si ha
\[
  x=\min M_A.
\]
Quindi \(x\in M_A\), ossia
\[
  x\ge a\qquad\forall a\in A.
\]
Per ogni \(\varepsilon>0\), inoltre,
\[
  x-\varepsilon<x=\min M_A,
\]
perciò \(x-\varepsilon\notin M_A\). Dalla definizione di maggiorante
segue che esiste \(\bar a\in A\) tale che
\[
  x-\varepsilon<\bar a.
\]
Entrambe le condizioni sono quindi verificate.

\emph{Seconda implicazione \((\Leftarrow)\).}
Si supponga che valgano le due condizioni. La prima implica
\(x\in M_A\). Poiché \(A\) è superiormente limitato,
per il \cref{thm:estremi-maggioranti-minoranti} esiste
\[
  y:=\min M_A=\sup A,
\]
e, per definizione di minimo, \(y\le x\).

Se fosse \(y<x\), scegliendo \(\varepsilon=x-y>0\),
la seconda condizione fornirebbe \(\bar a\in A\) tale che
\[
  \bar a>x-\varepsilon=y,
\]
in contraddizione con il fatto che \(y\) è un maggiorante di
\(A\). Pertanto \(y\ge x\) e, insieme a \(y\le x\), si ottiene
\[
  x=y=\sup A.
\]
\end{proof}

\begin{theorem}[Caratterizzazione dell'estremo inferiore]
\label{thm:caratterizzazione-inf}
Sia \(A\subseteq\R\) non vuoto e inferiormente limitato.
Per \(x\in\R\), l'uguaglianza
\[
  x=\inf A
\]
vale se e solo se sono soddisfatte entrambe le condizioni
\[
  x\le a\qquad\forall a\in A
\]
e
\[
  \forall\varepsilon>0\;\exists\bar a\in A:
  \quad\bar a<x+\varepsilon.
\]
\end{theorem}

Per la \cref{def:minorante}, le due condizioni equivalgono
a richiedere che, per ogni \(\varepsilon>0\), esista
\(\bar a\in A\) tale che
\[
  x\le\bar a<x+\varepsilon.
\]

\begin{figure}[htbp]
  \centering
  \begin{tikzpicture}[unipd/diagram]
  \draw[->,unipd-muted] (-0.8,0) -- (6.2,0) node[right] {\(\R\)};
  \draw[unipd-muted] (1,-0.08) -- (1,0.08);
  \draw[unipd-muted] (4,-0.08) -- (4,0.08);
  \node[below=6pt] at (1,0) {\(x\)};
  \node[below=18pt] at (1,0) {\(x=\inf A\)};
  \node[below=6pt] at (4,0) {\(x+\varepsilon\)};
  \fill[green!12,rounded corners=2pt] (1,0.18) rectangle (4,0.55);
  \draw[<->,unipd-muted] (1,0.82) -- (4,0.82);
  \draw[dashed,unipd-muted] (1,0.08) -- (1,0.9);
  \draw[dashed,unipd-muted] (4,0.08) -- (4,0.9);
  \foreach \p in {1.75,4.2,4.8,5.5}
    \fill[unipd-primary] (\p,0) circle (2.1pt);
  \node[above=6pt] at (4.8,0) {\(A\)};
  \node at (2.5,0.36) {\([x,\,x+\varepsilon)\)};
  \fill[green!45!black] (1.75,0) circle (2.4pt);
  \node[text=green!45!black,anchor=base] at (1.75,-0.41) {\(\bar a\)};
\end{tikzpicture}

  \caption{Caratterizzazione di \(x=\inf A\): per ogni
  \(\varepsilon>0\) esiste \(\bar a\in A\) tale che
  \(x\le\bar a<x+\varepsilon\).}
  \label{fig:caratterizzazione-inf}
\end{figure}

\begin{exercise}[Dimostrazione della caratterizzazione dell'estremo inferiore]
\label{ex:caratterizzazione-inf}
Dimostrare il \cref{thm:caratterizzazione-inf} seguendo le due
implicazioni della dimostrazione del \cref{thm:caratterizzazione-sup}.
Si utilizzi l'insieme \(m_A\) dei minoranti e la relazione
\[
  \inf A=\max m_A.
\]
\end{exercise}

\begin{proof}[Soluzione dell'esercizio]
Si dimostrano separatamente le due implicazioni.

\emph{Prima implicazione \((\Rightarrow)\).}
Si supponga \(x=\inf A\). Per la
\cref{def:estremo-inferiore}, si ha
\[
  x=\max m_A.
\]
Quindi \(x\in m_A\), ossia
\[
  x\le a\qquad\forall a\in A.
\]
Per ogni \(\varepsilon>0\), inoltre,
\[
  x+\varepsilon>x=\max m_A,
\]
perciò \(x+\varepsilon\notin m_A\). Dalla definizione di minorante
segue che esiste \(\bar a\in A\) tale che
\[
  \bar a<x+\varepsilon.
\]
Entrambe le condizioni sono quindi verificate.

\emph{Seconda implicazione \((\Leftarrow)\).}
Si supponga che valgano le due condizioni. La prima implica
\(x\in m_A\). Poiché \(A\) è inferiormente limitato,
per il \cref{thm:estremi-maggioranti-minoranti} esiste
\[
  y:=\max m_A=\inf A,
\]
e, per definizione di massimo, \(x\le y\).

Se fosse \(y>x\), scegliendo \(\varepsilon=y-x>0\),
la seconda condizione fornirebbe \(\bar a\in A\) tale che
\[
  \bar a<x+\varepsilon=y,
\]
in contraddizione con il fatto che \(y\) è un minorante di
\(A\). Pertanto \(y\le x\) e, insieme a \(x\le y\), si ottiene
\[
  x=y=\inf A.
\]
\end{proof}

\chapter[Lezione 5 --- 6 ottobre 2026]{Lezione 5: disuguaglianza di Bernoulli, binomio di Newton ed esercizi sugli estremi}
\label{chap:lezione-05}

\section{Disuguaglianza di Bernoulli}
\label{sec:lez05-bernoulli}

\begin{theorem}[Disuguaglianza di Bernoulli]
\label{thm:disuguaglianza-bernoulli}
Per ogni \(x\in\R\) con \(x>-1\) e per ogni \(n\in\N\) vale
\[
  (1+x)^n\ge 1+nx.
\]
\end{theorem}

\begin{proof}
Sia \(x>-1\) fissato. Si procede per induzione su \(n\),
applicando il \cref{thm:principio-induzione} con valore iniziale
\(p=0\).

\emph{Passo base.}
Per \(n=0\), i due membri coincidono:
\[
  (1+x)^0=1=1+0\cdot x.
\]

\emph{Passo induttivo.}
Sia \(n\in\N\) e si supponga che
\[
  (1+x)^n\ge 1+nx.
\]
Poiché \(x>-1\), si ha \(1+x>0\). Moltiplicando per
\(1+x\), il verso della disuguaglianza rimane invariato:
\begin{align*}
  (1+x)^{n+1}
    &\ge (1+nx)(1+x)\\
    &=1+x+nx+nx^2\\
    &=1+(n+1)x+nx^2\\
    &\ge 1+(n+1)x.
\end{align*}
L'ultima disuguaglianza segue da \(nx^2\ge0\), poiché
\(n\ge0\) e \(x^2\ge0\). Il passo induttivo è quindi verificato.

Per il principio di induzione, la disuguaglianza vale per ogni
\(n\in\N\). Poiché \(x>-1\) era arbitrario, segue la tesi.
\end{proof}

\section{Binomio di Newton}
\label{sec:lez05-binomio-newton}

I coefficienti binomiali introdotti nella
\cref{def:coefficiente-binomiale} permettono di esprimere la potenza
di una somma mediante una sommatoria finita.

\begin{theorem}[Binomio di Newton]
\label{thm:binomio-newton}
Per ogni \(a,b\in\R\) e per ogni \(n\in\N\) vale
\[
  (a+b)^n=\sum_{k=0}^{n}\binom{n}{k}a^k b^{n-k}.
\]
Si adotta la convenzione \(t^0=1\), anche per \(t=0\).
\end{theorem}

Per \(n\ge4\), lo sviluppo assume la forma
\[
  (a+b)^n
  =b^n+nab^{n-1}+\cdots+na^{n-1}b+a^n.
\]
In particolare, per \(n=2\) e \(n=3\) si ottengono le identità
\[
  (a+b)^2=b^2+2ab+a^2,
\]
\[
  (a+b)^3=b^3+3ab^2+3a^2b+a^3.
\]

Il teorema si dimostra mediante il \cref{thm:principio-induzione},
utilizzando l'identità di Pascal della \cref{prop:identita-pascal}.
La dimostrazione completa è riportata
nell'\cref{ex:tutorato01-binomio-newton} del primo tutorato.

\section{Esercizi sul fattoriale}
\label{sec:lez05-esercizi-fattoriale}

\begin{exercise}[Confronto tra fattoriale e potenza di \(2\)]
\label{ex:fattoriale-maggiore-due-n}
Dimostrare che, per ogni \(n\in\N\) con \(n\ge4\), vale
\[
  n!>2^n.
\]
\end{exercise}

\begin{proof}[Soluzione dell'esercizio]
Per \(n<4\) si ha
\[
\begin{aligned}
  0!&=1=2^0,\\
  1!&=1<2=2^1,\\
  2!&=2<4=2^2,\\
  3!&=6<8=2^3.
\end{aligned}
\]
La disuguaglianza richiesta non vale quindi per \(n<4\).
Si procede per induzione, applicando il
\cref{thm:principio-induzione} con valore iniziale \(p=4\).

\emph{Passo base.}
Per \(n=4\),
\[
  4!=24>16=2^4.
\]

\emph{Passo induttivo.}
Sia \(n\ge4\) e si supponga \(n!>2^n\).
Per la \cref{def:fattoriale}, moltiplicando per \(n+1>0\),
si ottiene
\[
  (n+1)!=(n+1)n!>(n+1)2^n.
\]
Poiché \(n+1\ge5>2\), si ha inoltre
\[
  (n+1)2^n>2\cdot2^n=2^{n+1}.
\]
Pertanto \((n+1)!>2^{n+1}\), e il passo induttivo è verificato.

Per il \cref{thm:principio-induzione}, la disuguaglianza
vale per ogni \(n\ge4\).
\end{proof}

\begin{exercise}[Confronto tra fattoriale e quadrato]
\label{ex:fattoriale-maggiore-n-quadro}
Dimostrare che, per ogni \(n\in\N\) con \(n\ge4\), vale
\[
  n!>n^2.
\]
\end{exercise}

\begin{proof}[Soluzione dell'esercizio]
Nel passo induttivo sarà necessario confrontare \(n^2\)
con \(n+1\). Per ogni \(n\ge2\) si ha
\[
  n^2-(n+1)=n(n-1)-1\ge2-1=1>0,
\]
quindi \(n^2>n+1\).

Si procede per induzione, applicando il
\cref{thm:principio-induzione} con valore iniziale \(p=4\).

\emph{Passo base.}
Per \(n=4\),
\[
  4!=24>16=4^2.
\]

\emph{Passo induttivo.}
Sia \(n\ge4\) e si supponga \(n!>n^2\).
Per la \cref{def:fattoriale}, moltiplicando per \(n+1>0\),
si ottiene
\[
  (n+1)!=(n+1)n!>(n+1)n^2.
\]
Poiché \(n\ge4\), dalla disuguaglianza appena dimostrata
segue \(n^2>n+1\). Pertanto
\[
  (n+1)!>(n+1)n^2>(n+1)^2.
\]
Il passo induttivo è verificato. Per il
\cref{thm:principio-induzione}, la disuguaglianza
vale per ogni \(n\ge4\).
\end{proof}

\section{Esercizi sugli estremi}
\label{sec:lez05-esercizi-estremi}

Negli esercizi si utilizzano le
\cref{def:estremo-superiore,def:estremo-inferiore} e le
convenzioni per gli estremi infiniti della \cref{sec:lez04-estremi},
con l'ordinamento della \cref{def:lez07-retta-estesa}.

\begin{exercise}[Estremi e inclusione]
\label{ex:estremi-inclusione}
Siano \(A,B\subseteq\R\) non vuoti, con \(A\subseteq B\).
\begin{enumerate}
  \item Stabilire se necessariamente \(\sup A\in B\).
  \item Dimostrare che \(\inf B\le\inf A\).
  \item Dimostrare che \(\sup A\le\sup B\).
\end{enumerate}
\end{exercise}

\begin{proof}[Soluzione dell'esercizio]
\emph{Primo punto.}
L'appartenenza \(\sup A\in B\) non è garantita, neppure quando
\(A\) è superiormente limitato. Si considerino, ad esempio,
\[
  A=(0,1),\qquad B=(0,1)\cup\{2\}.
\]
Si ha \(A\subseteq B\), ma \(\sup A=1\notin B\):
l'estremo superiore di \((0,1)\) è \(1\) per il
\cref{thm:caratterizzazione-sup}.
Se \(A\) non è superiormente limitato, \(\sup A=+\infty\)
non può appartenere a \(B\subseteq\R\).

\emph{Secondo punto.}
Se \(B\) non è inferiormente limitato, si ha \(\inf B=-\infty\),
quindi \(\inf B\le\inf A\) nell'ordinamento esteso.

Se \(B\) è inferiormente limitato, per la
\cref{prop:limitatezza-sottoinsiemi} lo è anche \(A\).
Per la \cref{def:estremo-inferiore}, \(\inf B\) è un minorante
di \(B\) e, poiché \(A\subseteq B\), si ha
\[
  \inf B\le a\qquad\forall a\in A.
\]
Dunque \(\inf B\in m_A\). Essendo \(\inf A=\max m_A\),
si conclude
\[
  \inf B\le\inf A.
\]

\emph{Terzo punto.}
Se \(B\) non è superiormente limitato, si ha \(\sup B=+\infty\),
quindi \(\sup A\le\sup B\) nell'ordinamento esteso.

Se \(B\) è superiormente limitato, per la
\cref{prop:limitatezza-sottoinsiemi} lo è anche \(A\).
Per la \cref{def:estremo-superiore}, \(\sup B\) è un maggiorante
di \(B\) e, poiché \(A\subseteq B\), si ha
\[
  a\le\sup B\qquad\forall a\in A.
\]
Dunque \(\sup B\in M_A\). Essendo \(\sup A=\min M_A\),
si conclude
\[
  \sup A\le\sup B.
\]
\end{proof}

\begin{exercise}[Estremi dell'insieme degli opposti]
\label{ex:estremi-insieme-opposto}
Sia \(A\subseteq\R\) non vuoto e si consideri l'insieme
\(-A\) definito nella \cref{sec:lez02-numeri-interi}.
Dimostrare che
\[
  \inf(-A)=-\sup A,
  \qquad
  \sup(-A)=-\inf A.
\]
Per gli estremi infiniti si adottano le convenzioni
\(-(+\infty)=-\infty\) e \(-(-\infty)=+\infty\).
\end{exercise}

\begin{proof}[Soluzione dell'esercizio]
Poiché \(A\) è non vuoto, lo è anche \(-A\).
Cambiando segno, il verso delle disuguaglianze si inverte:
\(a\le t\Longleftrightarrow-a\ge-t\).

\emph{Prima identità.}
Si supponga che \(A\) sia superiormente limitato e si ponga
\(s=\sup A\in\R\). Per la \cref{def:estremo-superiore},
si ha \(a\le s\) per ogni \(a\in A\), quindi
\[
  -s\le-a\qquad\forall a\in A.
\]
Pertanto \(-s\) è un minorante di \(-A\).

Sia \(\ell\in\R\) un qualsiasi minorante di \(-A\).
Allora \(\ell\le-a\) per ogni \(a\in A\), da cui
\[
  a\le-\ell\qquad\forall a\in A.
\]
Il numero \(-\ell\) è dunque un maggiorante di \(A\).
Poiché \(s\) è il minimo dei maggioranti, si ha
\(s\le-\ell\), ossia \(\ell\le-s\).
Pertanto \(-s\) è il massimo dei minoranti di \(-A\), e
per la \cref{def:estremo-inferiore} si conclude
\[
  \inf(-A)=-s=-\sup A.
\]

Se \(A\) non è superiormente limitato, per ogni \(t\in\R\)
esiste \(a\in A\) tale che \(a>-t\); altrimenti \(-t\)
sarebbe un maggiorante di \(A\). Cambiando segno, si ottiene
\[
  -a<t.
\]
L'insieme \(-A\) non è quindi inferiormente limitato.
Poiché \(\sup A=+\infty\), si conclude
\[
  \inf(-A)=-\infty=-\sup A.
\]

\emph{Seconda identità.}
Si supponga che \(A\) sia inferiormente limitato e si ponga
\(i=\inf A\in\R\). Per la \cref{def:estremo-inferiore},
si ha \(i\le a\) per ogni \(a\in A\), quindi
\[
  -a\le-i\qquad\forall a\in A.
\]
Pertanto \(-i\) è un maggiorante di \(-A\).

Sia \(u\in\R\) un qualsiasi maggiorante di \(-A\).
Allora \(-a\le u\) per ogni \(a\in A\), da cui
\[
  -u\le a\qquad\forall a\in A.
\]
Il numero \(-u\) è dunque un minorante di \(A\).
Poiché \(i\) è il massimo dei minoranti, si ha
\(-u\le i\), ossia \(-i\le u\).
Pertanto \(-i\) è il minimo dei maggioranti di \(-A\), e
per la \cref{def:estremo-superiore} si conclude
\[
  \sup(-A)=-i=-\inf A.
\]

Se \(A\) non è inferiormente limitato, per ogni \(t\in\R\)
esiste \(a\in A\) tale che \(a<-t\); altrimenti \(-t\)
sarebbe un minorante di \(A\). Cambiando segno, si ottiene
\[
  -a>t.
\]
L'insieme \(-A\) non è quindi superiormente limitato.
Poiché \(\inf A=-\infty\), si conclude
\[
  \sup(-A)=+\infty=-\inf A.
\]
\end{proof}

\chapter[Lezione 6 --- 13 ottobre 2026]{Lezione 6: applicazioni, immagini e controimmagini}
\label{chap:lezione-06}

\section{Un esercizio sugli estremi di insiemi separati}
\label{sec:lez06-insiemi-separati}

La separazione introdotta nella \cref{sec:lez03-numeri-reali}
si può esprimere mediante gli estremi della
\cref{sec:lez04-estremi}.

\begin{exercise}[Estremi di insiemi separati]
\label{ex:lez06-estremi-separati}
Siano \(A,B\subseteq\R\) non vuoti e si supponga che
\[
  a\le b\qquad\forall a\in A,\ \forall b\in B.
\]
Dimostrare che \(\sup A\le\inf B\).
\end{exercise}

\begin{proof}[Soluzione dell'esercizio]
Scelti \(a_0\in A\) e \(b_0\in B\), l'ipotesi implica
\[
  a_0\le b\quad\forall b\in B,
  \qquad a\le b_0\quad\forall a\in A.
\]
Per le \cref{def:maggiorante,def:minorante}, entrambi gli
estremi richiesti sono quindi reali.

Ogni \(b\in B\) è un maggiorante di \(A\).
Per la \cref{def:estremo-superiore}, si ha
\[
  \sup A\le b\qquad\forall b\in B.
\]
Ne segue che \(\sup A\) è un minorante di \(B\).
Per la \cref{def:estremo-inferiore}, si conclude
\[
  \sup A\le\inf B.
\]
\end{proof}

\section{Prodotto cartesiano e applicazioni}
\label{sec:lez06-applicazioni}

Per descrivere una corrispondenza tra due insiemi
si considerano le coppie ordinate formate dai
rispettivi elementi.

\begin{definition}[Prodotto cartesiano]
\label{def:lez06-prodotto-cartesiano}
Siano \(A\) e \(B\) due insiemi. Il loro
\emph{prodotto cartesiano} è l'insieme
\[
  A\times B:=\{(a,b):a\in A,\ b\in B\}.
\]
\end{definition}

Un sottoinsieme \(\Phi\subseteq A\times B\)
descrive una corrispondenza tra gli elementi
di \(A\) e quelli di \(B\). Non ogni insieme
di coppie ordinate definisce tuttavia
un'applicazione.

\begin{definition}[Applicazione]
\label{def:lez06-applicazione}
Un insieme \(\Phi\subseteq A\times B\) si dice
\emph{applicazione} da \(A\) a \(B\) se
\[
  \forall a\in A\;\exists!\,b\in B:
  (a,b)\in\Phi.
\]
\end{definition}

Indicando con \(\phi(a)\) l'unico elemento associato
a \(a\in A\), si scrive
\[
  \phi\colon A\longrightarrow B,
  \qquad a\longmapsto\phi(a).
\]
Per ogni \(a\in A\) e \(b\in B\), si ha
\[
  \phi(a)=b
  \quad\Longleftrightarrow\quad
  (a,b)\in\Phi.
\]
L'insieme \(\Phi\) costituisce quindi il
\emph{grafico} dell'applicazione \(\phi\):
\[
  \Phi=\{(a,\phi(a)):a\in A\}.
\]

\begin{figure}[htbp]
  \centering
  \begin{tikzpicture}[unipd/diagram]
  \foreach \shift in {0,4.4,8.8} {
    \begin{scope}[xshift=\shift cm]
      \draw[unipd/edge] (0,0) ellipse (0.55 and 1.15);
      \draw[unipd/edge] (2.35,0) ellipse (0.55 and 1.15);
      \node at (0,1.45) {\(A\)};
      \node at (2.35,1.45) {\(B\)};
      \foreach \y in {-.55,.55} {
        \fill[unipd-primary] (0,\y) circle (2pt);
        \fill[unipd-primary] (2.35,\y) circle (2pt);
      }
    \end{scope}
  }
  \draw[unipd/directed edge] (0,.55) -- (2.35,.55);
  \draw[unipd/directed edge] (4.4,.55) -- (6.75,.55);
  \draw[unipd/directed edge] (4.4,.55) -- (6.75,-.55);
  \draw[unipd/directed edge] (4.4,-.55) -- (6.75,-.55);
  \draw[unipd/directed edge] (8.8,.55) -- (11.15,.55);
  \draw[unipd/directed edge] (8.8,-.55) -- (11.15,.55);
  \node[align=center] at (1.18,-1.55) {valore mancante};
  \node[align=center] at (5.58,-1.55) {due valori};
  \node[align=center] at (9.98,-1.55) {applicazione};
\end{tikzpicture}

  \caption{Solo la corrispondenza a destra è
  un'applicazione; le altre violano rispettivamente
  l'esistenza e l'unicità del valore associato.}
  \label{fig:lez06-relazioni}
\end{figure}

Per una corrispondenza rappresentata nel piano,
il \emph{test della retta verticale} permette
di verificare l'unicità del valore associato:
nessuna retta verticale deve intersecare il
grafico in più di un punto. Nella
\cref{fig:lez06-test-verticale}, la parabola
\(y=x^2\), sul dominio rappresentato \([-1,1]\),
soddisfa tale condizione, mentre la curva
\(x=y^2\) e la circonferenza la violano.
Per definire un'applicazione occorre inoltre
che ogni retta con ascissa nel dominio
intersechi il grafico almeno una volta.

\begin{figure}[htbp]
  \centering
  \begin{tikzpicture}[unipd/diagram,scale=.9]
  \foreach \shift in {0,4.6,9.2} {
    \begin{scope}[xshift=\shift cm]
      \draw[->,unipd-muted] (-1.35,0) -- (1.45,0)
        node[right] {\(x\)};
      \draw[->,unipd-muted] (0,-1.25) -- (0,1.4)
        node[above] {\(y\)};
    \end{scope}
  }
  \draw[unipd/edge,very thick,domain=-1:1,samples=70]
    plot (\x,{\x*\x});
  \draw[unipd-accent,dashed] (.5,-1.15) -- (.5,1.25);
  \fill[unipd-accent] (.5,.25) circle (1.7pt);
  \node at (0,-1.65) {\(y=x^2\)};

  \begin{scope}[xshift=4.6cm]
    \draw[unipd/edge,very thick,domain=-1:1,samples=70]
      plot ({\x*\x},\x);
    \draw[unipd-accent,dashed] (.36,-1.15) -- (.36,1.25);
    \foreach \y in {-.6,.6} {
      \fill[unipd-accent] (.36,\y) circle (1.7pt);
    }
    \node at (0,-1.65) {\(x=y^2\)};
  \end{scope}

  \begin{scope}[xshift=9.2cm]
    \draw[unipd/edge,very thick] (0,0) circle (1);
    \draw[unipd-accent,dashed] (.6,-1.15) -- (.6,1.25);
    \foreach \y in {-.8,.8} {
      \fill[unipd-accent] (.6,\y) circle (1.7pt);
    }
    \node at (0,-1.65) {\(x^2+y^2=1\)};
  \end{scope}
\end{tikzpicture}

  \caption{Test della retta verticale: la prima
  curva presenta un'unica intersezione per ogni
  ascissa del dominio; le altre due ne presentano
  più di una per le ascisse evidenziate.}
  \label{fig:lez06-test-verticale}
\end{figure}

\section{Dominio, codominio e immagine}
\label{sec:lez06-dominio-immagine}

\begin{definition}[Dominio e codominio]
\label{def:lez06-dominio-codominio}
Sia \(f\colon A\longrightarrow B\). L'insieme
\(A\) si dice \emph{dominio} di \(f\) e si indica
con \(D(f)=A\); l'insieme \(B\) si dice
\emph{codominio} e si indica con \(C(f)=B\).
\end{definition}

L'espressione che determina i valori di una
funzione non basta a individuarla: occorre
specificarne anche il dominio e il codominio.
Si considerino, ad esempio, le applicazioni
\[
  \begin{aligned}
    f\colon\R&\longrightarrow\R,
      &x&\longmapsto x^2,\\
    g\colon\R&\longrightarrow[0,+\infty),
      &x&\longmapsto x^2,\\
    h\colon[0,+\infty)&\longrightarrow[0,+\infty),
      &x&\longmapsto x^2.
  \end{aligned}
\]
Le tre applicazioni hanno la stessa espressione
algebrica: \(f\) e \(g\) differiscono per il codominio,
\(g\) e \(h\) per il dominio.

\begin{definition}[Immagine di un'applicazione]
\label{def:lez06-immagine}
L'\emph{immagine} di \(f\colon A\to B\)
è l'insieme
\[
  \operatorname{Im}(f)
  :=\{b\in B:\exists a\in A,\ f(a)=b\}
  =f(A).
\]
\end{definition}

Per la \cref{def:lez06-immagine}, si ha
\[
  \operatorname{Im}(f)\subseteq C(f)=B,
\]
e l'inclusione può essere stretta. Per le tre
applicazioni precedenti,
\[
  \operatorname{Im}(f)
  =\operatorname{Im}(g)
  =\operatorname{Im}(h)=[0,+\infty).
\]
Infatti, il quadrato di un numero reale è
non negativo e, per ogni \(y\ge0\), si ha
\(y=(\sqrt y)^2\). La \cref{fig:lez06-quadrato}
mostra il grafico del quadrato e distingue,
nel caso di \(f\), l'immagine dal codominio.

\begin{figure}[htbp]
  \centering
  \begin{tikzpicture}[unipd/diagram,x=1.25cm,y=.78cm]
  \fill[unipd-accent!12] (-2.2,1) rectangle (2.2,4);
  \draw[unipd-muted,dashed] (-2.2,1) -- (2.2,1);
  \draw[unipd-muted,dashed] (-2.2,4) -- (2.2,4);
  \draw[->,unipd-muted] (-2.5,0) -- (2.75,0) node[right] {\(x\)};
  \draw[->,unipd-muted] (0,-1.05) -- (0,5.05) node[above] {\(y\)};
  \draw[very thick,unipd-primary,domain=-2.2:2.2,samples=90]
    plot (\x,{\x*\x});
  \draw[very thick,unipd-accent] (0,1) -- (0,4);
  \draw[very thick,unipd-accent] (-2,0) -- (-1,0);
  \draw[very thick,unipd-accent] (1,0) -- (2,0);
  \foreach \x in {-2,-1,1,2} {
    \draw[unipd-muted] (\x,.08) -- (\x,-.08)
      node[below] {\(\x\)};
  }
  \foreach \y in {-1,1,4} {
    \draw[unipd-muted] (.07,\y) -- (-.07,\y)
      node[left] {\(\y\)};
  }
  \node[anchor=west] at (2.45,3.8) {\(S=[1,4]\)};
  \node[anchor=west,text=unipd-primary] at (1.05,4.55)
    {\(y=x^2\)};
\end{tikzpicture}

  \caption{Per \(f\colon\R\to\R\), l'immagine
  \([0,+\infty)\) è distinta dal codominio \(\R\).
  All'intervallo \(S=[1,4]\) corrispondono
  le ascisse \([-2,-1]\cup[1,2]\).}
  \label{fig:lez06-quadrato}
\end{figure}

\section{Controimmagine di un sottoinsieme}
\label{sec:lez06-controimmagine}

\begin{definition}[Controimmagine]
\label{def:lez06-controimmagine}
Sia \(f\colon A\longrightarrow B\) e sia
\(S\subseteq B\). La \emph{controimmagine}
di \(S\) mediante \(f\) è l'insieme
\[
  f^{-1}(S):=\{a\in A:f(a)\in S\}.
\]
\end{definition}

\begin{figure}[htbp]
  \centering
  \begin{tikzpicture}[unipd/diagram]
  \fill[unipd-accent!13] (0,.5) ellipse (.49 and .75);
  \fill[unipd-accent!13] (5,.5) ellipse (.49 and .75);
  \draw[unipd/edge] (0,0) ellipse (.95 and 1.65);
  \draw[unipd/edge] (5,0) ellipse (.95 and 1.65);
  \draw[unipd-accent,dashed] (0,.5) ellipse (.49 and .75);
  \draw[unipd-accent,dashed] (5,.5) ellipse (.49 and .75);
  \node at (0,1.92) {\(A\)};
  \node at (5,1.92) {\(B\)};
  \node at (2.5,1.32) {\(f\)};
  \node[text=unipd-accent!80!black] at (0,-1.95)
    {\(f^{-1}(S)\)};
  \node[text=unipd-accent!80!black] at (5,-1.95) {\(S\)};

  \draw[unipd/directed edge] (.09,.91) -- (4.91,.91);
  \draw[unipd/directed edge] (.09,.19) -- (4.91,.19);
  \draw[unipd/directed edge] (.09,-.55) -- (4.91,-.98);
  \draw[unipd/directed edge] (.09,-1.2) -- (4.91,-.98);
  \foreach \y/\n in {.91/1,.19/2,-.55/3,-1.2/4} {
    \fill[unipd-primary] (0,\y) circle (2pt);
    \node[left] at (-.08,\y) {\(a_{\n}\)};
  }
  \foreach \y/\n in {.91/1,.19/2,-.98/3} {
    \fill[unipd-primary] (5,\y) circle (2pt);
    \node[right] at (5.09,\y) {\(b_{\n}\)};
  }
\end{tikzpicture}

  \caption{Per \(S=\{b_1,b_2\}\subseteq B\) si ha
  \(f^{-1}(S)=\{a_1,a_2\}\subseteq A\), anche se
  \(f\) non è iniettiva.}
  \label{fig:lez06-controimmagine-insiemi}
\end{figure}

Si indica con
\[
  \mathcal P(A):=\{S:S\subseteq A\}
\]
l'\emph{insieme delle parti} di \(A\), e
analogamente con \(\mathcal P(B)\) quello di
\(B\). Per ogni \(S\in\mathcal P(B)\), la
\cref{def:lez06-controimmagine} determina
un unico sottoinsieme di \(A\), eventualmente
vuoto. Si ottiene quindi un'applicazione
tra insiemi delle parti:
\[
  f^{-1}\colon\mathcal P(B)\longrightarrow\mathcal P(A),
  \qquad S\longmapsto f^{-1}(S).
\]
In particolare, se \(S\cap\operatorname{Im}(f)
=\varnothing\), allora \(f^{-1}(S)=\varnothing\).

\begin{example}[Controimmagini del quadrato]
\label{ex:lez06-controimmagini-quadrato}
Per l'applicazione \(f\colon\R\to\R\) definita
nella \cref{sec:lez06-dominio-immagine}, si ha
\[
  f^{-1}(\{-1\})=\varnothing,
  \qquad
  f^{-1}(\{0\})=\{0\},
\]
\[
  f^{-1}([1,4])=[-2,-1]\cup[1,2].
\]
La controimmagine dell'intervallo è
rappresentata nella \cref{fig:lez06-quadrato}.

L'equazione \(z^2=-1\) ammette soluzioni
complesse, ma \(z\mapsto z^2\) non definisce
un'applicazione da \(\C\) a \(\R\): ad esempio,
\[
  (1+\imaginaryunit)^2
  =2\imaginaryunit\notin\R.
\]
Considerando invece l'applicazione
\(q\colon\C\to\C\), \(q(z)=z^2\), si ottiene
\[
  q^{-1}(\{-1\})
  =\{\imaginaryunit,-\imaginaryunit\}.
\]
\end{example}

La notazione \(f^{-1}(S)\) indica una
controimmagine anche quando \(f\) non è
iniettiva né biiettiva. Non presuppone
l'esistenza di un'applicazione inversa
\(B\to A\).

\section{Applicazioni suriettive, iniettive e biiettive}
\label{sec:lez06-proprieta-applicazioni}

\begin{definition}[Suriettività]
\label{def:lez06-suriettiva}
Un'applicazione \(f\colon A\longrightarrow B\)
si dice \emph{suriettiva} se
\(\operatorname{Im}(f)=B\), ossia se
\[
  \forall b\in B\;\exists a\in A:
  f(a)=b.
\]
\end{definition}

La suriettività dipende dal codominio, come illustrato
nell'\cref{ex:lez06-quadrato-proprieta}. Per la
\cref{def:lez06-controimmagine}, la controimmagine di
\(\{b\}\) è non vuota per ogni \(b\in B\) se e solo se
\(f\) è suriettiva.

\begin{definition}[Iniettività]
\label{def:lez06-iniettiva}
Un'applicazione \(f\colon A\longrightarrow B\)
si dice \emph{iniettiva} se, per ogni
\(a_1,a_2\in A\), vale
\[
  a_1\ne a_2
  \quad\Longrightarrow\quad
  f(a_1)\ne f(a_2).
\]
\end{definition}

\begin{exercise}[Caratterizzazione dell'iniettività]
\label{ex:lez06-equivalenza-iniettiva}
Dimostrare che \(f\colon A\to B\) è iniettiva
se e solo se, per ogni \(a_1,a_2\in A\),
\[
  f(a_1)=f(a_2)
  \quad\Longrightarrow\quad
  a_1=a_2.
\]
\end{exercise}

\begin{proof}[Soluzione dell'esercizio]
Siano \(a_1,a_2\in A\) arbitrari.

Se \(f\) è iniettiva e
\(f(a_1)=f(a_2)\), non può valere
\(a_1\ne a_2\): altrimenti, per la
\cref{def:lez06-iniettiva}, si avrebbe
\(f(a_1)\ne f(a_2)\). Pertanto \(a_1=a_2\).

Viceversa, si supponga che
\(f(a_1)=f(a_2)\) implichi \(a_1=a_2\)
per ogni coppia di elementi di \(A\).
Se \(a_1\ne a_2\), non può valere
\(f(a_1)=f(a_2)\), perché si avrebbe
\(a_1=a_2\). Quindi \(f(a_1)\ne f(a_2)\)
e \(f\) è iniettiva.
\end{proof}

\begin{definition}[Biiettività]
\label{def:lez06-biettiva}
Un'applicazione \(f\colon A\longrightarrow B\)
si dice \emph{biiettiva} se è sia iniettiva
sia suriettiva.
\end{definition}

Per le \cref{def:lez06-suriettiva,def:lez06-iniettiva},
un'applicazione è biiettiva se e solo se
ogni elemento del codominio possiede uno
e un solo antecedente nel dominio:
\[
  \forall b\in B\;\exists!\,a\in A:
  f(a)=b.
\]
L'esistenza corrisponde alla \cref{def:lez06-suriettiva};
l'unicità segue dall'\cref{ex:lez06-equivalenza-iniettiva}.

\begin{figure}[htbp]
  \centering
  \begin{tikzpicture}[unipd/diagram]
  \foreach \shift in {0,4.4,8.8} {
    \begin{scope}[xshift=\shift cm]
      \draw[unipd/edge] (0,0) ellipse (0.55 and 1.25);
      \draw[unipd/edge] (2.35,0) ellipse (0.55 and 1.25);
      \node at (0,1.5) {\(A\)};
      \node at (2.35,1.5) {\(B\)};
    \end{scope}
  }
  \foreach \y in {-.65,.65} {
    \fill[unipd-primary] (0,\y) circle (2pt);
    \fill[unipd-primary] (8.8,\y) circle (2pt);
    \fill[unipd-primary] (11.15,\y) circle (2pt);
    \fill[unipd-primary] (6.75,\y) circle (2pt);
  }
  \foreach \y in {-.8,0,.8} {
    \fill[unipd-primary] (2.35,\y) circle (2pt);
    \fill[unipd-primary] (4.4,\y) circle (2pt);
  }
  \draw[unipd/directed edge] (0,.65) -- (2.35,.8);
  \draw[unipd/directed edge] (0,-.65) -- (2.35,-.8);
  \draw[unipd/directed edge] (4.4,.8) -- (6.75,.65);
  \draw[unipd/directed edge] (4.4,0) -- (6.75,.65);
  \draw[unipd/directed edge] (4.4,-.8) -- (6.75,-.65);
  \draw[unipd/directed edge] (8.8,.65) -- (11.15,.65);
  \draw[unipd/directed edge] (8.8,-.65) -- (11.15,-.65);
  \node at (1.18,-1.65) {iniettiva};
  \node at (5.58,-1.65) {suriettiva};
  \node at (9.98,-1.65) {biiettiva};
\end{tikzpicture}

  \caption{Da sinistra: applicazione iniettiva
  non suriettiva, suriettiva non iniettiva e
  biiettiva.}
  \label{fig:lez06-proprieta}
\end{figure}

\section{Tre applicazioni definite dal quadrato}
\label{sec:lez06-esempi-quadrato}

Si riprendono le tre applicazioni definite nella
\cref{sec:lez06-dominio-immagine}, distinguendone
iniettività e suriettività.

\begin{example}[Dominio e codominio del quadrato]
\label{ex:lez06-quadrato-proprieta}
Le applicazioni \(f\) e \(g\) non sono
iniettive, poiché \(1\ne-1\) e
\(1^2=(-1)^2\). Inoltre, \(f\) non è
suriettiva su \(\R\), mentre \(g\) è
suriettiva su \([0,+\infty)\). Nessuna
delle due è quindi biiettiva.

Per ogni \(a>0\), entrambi i domini
contengono le due radici reali di \(a\),
per cui
\[
  f^{-1}(\{a\})=g^{-1}(\{a\})
  =\{-\sqrt a,\sqrt a\}.
\]
Per \(a=0\), le controimmagini sono
\(\{0\}\); se \(a<0\), caso possibile
soltanto nel codominio di \(f\), si ha
\(f^{-1}(\{a\})=\varnothing\).

Restringendo il dominio a \([0,+\infty)\),
l'applicazione \(h\) è iniettiva.
Infatti, se \(x_1,x_2\ge0\) e
\(x_1^2=x_2^2\), allora
\[
  (x_1-x_2)(x_1+x_2)=0.
\]
Se \(x_1+x_2=0\), entrambi i numeri
sono nulli; altrimenti, \(x_1-x_2=0\).
In ogni caso, \(x_1=x_2\).

Poiché ogni \(a\ge0\) è immagine di
\(\sqrt a\in[0,+\infty)\), l'applicazione
\(h\) è anche suriettiva e quindi biiettiva.
In particolare, per ogni \(a\ge0\),
\[
  h^{-1}(\{a\})=\{\sqrt a\}.
\]
\end{example}

Per \(h\) esiste quindi anche
l'\emph{applicazione inversa}, distinta
dalla costruzione sulle parti introdotta
nella \cref{sec:lez06-controimmagine}:
\[
  h^{-1}\colon[0,+\infty)\longrightarrow[0,+\infty),
  \qquad y\longmapsto\sqrt y.
\]
L'applicazione inversa associa a ogni
valore il suo unico antecedente.
La scrittura \(h^{-1}(\{a\})\), in cui
l'argomento è un insieme, indica invece
la controimmagine di tale insieme.

\chapter[Lezione 7 --- 14 ottobre 2026]{Lezione 7: numerabilità, operazioni sulle funzioni e punti di accumulazione}
\label{chap:lezione-07}

\section{Numerabilità degli insiemi}
\label{sec:lez07-numerabilita}

La nozione di applicazione biiettiva, introdotta nella
\cref{def:lez06-biettiva}, permette di confrontare insiemi
anche quando contengono infiniti elementi. Due insiemi
hanno la stessa \emph{cardinalità} se esiste una biiezione
tra essi.

\begin{definition}[Insieme numerabile]
\label{def:lez07-numerabilita}
Un insieme \(A\) si dice \emph{numerabile} se esiste
un'applicazione biiettiva
\[
  \varphi\colon\N\longrightarrow A.
\]
Un insieme infinito che non ammette una tale biiezione
si dice \emph{non numerabile}.
\end{definition}

\begin{proposition}[Numerabilità dei naturali pari]
\label{prop:lez07-pari}
L'insieme dei numeri naturali pari
\[
  P:=\{2m:m\in\N\}
\]
è numerabile. In particolare, l'applicazione
\[
  \varphi\colon\N\longrightarrow P,
  \qquad m\longmapsto 2m,
\]
è biiettiva.
\end{proposition}

\begin{proof}
Si verificano separatamente la suriettività e
l'iniettività dell'applicazione.

\emph{Suriettività.}
Sia \(p\in P\) arbitrario. Per definizione di \(P\),
esiste \(m\in\N\) tale che
\[
  p=2m=\varphi(m).
\]
Ogni elemento di \(P\) possiede quindi almeno un
antecedente in \(\N\). Per la
\cref{def:lez06-suriettiva}, \(\varphi\) è suriettiva.

\emph{Iniettività.}
Siano \(m_1,m_2\in\N\) con \(m_1\ne m_2\).
Supponiamo per assurdo che
\(\varphi(m_1)=\varphi(m_2)\). Allora
\[
  2m_1=2m_2
  \quad\Longrightarrow\quad
  m_1=m_2,
\]
in contraddizione con l'ipotesi \(m_1\ne m_2\).
Pertanto \(\varphi(m_1)\ne\varphi(m_2)\),
quindi \(\varphi\) è iniettiva.

Essendo sia suriettiva sia iniettiva,
\(\varphi\) è biiettiva.
\end{proof}

Analogamente, l'insieme dei naturali dispari
\[
  D:=\{2m+1:m\in\N\}
\]
è numerabile mediante l'applicazione
\[
  \psi\colon\N\longrightarrow D,
  \qquad m\longmapsto 2m+1.
\]
La suriettività segue dalla definizione di \(D\);
l'iniettività dall'implicazione
\[
  2m_1+1=2m_2+1
  \quad\Longrightarrow\quad
  m_1=m_2.
\]

Poiché \(P\subsetneq\N\) (infatti \(1\notin P\)),
la \cref{prop:lez07-pari} mostra che può esistere
una biiezione tra un insieme infinito e un suo
sottoinsieme proprio.

\begin{proposition}[Numerabilità degli insiemi numerici]
\label{prop:lez07-insiemi-numerici}
Gli insiemi \(\N\), \(\Z\) e \(\Q\) sono numerabili,
mentre \(\R\) non è numerabile.
\end{proposition}

Per \(\N\) è sufficiente considerare l'applicazione
identità. La numerabilità di \(\Z\) e \(\Q\) e la
non numerabilità di \(\R\) sono qui enunciate
senza dimostrazione. Quest'ultimo risultato
è dovuto a Cantor.

La numerabilità non va confusa con la densità:
in particolare, \(\Q\) è numerabile pur essendo
denso in \(\R\), come stabilito nel
\cref{thm:densita-razionali}.

\begin{example}[Albergo di Hilbert]
\label{ex:lez07-hilbert}
Si consideri un albergo con infinite camere,
numerate mediante i numeri naturali, tutte occupate.

Sebbene non vi siano camere libere, è possibile
accogliere un nuovo ospite trasferendo
simultaneamente ciascun occupante dalla camera
\(n\) alla camera \(n+1\), per ogni \(n\in\N\).

L'assegnazione
\[
  n\longmapsto n+1
\]
è iniettiva, quindi nessuna camera riceve più
di un ospite. Dopo lo spostamento, la camera
\(0\) rimane libera e può accogliere il nuovo
arrivato.

L'assegnazione costituisce una biiezione tra
\(\N\) e \(\N\setminus\{0\}\).

\begin{figure}[H]
  \centering
  \begin{tikzpicture}[unipd/diagram,x=1cm,y=1cm]
  \node[left] at (-.6,0) {Prima};
  \node[left] at (-.6,-1.65) {Dopo};
  \foreach \n in {0,...,5} {
    \draw[unipd-primary] ({\n-.38},-.35) rectangle ({\n+.38},.35);
    \node at (\n,.13) {\(\n\)};
    \node[text=unipd-muted] at (\n,-.16) {\(\times\)};
    \draw[unipd-primary] ({\n-.38},-2) rectangle ({\n+.38},-1.3);
    \node at (\n,-1.52) {\(\n\)};
    \draw[->,unipd-muted] (\n,-.4) -- ({\n+1},-1.25);
  }
  \foreach \n in {1,...,5} {
    \node[text=unipd-muted] at (\n,-1.81) {\(\times\)};
  }
  \node at (6,0) {\(\dots\)};
  \node at (6,-1.65) {\(\dots\)};
  \node[right] at (6.2,-.8) {\(n\mapsto n+1\)};
  \node[below] at (0,-2.05) {libera};
\end{tikzpicture}

  \caption{Schema dell'albergo di Hilbert: dopo lo
  spostamento \(n\mapsto n+1\), la camera \(0\)
  risulta libera. Le crocette indicano le camere
  occupate.}
  \label{fig:lez07-hilbert}
\end{figure}
\end{example}


\section{Operazioni sulle funzioni}
\label{sec:lez07-operazioni}

Si considerino due funzioni a valori reali
\[
  \psi\colon A\longrightarrow\R,\qquad
  \varphi\colon B\longrightarrow\R,
  \qquad A,B\subseteq\R.
\]
Per la \cref{def:lez06-dominio-codominio}, i domini
di \(\psi\) e \(\varphi\) sono rispettivamente \(A\)
e \(B\). La somma e il prodotto dei loro valori
possono quindi essere calcolati per ogni
\(x\in A\cap B\).

Le funzioni così ottenute hanno codominio
\(\R\), mentre la loro immagine, secondo la
\cref{def:lez06-immagine}, è costituita dai
valori effettivamente assunti.

\begin{definition}[Somma di funzioni]
\label{def:lez07-somma}
La \emph{somma} di \(\psi\) e \(\varphi\)
è la funzione
\[
  \psi+\varphi\colon A\cap B\longrightarrow\R,
  \qquad
  (\psi+\varphi)(x):=\psi(x)+\varphi(x).
\]
In particolare,
\(D(\psi+\varphi)=D(\psi)\cap D(\varphi)\).
\end{definition}

La somma è definita punto per punto mediante
l'addizione in \(\R\). La proprietà commutativa,
già introdotta per \(\N\) nella
\cref{sec:lez01-numeri-naturali}, vale anche in
\(\R\). Poiché le funzioni \(\psi+\varphi\) e
\(\varphi+\psi\) hanno lo stesso dominio e
codominio, per ogni \(x\in A\cap B\) si ha
\[
  (\psi+\varphi)(x)
  =\psi(x)+\varphi(x)
  =\varphi(x)+\psi(x)
  =(\varphi+\psi)(x).
\]
Pertanto \(\psi+\varphi=\varphi+\psi\).

\begin{example}[Somma di una radice e di una funzione lineare]
\label{ex:lez07-somma-radice}
Siano \(\varphi\colon[0,+\infty)\to\R\),
\(\varphi(x)=\sqrt{x}\), e
\(\psi\colon\R\to\R\), \(\psi(x)=2x\). Si ha
\[
  D(\psi+\varphi)
  =\R\cap[0,+\infty)
  =[0,+\infty),
\]
e, per ogni \(x\ge0\),
\[
  (\psi+\varphi)(x)=2x+\sqrt{x}.
\]
La somma non è definita per \(x<0\), poiché
tali valori non appartengono al dominio
di \(\varphi\).
\end{example}

\begin{definition}[Prodotto di funzioni]
\label{def:lez07-prodotto}
Il \emph{prodotto} di \(\psi\) e \(\varphi\)
è la funzione
\[
  \psi\varphi\colon A\cap B\longrightarrow\R,
  \qquad
  (\psi\varphi)(x):=\psi(x)\varphi(x).
\]
In particolare,
\(D(\psi\varphi)=D(\psi)\cap D(\varphi)\).
\end{definition}

Il prodotto è definito punto per punto mediante
la moltiplicazione in \(\R\), che è commutativa.
Come per la somma, le funzioni \(\psi\varphi\) e
\(\varphi\psi\) hanno lo stesso dominio e
codominio e, per ogni \(x\in A\cap B\), si ha
\[
  (\psi\varphi)(x)
  =\psi(x)\varphi(x)
  =\varphi(x)\psi(x)
  =(\varphi\psi)(x).
\]
Pertanto \(\psi\varphi=\varphi\psi\).

\begin{definition}[Rapporto di funzioni]
\label{def:lez07-quoziente}
Il \emph{rapporto} di \(\psi\) e \(\varphi\)
è la funzione
\[
  \frac{\psi}{\varphi}\colon
  (A\cap B)\setminus Z_\varphi
  \longrightarrow\R,
  \qquad
  \left(\frac{\psi}{\varphi}\right)(x)
  :=\frac{\psi(x)}{\varphi(x)},
\]
dove
\[
  Z_\varphi:=\{x\in D(\varphi):\varphi(x)=0\}
\]
è l'insieme degli zeri di \(\varphi\).
In particolare,
\[
  D(\psi/\varphi)
  =(D(\psi)\cap D(\varphi))\setminus Z_\varphi.
\]
\end{definition}

Il denominatore deve essere diverso da zero:
per questo gli zeri di \(\varphi\) sono esclusi
dal dominio comune.

Scambiando numeratore e denominatore si ottiene
\[
  D(\varphi/\psi)
  =(D(\psi)\cap D(\varphi))\setminus Z_\psi,
\]
dove \(Z_\psi\) indica l'insieme degli zeri
di \(\psi\).

A differenza della somma e del prodotto, il rapporto
non è commutativo: i due domini possono essere
diversi e, anche nei punti in cui entrambi i
rapporti sono definiti, i loro valori non
coincidono in generale.

\section{Grafici e trasformazioni delle funzioni}
\label{sec:lez07-trasformazioni}

Nella \cref{sec:lez06-applicazioni}, il grafico
di un'applicazione è stato introdotto come insieme
delle coppie ordinate che la rappresentano.
Si considera ora il caso delle funzioni reali
di variabile reale.

\begin{definition}[Grafico di una funzione]
\label{def:lez07-grafico}
Sia \(\varphi\colon A\to\R\), con \(A\subseteq\R\).
Il \emph{grafico} di \(\varphi\) è l'insieme
\[
  G(\varphi):=\{(x,\varphi(x)):x\in A\}
  \subseteq\R^2.
\]
\end{definition}

Ogni punto del grafico ha ascissa \(x\in A\) e
ordinata \(\varphi(x)\). Le trasformazioni considerate
nel seguito agiscono sulle ascisse o sulle ordinate,
permettendo di ottenere nuovi grafici a partire da
\(G(\varphi)\).

\subsection{Traslazione rigida}
\label{sec:lez07-traslazioni-verticali}

Sia \(\varphi\colon A\longrightarrow\R\), con
\(A\subseteq\R\), e sia \(c\in\R\). Si consideri
la funzione
\[
  \varphi+c\colon A\longrightarrow\R,
  \qquad
  (\varphi+c)(x):=\varphi(x)+c.
\]
Poiché la costante si somma ai valori di
\(\varphi\), il dominio rimane invariato:
\[
  D(\varphi+c)=D(\varphi)=A.
\]

Per la \cref{def:lez07-grafico}, il grafico
della nuova funzione è
\[
  G(\varphi+c)
  =\{(x,\varphi(x)+c):x\in A\}.
\]
A ogni punto \((x,\varphi(x))\in G(\varphi)\)
corrisponde quindi il punto
\[
  (x,\varphi(x)+c)=(x,\varphi(x))+(0,c).
\]
Viceversa, ogni punto di \(G(\varphi+c)\)
si ottiene in questo modo da un punto di
\(G(\varphi)\). Il nuovo grafico è pertanto
ottenuto mediante una \emph{traslazione rigida}
verticale secondo il vettore \((0,c)\).

Se \(c>0\), si ha
\[
  \varphi(x)+c>\varphi(x),
\]
e il grafico viene traslato verso l'alto.
Se \(c<0\), la disuguaglianza si inverte
e la traslazione avviene verso il basso.

In entrambi i casi, la distanza verticale
tra i punti corrispondenti è costante:
\[
  |(\varphi+c)(x)-\varphi(x)|=|c|.
\]
Per \(c=0\), la funzione e il suo grafico
rimangono invariati.

\subsection{Riscalamento verticale e ribaltamento}
\label{sec:lez07-riscalamento}

Sia \(\varphi\colon A\longrightarrow\R\), con
\(A\subseteq\R\), e sia \(b\in\R\). Si consideri
la funzione
\[
  b\varphi\colon A\longrightarrow\R,
  \qquad
  (b\varphi)(x):=b\,\varphi(x).
\]
Poiché il coefficiente \(b\) moltiplica i valori
di \(\varphi\), il dominio rimane invariato:
\[
  D(b\varphi)=D(\varphi)=A.
\]

Per la \cref{def:lez07-grafico}, il grafico
della nuova funzione è
\[
  G(b\varphi)=\{(x,b\varphi(x)):x\in A\}.
\]
La corrispondenza tra i due grafici è quindi
\[
  (x,\varphi(x))\longmapsto(x,b\varphi(x)).
\]
L'ascissa rimane invariata, mentre l'ordinata
viene moltiplicata per \(b\).

La distanza di un punto dall'asse delle ascisse
passa da \(|\varphi(x)|\) a
\[
  |b\varphi(x)|=|b|\,|\varphi(x)|.
\]
Ne segue che, se \(|b|>1\), il grafico subisce
una dilatazione verticale; se \(0<|b|<1\), si ha
invece una compressione verticale. I punti di
ordinata nulla rimangono fissi.

Se \(b<0\), oltre al riscalamento si ha un
cambiamento di segno delle ordinate. In
particolare, per \(b=-1\) si ottiene la funzione
\[
  -\varphi\colon A\longrightarrow\R,
  \qquad
  (-\varphi)(x)=-\varphi(x),
\]
e il grafico corrispondente si ottiene mediante
un \emph{ribaltamento} rispetto all'asse delle
ascisse.

Per \(b=1\), il grafico rimane invariato. Per
\(b=0\), si ottiene invece
\[
  G(0\varphi)=\{(x,0):x\in A\},
\]
ossia la parte dell'asse delle ascisse le cui
ascisse appartengono al dominio di \(\varphi\).

Non si può dedurre da \(b>1\) che
\(b\varphi(x)>\varphi(x)\) per ogni \(x\), poiché
il confronto dipende dal segno di \(\varphi(x)\).
La distanza verticale tra i punti corrispondenti
dei grafici di \(\varphi\) e \(b\varphi\) è
\[
  |b\varphi(x)-\varphi(x)|
  =|b-1|\,|\varphi(x)|.
\]

\begin{figure}[htbp]
  \centering
  \begin{tikzpicture}[unipd/diagram,x=1.2cm,y=1.2cm]
  \begin{scope}
    \draw[->,unipd-muted] (-1.4,0) -- (1.65,0) node[right] {\(x\)};
    \draw[->,unipd-muted] (0,-.3) -- (0,2.5) node[above] {\(y\)};
    \draw[very thick,unipd-primary,domain=-1:1,samples=50]
      plot (\x,{\x*\x});
    \draw[very thick,unipd-accent,domain=-1:1,samples=50]
      plot (\x,{\x*\x+1});
    \draw[->,unipd-muted,dashed] (.65,.4225) -- (.65,1.4225);
    \node[right=5pt,text=unipd-primary] at (1.02,1.08) {\(\varphi\)};
    \node[right=5pt,text=unipd-accent] at (1.02,2.08) {\(\varphi+1\)};
    \draw[unipd-muted] (.05,1) -- (-.05,1);
    \node[below left=3pt] at (0,1) {\(1\)};
  \end{scope}
  \begin{scope}[xshift=6.5cm]
    \draw[->,unipd-muted] (-1.4,0) -- (1.65,0) node[right] {\(x\)};
    \draw[->,unipd-muted] (0,-1.4) -- (0,1.5) node[above] {\(y\)};
    \draw[very thick,unipd-primary,domain=-1:1,samples=50]
      plot (\x,{\x*\x});
    \draw[very thick,unipd-accent,domain=-1:1,samples=50]
      plot (\x,{-\x*\x});
    \node[right=5pt,text=unipd-primary] at (1.02,1.08) {\(\varphi\)};
    \node[right=5pt,text=unipd-accent] at (1.02,-1.08) {\(-\varphi\)};
    \draw[unipd-muted] (.05,1) -- (-.05,1) node[left] {\(1\)};
    \draw[unipd-muted] (.05,-1) -- (-.05,-1) node[left] {\(-1\)};
  \end{scope}
\end{tikzpicture}

  \caption{A sinistra, traslazione rigida del
  grafico di \(\varphi(x)=x^2\) verso l'alto di
  un'unità; a destra, ribaltamento rispetto
  all'asse delle ascisse, da \(\varphi(x)=x^2\)
  a \(-\varphi(x)=-x^2\). Le curve sono
  rappresentate per \(x\in[-1,1]\).}
  \label{fig:lez07-verticali}
\end{figure}

\subsection{Traslazione orizzontale}
\label{sec:lez07-traslazioni-orizzontali}

Sia \(\varphi\colon A\longrightarrow\R\), con
\(A\subseteq\R\), e sia \(c\in\R\). Si consideri
la funzione
\[
  \psi\colon A_c\longrightarrow\R,
  \qquad
  \psi(x):=\varphi(x+c),
\]
dove
\[
  A_c:=\{x\in\R:x+c\in A\}
      =\{t-c:t\in A\}.
\]
Poiché l'argomento di \(\varphi\) deve appartenere
ad \(A\), il dominio della nuova funzione è
\[
  D(\psi)=A_c.
\]

Per la \cref{def:lez07-grafico}, il grafico
della nuova funzione è
\[
  G(\psi)=\{(t-c,\varphi(t)):t\in A\}.
\]
A ogni punto \((t,\varphi(t))\in G(\varphi)\)
corrisponde quindi il punto
\[
  (t-c,\varphi(t))
  =(t,\varphi(t))+(-c,0).
\]
Viceversa, ogni punto di \(G(\psi)\) si ottiene
in questo modo da un punto di \(G(\varphi)\).
Il nuovo grafico è pertanto ottenuto mediante
una \emph{traslazione rigida} orizzontale secondo
il vettore \((-c,0)\).

Se \(c>0\), il grafico viene traslato verso
sinistra; se \(c<0\), la traslazione avviene
verso destra. In entrambi i casi, la distanza
orizzontale tra i punti corrispondenti è
\[
  |(t-c)-t|=|c|.
\]
Per \(c=0\), la funzione e il suo grafico
rimangono invariati.

\begin{example}[Traslazione del dominio di un intervallo]
\label{ex:lez07-dominio-traslato}
Sia \(\varphi\colon[-1,1]\longrightarrow\R\).
Si considerino le funzioni
\[
  \psi_1(x):=\varphi(x+2),
  \qquad
  \psi_2(x):=\varphi(x-2).
\]

Per \(\psi_1\), la condizione sull'argomento è
\[
  -1\le x+2\le1
  \quad\Longleftrightarrow\quad
  -3\le x\le-1,
\]
quindi
\[
  D(\psi_1)=[-3,-1].
\]

Analogamente, per \(\psi_2\) si ha
\[
  -1\le x-2\le1
  \quad\Longleftrightarrow\quad
  1\le x\le3,
\]
da cui
\[
  D(\psi_2)=[1,3].
\]
\end{example}

\begin{figure}[htbp]
  \centering
  \begin{tikzpicture}[unipd/diagram,x=1.55cm,y=1.3cm]
  \draw[->,unipd-muted] (-3.45,0) -- (3.55,0) node[right] {\(x\)};
  \draw[->,unipd-muted] (0,-.4) -- (0,1.8) node[above] {\(y\)};
  \draw[very thick,unipd-primary,domain=-1:1,samples=50]
    plot (\x,{\x*\x});
  \draw[very thick,unipd-accent,domain=-3:-1,samples=50]
    plot (\x,{(\x+2)*(\x+2)});
  \draw[very thick,unipd-accent,dashed,domain=1:3,samples=50]
    plot (\x,{(\x-2)*(\x-2)});
  \foreach \x in {-3,-1,1,3} {
    \fill[unipd-primary] (\x,1) circle (1.8pt);
  }
  \foreach \x in {-3,-2,-1,1,2,3} {
    \draw[unipd-muted] (\x,.05) -- (\x,-.05) node[below] {\(\x\)};
  }
  \node[above,text=unipd-accent] at (-2,1.25) {\(\varphi(x+2)\)};
  \node[above,text=unipd-primary] at (0.35,1.25) {\(\varphi(x)\)};
  \node[above,text=unipd-accent] at (2,1.25) {\(\varphi(x-2)\)};
\end{tikzpicture}

  \caption{Traslazioni orizzontali del grafico
  di \(\varphi(x)=x^2\), con dominio \([-1,1]\).
  La funzione \(\psi_1(x)=\varphi(x+2)\) ha dominio
  \([-3,-1]\), mentre \(\psi_2(x)=\varphi(x-2)\)
  ha dominio \([1,3]\). I punti pieni indicano
  gli estremi inclusi.}
  \label{fig:lez07-orizzontali}
\end{figure}

\section{Intorni e punti di accumulazione}
\label{sec:lez07-intorni-accumulazione}

Per studiare il comportamento di una funzione
in prossimità di un punto, si introducono le
nozioni di intorno e di punto di accumulazione,
che saranno utilizzate nella definizione di limite.

\subsection{Intorni di punti reali e dell'infinito}
\label{sec:lez07-intorni}

\begin{definition}[Intorno circolare]
\label{def:lez07-intorno}
Siano \(x_0\in\R\) e \(\delta>0\).
L'\emph{intorno circolare} di centro \(x_0\)
e raggio \(\delta\) è l'insieme
\[
  D(x_0,\delta)
  :=\{x\in\R:|x-x_0|<\delta\}
  =(x_0-\delta,x_0+\delta).
\]
\end{definition}

La condizione \(|x-x_0|<\delta\) individua
i punti della retta reale la cui distanza
da \(x_0\) è inferiore a \(\delta\).
Il centro appartiene all'intorno, mentre
gli estremi \(x_0-\delta\) e \(x_0+\delta\)
sono esclusi.

\begin{figure}[htbp]
  \centering
  \begin{tikzpicture}[unipd/diagram,x=1.6cm]
  \draw[->,unipd-muted] (-3,0) -- (3.2,0) node[right] {\(x\)};
  \draw[very thick,unipd-primary] (-2,0) -- (2,0);
  \foreach \x in {-2,2} {
    \draw[unipd-primary,fill=white,thick] (\x,0) circle (2.4pt);
  }
  \fill[unipd-primary] (0,0) circle (1.8pt);
  \node[below=5pt] at (-2,0) {\(x_0-\delta\)};
  \node[below=5pt] at (0,0) {\(x_0\)};
  \node[below=5pt] at (2,0) {\(x_0+\delta\)};
  \draw[<->,unipd-muted] (-2,.5) -- (0,.5) node[midway,above] {\(\delta\)};
  \draw[<->,unipd-muted] (0,.5) -- (2,.5) node[midway,above] {\(\delta\)};
\end{tikzpicture}

  \caption{Intorno circolare di centro \(x_0\)
  e raggio \(\delta>0\).}
  \label{fig:lez07-intorno}
\end{figure}

\begin{definition}[Intorni di infinito]
\label{def:lez07-intorni-infinito}
Si dice \emph{intorno di \(+\infty\)}
ogni intervallo della forma
\[
  (a,+\infty)=\{x\in\R:x>a\},
  \qquad a\in\R.
\]
Analogamente, si dice \emph{intorno di \(-\infty\)}
ogni intervallo della forma
\[
  (-\infty,a)=\{x\in\R:x<a\},
  \qquad a\in\R.
\]
\end{definition}

A differenza degli intorni circolari,
gli intorni di infinito sono intervalli
illimitati della retta reale. I simboli
\(+\infty\) e \(-\infty\) non rappresentano
numeri reali e non appartengono agli
intervalli così definiti.

\begin{definition}[Retta reale estesa]
\label{def:lez07-retta-estesa}
La \emph{retta reale estesa} è l'insieme
\[
  \overline{\R}:=\R\cup\{-\infty,+\infty\},
\]
ordinato ponendo
\[
  -\infty<x<+\infty
  \qquad\forall x\in\R.
\]
\end{definition}

\begin{example}[Estremi di intervalli illimitati]
\label{ex:lez07-estremi-infiniti}
Per ogni \(a\in\R\), secondo le convenzioni
introdotte nella \cref{sec:lez04-estremi}, si ha
\[
  \sup(a,+\infty)=+\infty,
  \qquad
  \inf(a,+\infty)=a,
\]
\[
  \sup(-\infty,a)=a,
  \qquad
  \inf(-\infty,a)=-\infty.
\]

L'intervallo \((a,+\infty)\) non è
superiormente limitato, mentre
\((-\infty,a)\) non è inferiormente limitato.
Gli estremi infiniti appartengono alla
retta reale estesa, mentre gli estremi
finiti coincidono con \(a\), che non
appartiene a nessuno dei due intervalli.

Pertanto, nessuno dei due intervalli
ammette massimo o minimo.
\end{example}


\subsection{Punti di accumulazione}
\label{sec:lez07-punti-accumulazione}

\begin{definition}[Punto di accumulazione]
\label{def:lez07-accumulazione}
Sia \(A\subseteq\R\) non vuoto. Un punto
\(x_0\in\R\) si dice \emph{punto di accumulazione}
di \(A\) se
\[
  \forall\delta>0,\qquad
  \bigl(D(x_0,\delta)\cap A\bigr)\setminus\{x_0\}
  \neq\varnothing.
\]
L'insieme dei punti di accumulazione di \(A\)
si indica con \(A'\).
\end{definition}

Per la \cref{def:lez07-intorno}, la condizione
richiede che ogni intorno circolare di centro
\(x_0\) contenga almeno un elemento di \(A\)
distinto dal centro. Equivalentemente,
\[
  \forall\delta>0\ \exists a\in A:
  \qquad 0<|a-x_0|<\delta.
\]

La definizione non richiede che \(x_0\in A\).
Un punto può quindi essere di accumulazione
anche senza appartenere all'insieme.

Viceversa, un punto \(x_0\in A\) si dice
\emph{isolato} se esiste un intorno circolare
che non contiene altri elementi di \(A\).
In generale, negando la condizione della
\cref{def:lez07-accumulazione}, si ottiene
\[
  x_0\notin A'
  \quad\Longleftrightarrow\quad
  \exists\delta>0:
  \bigl(D(x_0,\delta)\cap A\bigr)
  \setminus\{x_0\}=\varnothing.
\]

Per dimostrare che un punto è di accumulazione
occorre dunque verificare la condizione per
ogni \(\delta>0\); per escluderlo è sufficiente
individuare un solo raggio positivo per cui
la condizione non sia soddisfatta.

\begin{example}[Un intervallo aperto e un punto isolato]
\label{ex:lez07-accumulazione-intervallo}
Sia
\[
  A=(1,2)\cup\{3\}.
\]
L'insieme dei suoi punti di accumulazione è
\[
  A'=[1,2].
\]
\end{example}

\begin{proof}
Si dimostra anzitutto che ogni punto di
\([1,2]\) è di accumulazione.

Sia \(x_0\in(1,2)\). Per ogni \(\delta>0\),
si ponga
\[
  \eta=
  \min\left\{
    \frac{\delta}{2},
    \frac{2-x_0}{2}
  \right\},
  \qquad
  a=x_0+\eta.
\]
Poiché \(\eta>0\), si ha \(a\in(1,2)\)
e
\[
  0<|a-x_0|=\eta<\delta.
\]
Per la \cref{def:lez07-accumulazione},
\(x_0\in A'\).

Anche gli estremi \(1\) e \(2\) sono punti
di accumulazione. Per ogni \(\delta>0\),
posto
\[
  \eta=\min\left\{\frac{\delta}{2},
                  \frac12\right\},
\]
i punti \(1+\eta\) e \(2-\eta\) appartengono
ad \(A\) e distano rispettivamente da
\(1\) e da \(2\) una quantità positiva
minore di \(\delta\).

Pertanto
\[
  [1,2]\subseteq A'.
\]

Si dimostra ora che nessun punto esterno
a \([1,2]\) è di accumulazione.

Se \(x_0<1\), scegliendo
\[
  \delta=\frac{1-x_0}{2}>0
\]
si ottiene
\[
  D(x_0,\delta)\cap A=\varnothing.
\]
Quindi \(x_0\notin A'\).

Se \(x_0>2\) e \(x_0\ne3\), si scelga
\[
  \delta=
  \frac12\min\{x_0-2,|x_0-3|\}>0.
\]
L'intorno \(D(x_0,\delta)\) non interseca
l'intervallo \((1,2)\) e non contiene
il punto \(3\). Pertanto
\[
  D(x_0,\delta)\cap A=\varnothing,
\]
da cui \(x_0\notin A'\).

Infine, per \(x_0=3\), scegliendo
\(\delta=1/2\) si ha
\[
  D(3,1/2)\cap A=\{3\},
\]
quindi
\[
  \bigl(D(3,1/2)\cap A\bigr)
  \setminus\{3\}=\varnothing.
\]
Il punto \(3\) è dunque isolato e non
è di accumulazione.

Si è così dimostrato che
\(A'\subseteq[1,2]\). Dalla doppia
inclusione segue
\[
  A'=[1,2].
\]
\end{proof}

\begin{figure}[htbp]
  \centering
  \begin{tikzpicture}[unipd/diagram,x=2.4cm,y=1cm]
  \foreach \y in {0,-1.3} {
    \draw[->,unipd-muted] (.55,\y) -- (3.5,\y) node[right] {\(x\)};
    \foreach \x in {1,2,3} {
      \draw[unipd-muted] (\x,{\y+.06}) -- (\x,{\y-.06})
        node[below=4pt] {\(\x\)};
    }
  }
  \node[left] at (.45,0) {\(A\)};
  \node[left] at (.45,-1.3) {\(A'\)};
  \draw[very thick,unipd-primary] (1,0) -- (2,0);
  \foreach \x in {1,2} {
    \draw[unipd-primary,fill=white,thick] (\x,0) circle (2.4pt);
  }
  \fill[unipd-primary] (3,0) circle (2.4pt);
  \draw[very thick,unipd-accent] (1,-1.3) -- (2,-1.3);
  \foreach \x in {1,2} {
    \fill[unipd-accent] (\x,-1.3) circle (2.4pt);
  }
\end{tikzpicture}

  \caption{L'insieme \(A=(1,2)\cup\{3\}\)
  e il suo insieme dei punti di accumulazione
  \(A'=[1,2]\). I cerchi vuoti indicano
  estremi esclusi, mentre i punti pieni
  indicano elementi inclusi.}
  \label{fig:lez07-accumulazione}
\end{figure}

\begin{exercise}[Verifica dei punti di accumulazione]
\label{ex:lez07-verifica-accumulazione}
Per \(A=(1,2)\cup\{3\}\), verificare direttamente
che \(2\) e ogni \(x_0\in(1,2)\) sono punti
di accumulazione, mentre ogni \(x_0>2\),
con \(x_0\ne3\), non lo è.
\end{exercise}

\begin{proof}[Soluzione dell'esercizio]
Nella dimostrazione dell'\cref{ex:lez07-accumulazione-intervallo},
per ogni \(x_0\in(1,2]\) e ogni \(\delta>0\)
è stato costruito un elemento \(a\in A\) tale che
\(0<|a-x_0|<\delta\). Per \(x_0>2\), con
\(x_0\ne3\), è stato invece individuato un
intorno disgiunto da \(A\). Le conclusioni
seguono quindi dalla \cref{def:lez07-accumulazione}.
\end{proof}

\begin{exercise}[Variante con estremi inclusi]
\label{ex:lez07-accumulazione-chiuso}
Ripetere la verifica dell'\cref{ex:lez07-verifica-accumulazione}
per l'insieme
\[
  B=[1,2]\cup\{3\}
\]
e determinarne tutti i punti di accumulazione.
\end{exercise}

\begin{proof}[Soluzione dell'esercizio]
Si ha ancora
\[
  B'=[1,2].
\]

Infatti, poiché \((1,2)\subseteq B\),
le costruzioni utilizzate nella dimostrazione
dell'\cref{ex:lez07-accumulazione-intervallo}
rimangono valide per ogni \(x_0\in[1,2]\).
Anche per gli estremi \(1\) e \(2\), gli elementi
scelti sono distinti dal centro dell'intorno.

Se \(x_0<1\) oppure \(x_0>2\), con \(x_0\ne3\),
gli stessi raggi consentono di individuare
un intorno disgiunto da \(B\).

Infine,
\[
  D(3,1/2)\cap B=\{3\},
\]
quindi \(3\) è ancora un punto isolato.

L'inclusione degli estremi nell'intervallo
modifica la loro appartenenza all'insieme,
ma non la loro proprietà di essere punti
di accumulazione.
\end{proof}

\chapter[Lezione 8 --- 15 ottobre 2026]{Lezione 8: punti isolati, accumulazione all'infinito e limiti di funzioni}
\label{chap:lezione-08}

\section{Punti isolati e accumulazione all'infinito}
\label{sec:lez08-accumulazione-infinito}

La \cref{def:lez07-accumulazione} si completa introducendo
i punti isolati e l'accumulazione nei punti infiniti della
\cref{def:lez07-retta-estesa}.

\subsection{Punti isolati}
\label{sec:lez08-isolati}

\begin{definition}[Punto isolato]
\label{def:lez08-isolato}
Sia \(A\subseteq\R\) non vuoto. Un punto \(x_0\in A\)
si dice \emph{isolato} in \(A\) se esiste \(\delta>0\)
tale che
\[
  D(x_0,\delta)\cap A=\{x_0\}.
\]
\end{definition}

Negando la condizione della \cref{def:lez07-accumulazione},
per ogni \(x_0\in\R\) si ottiene
\[
  x_0\notin A'
  \quad\Longleftrightarrow\quad
  \exists\delta>0:
  \bigl(D(x_0,\delta)\cap A\bigr)
  \setminus\{x_0\}=\varnothing.
\]
Se \(x_0\in A\), questa condizione equivale alla
\cref{def:lez08-isolato}, poiché il centro appartiene
a ogni suo intorno circolare. Si ottiene quindi, per ogni \(x_0\in A\), l'equivalenza
\[
  x_0\text{ è isolato in }A
  \quad\Longleftrightarrow\quad
  x_0\notin A'.
\]
L'ipotesi \(x_0\in A\) è essenziale: si confrontino i punti
\(3\) e \(1\) nell'\cref{ex:lez07-accumulazione-intervallo}:
il primo è isolato in \(A=(1,2)\cup\{3\}\),
mentre il secondo è di accumulazione pur non appartenendo ad \(A\).

\subsection{Punti di accumulazione infiniti}
\label{sec:lez08-accumulazione-estesa}

Gli intorni di \(+\infty\) e \(-\infty\), introdotti nella
\cref{def:lez07-intorni-infinito}, permettono di estendere
la \cref{def:lez07-accumulazione} ai punti infiniti. In questo
caso non occorre escludere il centro dell'intorno, poiché
\(+\infty\) e \(-\infty\) non sono numeri reali.

\begin{definition}[Punti di accumulazione infiniti]
\label{def:lez08-accumulazione-infinito}
Sia \(A\subseteq\R\) non vuoto. Il simbolo \(+\infty\) si dice
\emph{punto di accumulazione} di \(A\) se
\[
  \forall b\in\R,\qquad A\cap(b,+\infty)\ne\varnothing.
\]
Analogamente, \(-\infty\) si dice \emph{punto di accumulazione}
di \(A\) se
\[
  \forall b\in\R,\qquad A\cap(-\infty,b)\ne\varnothing.
\]
\end{definition}

Il collegamento con la \cref{def:limitatezza-superiore-inferiore}
è dato dalla seguente proposizione.

\begin{proposition}[Punti di accumulazione infiniti e limitatezza]
\label{prop:lez08-accumulazione-limitatezza}
Sia \(A\subseteq\R\) non vuoto. Allora:
\begin{enumerate}
  \item \(+\infty\) è un punto di accumulazione di \(A\) se e solo
  se \(A\) non è superiormente limitato;
  \item \(-\infty\) è un punto di accumulazione di \(A\) se e solo
  se \(A\) non è inferiormente limitato.
\end{enumerate}
\end{proposition}

\begin{proof}
\emph{Primo punto.}
Se \(A\) è superiormente limitato, per la \cref{def:maggiorante}
esiste \(\lambda\in\R\) tale che
\[
  a\le\lambda \qquad \forall a\in A.
\]
Scegliendo \(b=\lambda+1\), si ottiene
\[
  A\cap(b,+\infty)=\varnothing.
\]
Per la \cref{def:lez08-accumulazione-infinito},
\(+\infty\) non è quindi di accumulazione.

Viceversa, se \(+\infty\) non è di accumulazione,
per la \cref{def:lez08-accumulazione-infinito}
esiste \(b\in\R\) tale che
\[
  A\cap(b,+\infty)=\varnothing.
\]
Ne segue che \(a\le b\) per ogni \(a\in A\).
Il numero \(b\) è dunque un maggiorante di \(A\),
che risulta superiormente limitato.

\emph{Secondo punto.}
Se \(A\) è inferiormente limitato, per la \cref{def:minorante}
esiste \(\mu\in\R\) tale che
\[
  \mu\le a \qquad \forall a\in A.
\]
Scegliendo \(b=\mu-1\), si ottiene
\[
  A\cap(-\infty,b)=\varnothing.
\]
Per la \cref{def:lez08-accumulazione-infinito},
\(-\infty\) non è quindi di accumulazione.

Viceversa, se \(-\infty\) non è di accumulazione,
per la \cref{def:lez08-accumulazione-infinito}
esiste \(b\in\R\) tale che
\[
  A\cap(-\infty,b)=\varnothing.
\]
Ne segue che \(b\le a\) per ogni \(a\in A\).
Il numero \(b\) è dunque un minorante di \(A\),
che risulta inferiormente limitato.
\end{proof}

\begin{figure}[htbp]
  \centering
  \begin{tikzpicture}[unipd/diagram]
  \draw[->,unipd-muted] (-.8,0) -- (7.5,0) node[right] {\(x\)};
  \draw[very thick,unipd-primary,<-] (-.65,0) -- (2.7,0);
  \draw[unipd-primary,fill=white,thick] (2.7,0) circle (2.4pt);
  \draw[very thick,unipd-accent,->] (4.4,0) -- (7.25,0);
  \draw[unipd-accent,fill=white,thick] (4.4,0) circle (2.4pt);
  \draw[unipd-muted] (2.7,-.1) -- (2.7,.1);
  \draw[unipd-muted] (4.4,-.1) -- (4.4,.1);
  \node[below=6pt] at (2.7,0) {\(\lambda\)};
  \node[below=6pt] at (4.4,0) {\(b\)};
  \node[above=7pt,text=unipd-primary] at (.9,0) {\(A\)};
  \node[above=7pt,text=unipd-accent] at (5.8,0) {\((b,+\infty)\)};
  \node[below=21pt] at (3.55,0) {\(A\cap(b,+\infty)=\varnothing\)};
\end{tikzpicture}

  \caption{Un insieme \(A=(-\infty,\lambda)\) superiormente limitato
  non interseca l'intorno \((b,+\infty)\) di \(+\infty\), con
  \(b>\lambda\). Pertanto \(+\infty\) non è un suo punto di
  accumulazione.}
  \label{fig:lez08-accumulazione-infinito}
\end{figure}

La \cref{prop:lez08-accumulazione-limitatezza} non richiede
l'esistenza di massimi o minimi: ad esempio,
\((-\infty,\lambda)\) non ammette massimo
(\cref{fig:lez08-accumulazione-infinito}).

\section{Definizione generale di limite}
\label{sec:lez08-limiti-definizione}

Gli intorni della \cref{sec:lez07-intorni} e le nozioni di
accumulazione introdotte nelle
\cref{def:lez07-accumulazione,def:lez08-accumulazione-infinito}
permettono ora di definire il limite di una funzione.

\begin{definition}[Limite di una funzione]
\label{def:lez08-limite}
Sia \(f\colon A\longrightarrow\R\), con \(A\subseteq\R\)
non vuoto, e sia \(x_0\in\overline{\R}\) un punto di
accumulazione per \(A\). Si dice che \(f\) \emph{ammette limite}
\(\ell\in\overline{\R}\) per \(x\) che tende a \(x_0\),
e si scrive
\[
  \lim_{x\to x_0}f(x)=\ell,
\]
se, per ogni intorno \(V_\ell\) di \(\ell\), esiste un intorno
\(U_{x_0}\) di \(x_0\) tale che
\[
  f(x)\in V_\ell
  \qquad\forall x\in(U_{x_0}\cap A)\setminus\{x_0\}.
\]
\end{definition}

Non è necessario che \(x_0\in A\); anche quando
\(x_0\in A\), il valore \(f(x_0)\) non influisce sul limite.
Per \(x_0=\pm\infty\), l'esclusione del centro è superflua
poiché \(A\subseteq\R\).

L'ordine dei quantificatori è essenziale: per \emph{ogni}
intorno del valore limite deve esistere un intorno di
\(x_0\) nel quale \emph{tutti} i valori di \(f\) nei
punti ammessi del dominio, distinti da \(x_0\),
appartengano all'intorno assegnato.

\subsection{Limite finito in un punto reale}
\label{sec:lez08-limite-finito}

Nel caso \(x_0,\ell\in\R\), per la \cref{def:lez07-intorno},
si possono scegliere gli intorni circolari
\[
  V_\ell=(\ell-\varepsilon,\ell+\varepsilon),
  \qquad
  U_{x_0}=D(x_0,\delta)=(x_0-\delta,x_0+\delta),
\]
con \(\varepsilon,\delta>0\). La \cref{def:lez08-limite}
assume pertanto la forma
\[
  \forall\varepsilon>0\;\exists\delta_\varepsilon>0:
  \quad
  0<|x-x_0|<\delta_\varepsilon
  \ \Longrightarrow\ |f(x)-\ell|<\varepsilon
  \quad\forall x\in A.
\]
Il raggio \(\delta_\varepsilon\) può dipendere da
\(\varepsilon\), mentre il centro \(x_0\) resta escluso.

La \cref{fig:lez08-limite-finito} rappresenta questa condizione
per una funzione razionale, il cui limite sarà verificato
nell'\cref{ex:lez08-razionale-finito}.

\begin{figure}[htbp]
  \centering
  \begin{tikzpicture}[unipd/diagram,x=1.55cm,y=3.1cm]
  \fill[unipd-soft] (1.5,.5) rectangle (2.5,1);
  \draw[->,unipd-muted] (.35,0) -- (3.72,0) node[right] {\(x\)};
  \draw[->,unipd-muted] (.5,-.04) -- (.5,1.3) node[above] {\(y\)};
  \foreach \xx in {1.5,2,2.5} {
    \draw[dashed,unipd-muted] (\xx,0) -- (\xx,1.09);
  }
  \foreach \yy in {.5,.75,1} {
    \draw[densely dotted,unipd-muted] (.5,\yy) -- (3.4,\yy);
  }
  \draw[very thick,unipd-primary,domain=.72:3.37,samples=90]
    plot (\x,{(2*\x-1)/(\x+2)});
  \fill[unipd-primary] (2,.75) circle (1.9pt);
  \draw[unipd-muted] (1.5,-.025) -- (1.5,.025);
  \draw[unipd-muted] (2,-.025) -- (2,.025);
  \draw[unipd-muted] (2.5,-.025) -- (2.5,.025);
  \node[below=5pt] at (1.5,0) {\(x_0-\delta\)};
  \node[below=5pt] at (2,0) {\(x_0\)};
  \node[below=5pt] at (2.5,0) {\(x_0+\delta\)};
  \node[left=5pt] at (.5,.5) {\(\ell-\varepsilon\)};
  \node[left=5pt] at (.5,.75) {\(\ell\)};
  \node[left=5pt] at (.5,1) {\(\ell+\varepsilon\)};
  \node[above right=2pt,text=unipd-primary] at (2.72,.93) {\(f\)};
\end{tikzpicture}

  \caption{Rappresentazione della condizione di limite finito:
  per \(0<|x-x_0|<\delta\), i valori \(f(x)\) appartengono
  alla fascia \((\ell-\varepsilon,\ell+\varepsilon)\).
  Il grafico illustra \(f(x)=(2x-1)/(x+2)\) in prossimità
  di \(x_0=2\), con \(\ell=3/4\), \(\varepsilon=1/4\) e
  \(\delta=1/2\).}
  \label{fig:lez08-limite-finito}
\end{figure}

\subsection{Altri casi particolari}
\label{sec:lez08-limite-casi}

Poiché \(x_0,\ell\in\overline{\R}\), la
\cref{def:lez08-limite} comprende nove casi, secondo che
ciascuno dei due elementi sia reale, \(+\infty\) o
\(-\infty\). Il caso \(x_0,\ell\in\R\) è stato
esplicitato nella \cref{sec:lez08-limite-finito};
gli altri si ricavano dagli intorni della
\cref{def:lez07-intorni-infinito}.

\begin{exercise}[Forme esplicite della definizione di limite]
\label{ex:lez08-casi-limite}
Scrivere la \cref{def:lez08-limite} per gli altri otto casi
possibili, distinguendo punti di accumulazione finiti e
infiniti e valori limite finiti e infiniti.
\end{exercise}

\begin{proof}[Soluzione dell'esercizio]
Si assume in ogni caso che \(x_0\) sia di accumulazione
per \(A\) e che le implicazioni valgano per ogni \(x\in A\).
Quando \(x_0\in\R\), il punto \(x=x_0\) deve essere escluso.

\emph{Punto \(x_0\in\R\), limite infinito.}
La condizione \(\lim_{x\to x_0}f(x)=+\infty\) equivale a
\[
  \forall M\in\R\;\exists\delta>0:\quad
  0<|x-x_0|<\delta\ \Longrightarrow\ f(x)>M.
\]
Analogamente, per \(\lim_{x\to x_0}f(x)=-\infty\), si richiede
\[
  \forall M\in\R\;\exists\delta>0:\quad
  0<|x-x_0|<\delta\ \Longrightarrow\ f(x)<M.
\]

\emph{Punto \(x_0=+\infty\).}
Per \(\ell\in\R\), si richiede
\[
  \forall\varepsilon>0\;\exists b\in\R:\quad
  x>b\ \Longrightarrow\ |f(x)-\ell|<\varepsilon.
\]
Per i due limiti infiniti si ottengono rispettivamente
\[
  \forall M\in\R\;\exists b\in\R:\quad
  x>b\ \Longrightarrow\ f(x)>M
\]
e
\[
  \forall M\in\R\;\exists b\in\R:\quad
  x>b\ \Longrightarrow\ f(x)<M.
\]

\emph{Punto \(x_0=-\infty\).}
Per \(\ell\in\R\), si richiede
\[
  \forall\varepsilon>0\;\exists b\in\R:\quad
  x<b\ \Longrightarrow\ |f(x)-\ell|<\varepsilon.
\]
Per i due limiti infiniti si ottengono rispettivamente
\[
  \forall M\in\R\;\exists b\in\R:\quad
  x<b\ \Longrightarrow\ f(x)>M
\]
e
\[
  \forall M\in\R\;\exists b\in\R:\quad
  x<b\ \Longrightarrow\ f(x)<M.
\]
In tutti i casi si applica la \cref{def:lez08-limite}
con gli intorni del punto \(x_0\) e del valore limite \(\ell\).
\end{proof}

\section{Esercizi sulla definizione di limite}
\label{sec:lez08-esercizi-limiti}

Nelle verifiche della \cref{def:lez08-limite} occorre
anzitutto controllare il dominio e la condizione
di accumulazione.

\begin{exercise}[Un punto non di accumulazione del dominio]
\label{ex:lez08-radice-meno-infinito}
Stabilire se, nel senso della \cref{def:lez08-limite},
è possibile considerare il limite
\[
  \lim_{x\to-\infty}\frac{2x-1}{\sqrt{x+2}}.
\]
\end{exercise}

\begin{proof}[Soluzione dell'esercizio]
Perché il denominatore sia reale e non nullo, occorre
\[
  x+2>0\quad\Longleftrightarrow\quad x>-2.
\]
Il dominio naturale della funzione è pertanto
\[
  A=(-2,+\infty).
\]
Essendo \(A\) inferiormente limitato, per la
\cref{prop:lez08-accumulazione-limitatezza} il punto
\(-\infty\) non è di accumulazione. Manca quindi
un'ipotesi della \cref{def:lez08-limite}.
\end{proof}

\begin{exercise}[Verifica di un limite finito]
\label{ex:lez08-razionale-finito}
Per la funzione
\[
  f(x)=\frac{2x-1}{x+2},
\]
verificare mediante la \cref{def:lez08-limite} che
\[
  \lim_{x\to2}f(x)=\frac34.
\]
\end{exercise}

\begin{proof}[Soluzione dell'esercizio]
\emph{Dominio e punto di accumulazione.}
La funzione è definita esattamente quando \(x+2\ne0\),
quindi
\[
  A=\R\setminus\{-2\}
   =(-\infty,-2)\cup(-2,+\infty).
\]
Il punto \(2\) è di accumulazione per \(A\). Infatti,
per ogni \(\delta>0\) il numero \(a=2+\delta/2\)
appartiene ad \(A\) e soddisfa
\[
  0<|a-2|=\frac{\delta}{2}<\delta.
\]
Più in generale, ogni punto di \(\overline{\R}\) è di
accumulazione per \(A\): per ogni \(t\in\R\) e
\(\delta>0\), almeno uno dei punti \(t-\delta/2\)
e \(t+\delta/2\) appartiene ad \(A\) ed è distinto
da \(t\). Inoltre, \(A\) è illimitato nei due versi,
quindi \(\pm\infty\) sono di accumulazione per la
\cref{prop:lez08-accumulazione-limitatezza}.

\emph{Verifica delle disuguaglianze.}
Fissato \(\varepsilon>0\), occorre trovare
\(\delta_\varepsilon>0\) tale che
\[
  \left|\frac{2x-1}{x+2}-\frac34\right|<\varepsilon
  \qquad\forall x\in A:\ 0<|x-2|<\delta_\varepsilon.
\]
Equivale al sistema
\begin{equation*}
  \begin{cases}
    \displaystyle\frac34-\varepsilon
      <\frac{2x-1}{x+2}, &\text{(I)},\\[7pt]
    \displaystyle\frac{2x-1}{x+2}
      <\frac34+\varepsilon, &\text{(II)}.
  \end{cases}
\end{equation*}
Si cercherà un raggio \(\delta_\varepsilon\le1\),
in modo che \(|x-2|<\delta_\varepsilon\) implichi
\(x>1\), e quindi \(x+2>0\). In tale intorno il
denominatore è positivo e si possono risolvere le
disuguaglianze senza considerare il ramo \(x<-2\)
del dominio.

Per la disuguaglianza (I), moltiplicando per
\(4(x+2)>0\) si ottiene
\begin{align*}
  (3-4\varepsilon)(x+2)&<4(2x-1)\\
  3x-4\varepsilon x+6-8\varepsilon&<8x-4\\
  10-8\varepsilon&<(5+4\varepsilon)x\\
  x&>\frac{10-8\varepsilon}{5+4\varepsilon}.
\end{align*}
Poiché \(5+4\varepsilon>0\),
\[
  2-\frac{10-8\varepsilon}{5+4\varepsilon}
  =\frac{16\varepsilon}{5+4\varepsilon}>0.
\]
Per la disuguaglianza (II), con lo stesso vincolo \(x+2>0\),
si ha
\begin{align*}
  4(2x-1)&<(3+4\varepsilon)(x+2)\\
  8x-4&<3x+4\varepsilon x+6+8\varepsilon\\
  (5-4\varepsilon)x&<10+8\varepsilon.
\end{align*}
Per risolvere quest'ultima disuguaglianza occorre
considerare il segno di \(5-4\varepsilon\):
\begin{enumerate}
  \item Se \(0<\varepsilon<5/4\), si ottiene
  \[
    x<\frac{10+8\varepsilon}{5-4\varepsilon},
  \]
  e la soglia è maggiore di \(2\), poiché
  \[
    \frac{10+8\varepsilon}{5-4\varepsilon}-2
    =\frac{16\varepsilon}{5-4\varepsilon}>0.
  \]
  \item Se \(\varepsilon=5/4\), la disuguaglianza si
  riduce a \(0<20\), verificata per ogni \(x>1\).
  \item Se \(\varepsilon>5/4\), il coefficiente è
  negativo e il verso si inverte:
  \[
    x>\frac{10+8\varepsilon}{5-4\varepsilon}.
  \]
  La soglia è negativa, quindi la condizione è
  soddisfatta in particolare per ogni \(x>1\).
\end{enumerate}

\emph{Scelta del raggio e conclusione.}
Per scegliere un raggio valido per ogni
\(\varepsilon>0\), si ponga
\[
  \eta:=\min\{\varepsilon,1\}.
\]
Poiché \(0<\eta\le1<5/4\), è sufficiente
soddisfare entrambe le disuguaglianze con \(\eta\)
al posto di \(\varepsilon\). Le soglie appena calcolate
permettono di scegliere
\[
  \delta_\varepsilon:=\min\left\{
    1,\frac{16\eta}{5+4\eta},
      \frac{16\eta}{5-4\eta}
  \right\}>0.
\]
Se \(x\in A\) e \(0<|x-2|<\delta_\varepsilon\), allora
\[
  2-\delta_\varepsilon<x<2+\delta_\varepsilon,
\]
e, per la scelta del raggio,
\[
  x>\frac{10-8\eta}{5+4\eta},
  \qquad
  x<\frac{10+8\eta}{5-4\eta}.
\]
Sono quindi verificate entrambe le disuguaglianze
con \(\eta\) al posto di \(\varepsilon\), e ne segue
\[
  \left|\frac{2x-1}{x+2}-\frac34\right|
  <\eta\le\varepsilon.
\]
Poiché \(\varepsilon>0\) era arbitrario, per la
\cref{def:lez08-limite} si conclude
\[
  \lim_{x\to2}\frac{2x-1}{x+2}=\frac34.
\]
\end{proof}

\begin{exercise}[Verifica di un limite infinito]
\label{ex:lez08-razionale-meno-infinito}
Verificare mediante la \cref{def:lez08-limite} che
\[
  \lim_{x\to-4}\frac{-2}{(x+4)^2}=-\infty.
\]
\end{exercise}

\begin{proof}[Soluzione dell'esercizio]
\emph{Dominio e punto di accumulazione.}
La funzione
\[
  f(x):=\frac{-2}{(x+4)^2}
\]
ha dominio
\[
  A=\R\setminus\{-4\}
   =(-\infty,-4)\cup(-4,+\infty).
\]
Per ogni \(\delta>0\), il punto
\(a=-4+\delta/2\) appartiene ad \(A\) e verifica
\(0<|a+4|<\delta\). Pertanto \(-4\) è di
accumulazione per \(A\), secondo la
\cref{def:lez07-accumulazione}.

\emph{Verifica della definizione.}
Per dimostrare che il limite vale \(-\infty\),
per ogni \(K\in\R\) si deve trovare \(\delta_K>0\)
tale che
\[
  \frac{-2}{(x+4)^2}<K
  \qquad\forall x\in A:\ 0<|x+4|<\delta_K.
\]
Poiché \(x\in A\), si ha \(x\ne-4\) e quindi
\((x+4)^2>0\). Moltiplicando per tale numero,
il verso della disuguaglianza rimane invariato:
\[
  \frac{-2}{(x+4)^2}<K
  \quad\Longleftrightarrow\quad
  -2<K(x+4)^2.
\]

Se \(K\ge0\), quest'ultima disuguaglianza è vera
per ogni \(x\in A\), poiché il membro destro è
non negativo. Basta quindi scegliere, ad esempio,
\(\delta_K=1\).

Se \(K<0\), dividendo per \(K\) si inverte il verso:
\[
  (x+4)^2<\frac{-2}{K}
  \quad\Longleftrightarrow\quad
  |x+4|<\sqrt{\frac{-2}{K}}.
\]
Poiché \(-2/K>0\), è sufficiente scegliere
\[
  \delta_K:=\sqrt{\frac{-2}{K}}>0.
\]
In entrambi i casi esiste il raggio richiesto.
Pertanto, per la \cref{def:lez08-limite}, si conclude
\[
  \lim_{x\to-4}\frac{-2}{(x+4)^2}=-\infty.
\]
\end{proof}

\chapter[Lezione 9 --- 16 ottobre 2026]{Lezione 9: limiti all'infinito, composizione e applicazioni inverse}
\label{chap:lezione-09}

\section[Verifica di un limite all'infinito]{Verifica di un limite per \(x\to+\infty\)}
\label{sec:lez09-limite-infinito}

Si completa lo studio degli esempi della \cref{sec:lez08-esercizi-limiti}
considerando un limite finito per \(x\to+\infty\).

\begin{exercise}[Verifica di un limite all'infinito]
\label{ex:lez09-limite-razionale-infinito}
Verificare mediante la \cref{def:lez08-limite} che
\[
  \lim_{x\to+\infty}\frac{3x-1}{x+2}=3.
\]
\end{exercise}

\begin{proof}[Soluzione dell'esercizio]
\emph{Dominio e punto di accumulazione.}
La funzione
\[
  f(x):=\frac{3x-1}{x+2}
\]
ha dominio
\[
  A=\R\setminus\{-2\}
   =(-\infty,-2)\cup(-2,+\infty).
\]
Poiché \(A\) non è superiormente limitato,
\(+\infty\) è un punto di accumulazione per \(A\)
in virtù della \cref{prop:lez08-accumulazione-limitatezza}.

\emph{Verifica della definizione.}
Fissato \(\varepsilon>0\), occorre trovare
\(k_\varepsilon\in\R\) tale che
\[
  \left|\frac{3x-1}{x+2}-3\right|<\varepsilon
  \qquad\forall x\in A\cap(k_\varepsilon,+\infty).
\]
Per ogni \(x\in A\), si ha
\[
  \left|\frac{3x-1}{x+2}-3\right|
  =\left|\frac{-7}{x+2}\right|
  =\frac7{|x+2|}.
\]
Se \(x>-2\), allora \(|x+2|=x+2\), e la condizione
\[
  \frac7{x+2}<\varepsilon
\]
è equivalente a \(x>-2+7/\varepsilon\). È pertanto
sufficiente scegliere
\[
  k_\varepsilon:=-2+\frac7\varepsilon.
\]
Infatti, per ogni \(x\in A\cap(k_\varepsilon,+\infty)\)
si ha \(x>-2\) e
\[
  \left|f(x)-3\right|
  =\frac7{x+2}<\varepsilon.
\]
Poiché \(\varepsilon>0\) era arbitrario, per la
\cref{def:lez08-limite} si conclude
\[
  \lim_{x\to+\infty}\frac{3x-1}{x+2}=3.
\]
\end{proof}

\clearpage

\begin{figure}[htbp]
  \centering
  \begin{tikzpicture}[unipd/diagram,scale=0.95]
  \draw[unipd/directed edge] (-4.5,0) -- (6.2,0);
  \node[below] at (6.1,-.02) {\(x\)};
  \coordinate (excluded) at (-2.6,0);
  \coordinate (threshold) at (.65,0);
  \draw[unipd-accent,line width=2.7pt,opacity=.45] (threshold) -- (5.85,0);
  \draw[unipd-accent,line width=.75pt,dashed] (.65,0) -- (.65,1.05);
  \draw[unipd-primary,line width=.75pt] (-2.6,-.10) -- (-2.6,.10);
  \draw[unipd-primary,line width=.75pt] (.65,-.10) -- (.65,.10);
  \fill[white] (-2.6,0) circle (2.9pt);
  \draw[unipd-primary,line width=.75pt] (-2.6,0) circle (2.9pt);
  \fill[white] (.65,0) circle (2.9pt);
  \draw[unipd-accent,line width=.75pt] (.65,0) circle (2.9pt);
  \node[below=6pt] at (-2.6,0) {\(-2\)};
  \node[below=6pt] at (.65,0) {\(k_\varepsilon=-2+\dfrac7\varepsilon\)};
  \node[above,text=unipd-primary] at (-2.6,.18) {punto escluso};
  \node[above,text=unipd-accent!80!black] at (3.15,.18) {\(x>k_\varepsilon\)};
\end{tikzpicture}

  \caption{La soglia \(k_\varepsilon=-2+7/\varepsilon\)
  individua un intorno di \(+\infty\) nel quale
  \(x+2>0\) e \(|f(x)-3|<\varepsilon\).}
  \label{fig:lez09-limite-razionale-infinito}
\end{figure}

\section{Composizione di applicazioni}
\label{sec:lez09-composizione}

\begin{definition}[Applicazione composta]
\label{def:lez09-composizione}
Siano \(A,B,C,D\) insiemi non vuoti e siano
\[
  \varphi\colon A\longrightarrow B,
  \qquad \psi\colon C\longrightarrow D,
\]
tali che \(\operatorname{Im}(\varphi)\subseteq C\).
Si definisce \emph{applicazione composta} di \(\varphi\)
e \(\psi\) l'applicazione
\[
  \psi\circ\varphi\colon A\longrightarrow D,
  \qquad a\longmapsto\psi\bigl(\varphi(a)\bigr).
\]
\end{definition}

L'ordine della composizione è indicato nella
\cref{fig:lez09-composizione}: si applica prima \(\varphi\)
e successivamente \(\psi\).

\begin{figure}[htbp]
  \centering
  \begin{tikzpicture}[unipd/diagram]
  \foreach \pos/\set in {0/A,4.3/C,8.6/D} {
    \draw[unipd/edge] (\pos,0) ellipse (.83 and 1.0);
    \node at (\pos,1.26) {\(\set\)};
    \fill[unipd-primary] (\pos,0) circle (2pt);
  }
  \node[below=7pt] at (0,0) {\(a\)};
  \node[below=7pt] at (4.3,0) {\(\varphi(a)\)};
  \node[below=7pt] at (8.6,0) {\(\psi(\varphi(a))\)};
  \draw[unipd/directed edge] (.25,.20) -- (4.05,.20)
    node[midway,above=3pt] {\(\varphi\)};
  \draw[unipd/directed edge] (4.55,.20) -- (8.35,.20)
    node[midway,above=3pt] {\(\psi\)};
  \draw[unipd-accent,dashed,-{Latex[length=2.2mm]},line width=.85pt]
    (.08,1.46) to[bend left=22]
    node[midway,above=1pt] {\(\psi\circ\varphi\)} (8.52,1.46);
  \node[below=7pt,text=unipd-primary] at (4.3,-1.06)
    {\(\operatorname{Im}(\varphi)\subseteq C\)};
\end{tikzpicture}

  \caption{Composizione \(\psi\circ\varphi\), con
  \(\operatorname{Im}(\varphi)\subseteq D(\psi)=C\).}
  \label{fig:lez09-composizione}
\end{figure}

\begin{example}[L'ordine della composizione]
\label{ex:lez09-ordine-composizione}
Si considerino le applicazioni
\[
  \varphi\colon\R\longrightarrow\R,
  \quad x\longmapsto x^2,
  \qquad
  \psi\colon\R\longrightarrow\R,
  \quad x\longmapsto x+1.
\]
Entrambe le composizioni sono definite e, per ogni \(x\in\R\),
\[
  (\psi\circ\varphi)(x)=\psi(x^2)=x^2+1,
  \qquad
  (\varphi\circ\psi)(x)=\varphi(x+1)=(x+1)^2.
\]
Per \(x=10\), in particolare,
\[
  (\psi\circ\varphi)(10)=101,
  \qquad (\varphi\circ\psi)(10)=121.
\]
Pertanto \(\psi\circ\varphi\ne\varphi\circ\psi\):
la composizione non è, in generale, commutativa.
\end{example}

L'esistenza delle due composizioni richiede inoltre
condizioni distinte, come mostra l'esempio seguente.

\clearpage

\begin{example}[Composizione non definita]
\label{ex:lez09-composizione-non-definita}
Siano
\[
  \varphi\colon\N\longrightarrow\Q,
  \quad n\longmapsto\frac n2,
  \qquad
  \psi\colon\N\longrightarrow\R,
  \quad n\longmapsto n^2.
\]
Poiché \(\operatorname{Im}(\psi)\subseteq\N=D(\varphi)\),
la composizione \(\varphi\circ\psi\) è definita e
\[
  (\varphi\circ\psi)(n)=\frac{n^2}{2}
  \qquad\forall n\in\N.
\]
Viceversa, \(\varphi(1)=1/2\notin\N=D(\psi)\);
pertanto \(\operatorname{Im}(\varphi)\not\subseteq D(\psi)\)
e \(\psi\circ\varphi\) non è definita su tutto \(\N\).
\end{example}

\section{Applicazione identità}
\label{sec:lez09-identita}

\begin{definition}[Applicazione identità]
\label{def:lez09-identita}
Sia \(A\) un insieme non vuoto. Si chiama
\emph{applicazione identità} su \(A\) l'applicazione
\[
  \operatorname{id}_A\colon A\longrightarrow A,
  \qquad a\longmapsto a.
\]
\end{definition}

Per \(A=\R\), il grafico definito nella
\cref{def:lez07-grafico} è
\[
  G(\operatorname{id}_{\R})
  =\{(a,a):a\in\R\}.
\]

\begin{proposition}[Identità e composizione]
\label{prop:lez09-identita-composizione}
Sia \(\varphi\colon A\longrightarrow B\) un'applicazione
tra insiemi non vuoti. Allora
\[
  \varphi\circ\operatorname{id}_A=\varphi,
  \qquad
  \operatorname{id}_B\circ\varphi=\varphi.
\]
\end{proposition}

\begin{proof}
Per ogni \(a\in A\), applicando le
\cref{def:lez09-composizione,def:lez09-identita}, si ha
\[
  (\varphi\circ\operatorname{id}_A)(a)
  =\varphi(\operatorname{id}_A(a))=\varphi(a),
\]
\[
  (\operatorname{id}_B\circ\varphi)(a)
  =\operatorname{id}_B(\varphi(a))=\varphi(a).
\]
Le due composizioni hanno dominio \(A\) e codominio \(B\),
come \(\varphi\), dunque valgono entrambe le identità.
\end{proof}

Considerando la stessa legge di \(\varphi\) con codominio
\(\operatorname{Im}(\varphi)\), si ottiene l'applicazione
\(\widetilde{\varphi}\colon A\to\operatorname{Im}(\varphi)\).
Per la \cref{prop:lez09-identita-composizione} vale anche
\[
  \operatorname{id}_{\operatorname{Im}(\varphi)}
  \circ\widetilde{\varphi}=\widetilde{\varphi}.
\]

\begin{figure}[H]
  \centering
  \begin{tikzpicture}[unipd/diagram]
  \foreach \yy in {0,-2.05} {
    \foreach \xx in {0,3.35,6.7} {
      \draw[unipd/edge] (\xx,\yy) ellipse (.65 and .65);
      \fill[unipd-primary] (\xx,\yy) circle (2pt);
    }
  }
  \node at (0,.9) {\(A\)};
  \node at (3.35,.9) {\(A\)};
  \node at (6.7,.9) {\(B\)};
  \node at (0,-1.15) {\(A\)};
  \node at (3.35,-1.15) {\(B\)};
  \node at (6.7,-1.15) {\(B\)};
  \draw[unipd/directed edge] (.15,.12) -- (3.2,.12)
    node[midway,above] {\(\operatorname{id}_A\)};
  \draw[unipd/directed edge] (3.5,.12) -- (6.55,.12)
    node[midway,above] {\(\varphi\)};
  \draw[unipd/directed edge] (.15,-1.93) -- (3.2,-1.93)
    node[midway,above] {\(\varphi\)};
  \draw[unipd/directed edge] (3.5,-1.93) -- (6.55,-1.93)
    node[midway,above] {\(\operatorname{id}_B\)};
  \foreach \xx in {0,3.35} {
    \node[below=5pt] at (\xx,0) {\(a\)};
  }
  \node[below=5pt] at (6.7,0) {\(\varphi(a)\)};
  \node[below=5pt] at (0,-2.05) {\(a\)};
  \foreach \xx in {3.35,6.7} {
    \node[below=5pt] at (\xx,-2.05) {\(\varphi(a)\)};
  }
\end{tikzpicture}

  \caption{L'applicazione identità è elemento neutro
  a destra e a sinistra per la composizione.}
  \label{fig:lez09-identita-composizione}
\end{figure}

\clearpage

\section{Applicazione inversa}
\label{sec:lez09-inversa}

Per invertire un'applicazione occorre poter risalire
in modo univoco da ogni elemento dell'immagine al suo
antecedente. La \cref{def:lez06-iniettiva} assicura tale unicità.

\begin{example}[Ostacolo all'esistenza dell'inversa]
\label{ex:lez09-non-iniettiva-inversa}
Per l'applicazione
\(\varphi\colon\R\to\R\), \(\varphi(x)=x^2\), si ha
\[
  \varphi(-1)=\varphi(1)=1.
\]
La \cref{def:lez06-controimmagine} fornisce
\[
  \varphi^{-1}(\{1\})=\{-1,1\}.
\]
Il valore \(1\) possiede quindi due antecedenti distinti:
non è possibile definirne un'inversa che associ a \(1\)
un unico elemento di \(\R\).
\end{example}

\begin{figure}[htbp]
  \centering
  \begin{tikzpicture}[unipd/diagram]
  \draw[unipd/edge] (0,0) ellipse (1 and 1.35);
  \draw[unipd/edge] (4.5,0) ellipse (.9 and 1.35);
  \node at (0,1.65) {\(\R\)};
  \node at (4.5,1.65) {\(\R\)};
  \fill[unipd-primary] (0,.65) circle (2pt);
  \fill[unipd-primary] (0,-.65) circle (2pt);
  \fill[unipd-primary] (4.5,0) circle (2pt);
  \node[left=5pt] at (0,.65) {\(-1\)};
  \node[left=5pt] at (0,-.65) {\(1\)};
  \node[right=5pt] at (4.5,0) {\(1\)};
  \draw[unipd/directed edge] (.15,.62) -- (4.35,.04);
  \draw[unipd/directed edge] (.15,-.62) -- (4.35,-.04);
  \node at (2.25,1.05) {\(\varphi(x)=x^2\)};
  \node[align=center,text=unipd-accent!80!black] at (2.25,-1.65)
    {due preimmagini dello stesso valore};
\end{tikzpicture}

  \caption{La funzione \(\varphi(x)=x^2\) non è iniettiva:
  \(-1\) e \(1\) hanno la stessa immagine.}
  \label{fig:lez09-inversa-non-iniettiva}
\end{figure}

\begin{definition}[Applicazione inversa]
\label{def:lez09-inversa}
Siano \(A,B\) insiemi non vuoti e sia
\(\varphi\colon A\longrightarrow B\) iniettiva.
Si chiama \emph{applicazione inversa} di \(\varphi\)
l'applicazione
\[
  \varphi^{-1}\colon\operatorname{Im}(\varphi)
  \longrightarrow A,
  \qquad b\longmapsto a,
\]
dove \(a\in A\) è l'unico elemento tale che
\(\varphi(a)=b\).
\end{definition}

L'esistenza di \(a\) segue dalla
\cref{def:lez06-immagine}, mentre l'unicità segue dalla
\cref{def:lez06-iniettiva}. In particolare,
\(D(\varphi^{-1})=\operatorname{Im}(\varphi)\) e
\(C(\varphi^{-1})=A\).

\begin{figure}[H]
  \centering
  \begin{tikzpicture}[unipd/diagram]
  \foreach \xx/\set in {0/A,4.1/{\operatorname{Im}(\varphi)},8.2/A} {
    \draw[unipd/edge] (\xx,0) ellipse (.75 and 1);
    \node at (\xx,1.32) {\(\set\)};
    \fill[unipd-primary] (\xx,0) circle (2pt);
  }
  \node[below=7pt] at (0,0) {\(a\)};
  \node[below=7pt] at (4.1,0) {\(\varphi(a)\)};
  \node[below=7pt] at (8.2,0) {\(a\)};
  \draw[unipd/directed edge] (.2,.19) -- (3.9,.19)
    node[midway,above=3pt] {\(\varphi\)};
  \draw[unipd/directed edge] (4.3,.19) -- (8,.19)
    node[midway,above=3pt] {\(\varphi^{-1}\)};
  \draw[unipd-accent,dashed,-{Latex[length=2.2mm]},line width=.85pt]
    (.08,-1.2) to[bend right=20]
    node[midway,below] {\(\varphi^{-1}\circ\varphi=\operatorname{id}_A\)}
    (8.12,-1.2);
\end{tikzpicture}

  \caption{L'applicazione inversa riporta
  \(\varphi(a)\) ad \(a\), realizzando
  \(\varphi^{-1}\circ\varphi=\operatorname{id}_A\).}
  \label{fig:lez09-composizione-inversa}
\end{figure}

\begin{proposition}[Composizione con l'applicazione inversa]
\label{prop:lez09-inversa-composizione}
Sia \(\varphi\colon A\longrightarrow B\) iniettiva.
Con la notazione introdotta nella \cref{sec:lez09-identita},
si ha
\[
  \varphi^{-1}\circ\varphi=\operatorname{id}_A,
  \qquad
  \widetilde{\varphi}\circ\varphi^{-1}
  =\operatorname{id}_{\operatorname{Im}(\varphi)}.
\]
\end{proposition}

\begin{proof}
Le due composizioni sono definite rispettivamente
da \(A\) in \(A\) e da \(\operatorname{Im}(\varphi)\)
in sé.

Per ogni \(a\in A\), dalla \cref{def:lez09-inversa}
si ottiene
\[
  (\varphi^{-1}\circ\varphi)(a)
  =\varphi^{-1}(\varphi(a))=a
  =\operatorname{id}_A(a).
\]
Analogamente, per ogni \(b\in\operatorname{Im}(\varphi)\),
\[
  (\widetilde{\varphi}\circ\varphi^{-1})(b)
  =\widetilde{\varphi}(\varphi^{-1}(b))=b
  =\operatorname{id}_{\operatorname{Im}(\varphi)}(b).
\]
Seguono entrambe le uguaglianze.
\end{proof}

\begin{example}[Esponenziale e logaritmo]
\label{ex:lez09-inversa-esponenziale}
Si considerino le applicazioni inverse
\[
  \varphi\colon\R\longrightarrow(0,+\infty),
  \quad x\longmapsto e^x,
\]
\[
  \varphi^{-1}\colon(0,+\infty)\longrightarrow\R,
  \quad y\longmapsto\log y.
\]
Per ogni \(x\in\R\), si ha
\[
  (\varphi^{-1}\circ\varphi)(x)
  =\log(e^x)=x.
\]
Analogamente, per ogni \(y\in(0,+\infty)\),
\[
  (\varphi\circ\varphi^{-1})(y)
  =e^{\log y}=y.
\]
Pertanto
\[
  \varphi^{-1}\circ\varphi=\operatorname{id}_{\R},
  \qquad
  \varphi\circ\varphi^{-1}
  =\operatorname{id}_{(0,+\infty)}.
\]
Le due composizioni hanno domini differenti:
\[
  D(\varphi^{-1}\circ\varphi)=\R,
  \qquad
  D(\varphi\circ\varphi^{-1})=(0,+\infty).
\]
\end{example}

\chapter[Lezione 10 --- 20 ottobre 2026]{Lezione 10: applicazioni inverse e monotonia}
\label{chap:lezione-10}

\section{Dominio, immagine e grafico dell'applicazione inversa}
\label{sec:lez10-proprieta-inversa}

La \cref{def:lez09-inversa} permette di precisare il rapporto
tra dominio e immagine di un'applicazione iniettiva e della sua inversa.

\begin{proposition}[Dominio e immagine dell'applicazione inversa]
\label{prop:lez10-dominio-immagine-inversa}
Siano \(A,B\) insiemi non vuoti e
\(\varphi\colon A\longrightarrow B\) un'applicazione iniettiva.
Allora
\[
  D(\varphi^{-1})=\operatorname{Im}(\varphi),
  \qquad
  \operatorname{Im}(\varphi^{-1})=D(\varphi).
\]
\end{proposition}

\begin{proof}
\emph{Prima uguaglianza.}
Segue direttamente dalla \cref{def:lez09-inversa}.

\emph{Seconda uguaglianza.}
Poiché \(C(\varphi^{-1})=A=D(\varphi)\),
per la \cref{def:lez06-immagine} si ha
\[
  \operatorname{Im}(\varphi^{-1})\subseteq D(\varphi).
\]
Viceversa, sia \(a_0\in D(\varphi)\).
Ponendo \(b_0:=\varphi(a_0)\), si ottiene
\(b_0\in\operatorname{Im}(\varphi)=D(\varphi^{-1})\).
Dalla \cref{def:lez09-inversa} segue
\[
  \varphi^{-1}(b_0)=a_0,
\]
quindi \(a_0\in\operatorname{Im}(\varphi^{-1})\).
Pertanto \(D(\varphi)\subseteq\operatorname{Im}(\varphi^{-1})\),
da cui la seconda uguaglianza.
\end{proof}

\begin{example}[Esponenziale e logaritmo]
\label{ex:lez10-dominio-immagine-logaritmo}
L'applicazione
\[
  \varphi\colon\R\longrightarrow(0,+\infty),
  \qquad x\longmapsto e^x,
\]
è invertibile e la sua inversa è
\[
  \varphi^{-1}\colon(0,+\infty)\longrightarrow\R,
  \qquad y\longmapsto\log y.
\]
In accordo con la \cref{prop:lez10-dominio-immagine-inversa},
\[
  D(\varphi^{-1})=(0,+\infty),
  \qquad \operatorname{Im}(\varphi^{-1})=\R.
\]
\end{example}

Se \(A,B\subseteq\R\), il grafico dell'inversa si ottiene
scambiando le coordinate dei punti del grafico originario.
Infatti, per la \cref{def:lez07-grafico},
\[
\begin{aligned}
  G(\varphi)
    &=\{(a,\varphi(a)):a\in A\},\\
  G(\varphi^{-1})
    &=\{(b,\varphi^{-1}(b)):b\in\operatorname{Im}(\varphi)\}\\
    &=\{(\varphi(a),a):a\in A\}.
\end{aligned}
\]

\clearpage

I due grafici sono quindi simmetrici rispetto alla retta
\(y=x\), rappresentata nella \cref{fig:lez10-grafico-inversa}.

\begin{figure}[htbp]
  \centering
  \begin{tikzpicture}[unipd/diagram,x=1.12cm,y=1.12cm]
  \clip (-2.35,-2.35) rectangle (4.35,4.25);
  \draw[->,unipd-muted] (-2.25,0) -- (4.23,0)
    node[below right] {\(x\)};
  \draw[->,unipd-muted] (0,-2.25) -- (0,4.12)
    node[above left] {\(y\)};
  \draw[unipd-muted,densely dashed,thick] (-2.15,-2.15) -- (4.05,4.05)
    node[above left] {\(y=x\)};
  \draw[very thick,unipd-primary,domain=-2.2:1.39,samples=110]
    plot (\x,{exp(\x)});
  \draw[very thick,unipd-accent,domain=.11:4.13,samples=110]
    plot (\x,{ln(\x)});
  \fill[unipd-primary] (0,1) circle (1.9pt);
  \fill[unipd-accent] (1,0) circle (1.9pt);
  \node[left=4pt] at (0,1) {\((0,1)\)};
  \node[below=5pt] at (1,0) {\((1,0)\)};
  \node[text=unipd-primary,above left] at (1.32,3.74) {\(y=e^x\)};
  \node[text=unipd-accent,above left] at (3.65,1.28) {\(y=\log x\)};
\end{tikzpicture}

  \caption{I grafici di \(e^x\) e \(\log x\) sono simmetrici
  rispetto alla retta \(y=x\).}
  \label{fig:lez10-grafico-inversa}
\end{figure}

\section{Monotonia delle funzioni}
\label{sec:lez10-monotonia}

Siano \(A\subseteq\R\) non vuoto,
\(\varphi\colon A\longrightarrow\R\) e
\(S\subseteq A\) non vuoto. Le definizioni seguenti
si applicano, in particolare, quando \(S=I\) è un intervallo,
nel senso della \cref{def:intervallo-retta-reale}.

\begin{definition}[Crescenza, decrescenza e monotonia]
\label{def:lez10-monotonia}
Per ogni \(x_1,x_2\in S\) con \(x_1<x_2\),
si dice che \(\varphi\) è su \(S\):
\begin{enumerate}
  \item \emph{strettamente crescente} se
  \(\varphi(x_1)<\varphi(x_2)\);
  \item \emph{debolmente crescente} se
  \(\varphi(x_1)\le\varphi(x_2)\);
  \item \emph{strettamente decrescente} se
  \(\varphi(x_1)>\varphi(x_2)\);
  \item \emph{debolmente decrescente} se
  \(\varphi(x_1)\ge\varphi(x_2)\).
\end{enumerate}
Si dice che \(\varphi\) è \emph{strettamente monotona} su \(S\)
se è strettamente crescente oppure strettamente decrescente,
e \emph{debolmente monotona} su \(S\) se è debolmente crescente
oppure debolmente decrescente.
\end{definition}

\begin{example}[Un'unione che non è un intervallo]
\label{ex:lez10-unione-non-intervallo}
L'insieme
\[
  J=(1,2)\cup(2,4)
\]
non è un intervallo, benché sia unione di due intervalli.
Infatti, scegliendo
\[
  x_1=\frac32\in J,
  \qquad x_2=3\in J,
\]
si ha \(x_1<2<x_2\) ma \(2\notin J\),
in contrasto con la \cref{def:intervallo-retta-reale}.
\end{example}

\clearpage

\begin{proposition}[Monotonia della funzione opposta]
\label{prop:lez10-monotonia-opposta}
Siano \(I\subseteq\R\) un intervallo non vuoto e
\(\varphi\colon I\longrightarrow\R\).

Se \(\varphi\) è debolmente crescente su \(I\),
allora \(-\varphi\) è debolmente decrescente su \(I\).
Se \(\varphi\) è debolmente decrescente su \(I\),
allora \(-\varphi\) è debolmente crescente su \(I\).
Le stesse proprietà valgono per la monotonia stretta.
\end{proposition}

\begin{proof}
Siano \(x_1,x_2\in I\) con \(x_1<x_2\).

\emph{Caso crescente.}
Se \(\varphi(x_1)\le\varphi(x_2)\), cambiando
segno a entrambi i membri si inverte il verso
della disuguaglianza e si ottiene
\[
  (-\varphi)(x_1)=-\varphi(x_1)
  \ge-\varphi(x_2)=(-\varphi)(x_2).
\]
Pertanto \(-\varphi\) è debolmente decrescente.

\emph{Caso decrescente.}
Se \(\varphi(x_1)\ge\varphi(x_2)\),
con lo stesso procedimento si ottiene
\[
  (-\varphi)(x_1)=-\varphi(x_1)
  \le-\varphi(x_2)=(-\varphi)(x_2).
\]
Pertanto \(-\varphi\) è debolmente crescente.
Se le disuguaglianze iniziali sono strette,
restano strette anche quelle ottenute moltiplicando per \(-1\).
\end{proof}

\begin{proposition}[Monotonia del prodotto di funzioni non negative]
\label{prop:lez10-monotonia-prodotto}
Siano \(I\subseteq\R\) un intervallo non vuoto e
\(\varphi,\psi\colon I\longrightarrow\R\), con
\[
  \operatorname{Im}(\varphi)\subseteq[0,+\infty),
  \qquad
  \operatorname{Im}(\psi)\subseteq[0,+\infty).
\]
Se \(\varphi\) e \(\psi\) sono entrambe debolmente crescenti
su \(I\), allora il prodotto \(\varphi\psi\) è debolmente
crescente su \(I\).

Se \(\varphi\) e \(\psi\) sono entrambe debolmente decrescenti
su \(I\), allora il prodotto \(\varphi\psi\) è debolmente
decrescente su \(I\).
\end{proposition}

\begin{proof}
Siano \(x_1,x_2\in I\) con \(x_1<x_2\).

\emph{Caso crescente.}
Per ipotesi,
\[
  0\le\varphi(x_1)\le\varphi(x_2),
  \qquad 0\le\psi(x_1)\le\psi(x_2).
\]
Ne segue che
\[
  \varphi(x_1)\psi(x_1)
  \le\varphi(x_2)\psi(x_1)
  \le\varphi(x_2)\psi(x_2),
\]
quindi \(\varphi\psi\) è debolmente crescente.

\emph{Caso decrescente.}
Per ipotesi,
\[
  \varphi(x_1)\ge\varphi(x_2)\ge0,
  \qquad \psi(x_1)\ge\psi(x_2)\ge0.
\]
Ne segue che
\[
  \varphi(x_1)\psi(x_1)
  \ge\varphi(x_2)\psi(x_1)
  \ge\varphi(x_2)\psi(x_2),
\]
quindi \(\varphi\psi\) è debolmente decrescente.
\end{proof}

Quando una funzione è crescente e l'altra decrescente,
non si può dedurre, in generale, la monotonia del prodotto.
Ad esempio, su \(I=[0,1]\) le funzioni non negative
\(\varphi(x)=x\) e \(\psi(x)=1-x\) sono rispettivamente
crescenti e decrescenti, ma
\[
  (\varphi\psi)(0)=(\varphi\psi)(1)=0,
  \qquad (\varphi\psi)\!\left(\frac12\right)=\frac14,
\]
per cui il prodotto non è monotono su \(I\).

\begin{proposition}[Monotonia stretta e applicazione inversa]
\label{prop:lez10-monotonia-inversa}
Siano \(I\subseteq\R\) un intervallo non vuoto e
\(\varphi\colon I\longrightarrow\R\) strettamente monotona su \(I\).
Allora:
\begin{enumerate}
  \item \(\varphi\) è iniettiva;
  \item l'applicazione inversa
  \(\varphi^{-1}\colon\operatorname{Im}(\varphi)\longrightarrow I\)
  è strettamente monotona su \(\operatorname{Im}(\varphi)\).
  In particolare, se \(\varphi\) è strettamente crescente,
  anche \(\varphi^{-1}\) è strettamente crescente;
  se \(\varphi\) è strettamente decrescente,
  anche \(\varphi^{-1}\) è strettamente decrescente.
\end{enumerate}
\end{proposition}

\begin{proof}[Dimostrazione del primo punto]
Siano \(x_1,x_2\in I\) distinti.
Se \(x_1<x_2\), la stretta monotonia implica
\(\varphi(x_1)<\varphi(x_2)\), nel caso crescente, e
\(\varphi(x_1)>\varphi(x_2)\), nel caso decrescente.
Se \(x_2<x_1\) si procede analogamente.
In entrambi i casi \(\varphi(x_1)\ne\varphi(x_2)\),
dunque \(\varphi\) è iniettiva per la
\cref{def:lez06-iniettiva}.
\end{proof}

La seconda affermazione della
\cref{prop:lez10-monotonia-inversa} è enunciata senza dimostrazione.

\begin{example}[Monotonia del logaritmo]
\label{ex:lez10-logaritmo-monotono}
Per ogni \(a>1\), la funzione
\[
  \varphi\colon\R\longrightarrow(0,+\infty),
  \qquad x\longmapsto a^x,
\]
è strettamente crescente su \(\R\).
La sua inversa
\[
  \varphi^{-1}\colon(0,+\infty)\longrightarrow\R,
  \qquad x\longmapsto\log_a x,
\]
è quindi strettamente crescente su \((0,+\infty)\)
per la \cref{prop:lez10-monotonia-inversa}.
\end{example}

\chapter[Lezione 11 --- 21 ottobre 2026]{Lezione 11: intorni e proprietà dei limiti}
\label{chap:lezione-11}

\section{Intersezioni di intorni}
\label{sec:lez11-intersezione-intorni}

Sia \(x_0\in\overline{\R}\). L'intersezione di due intorni di
\(x_0\) è ancora un intorno di \(x_0\).

Se \(x_0\in\R\), per la \cref{def:lez07-intorno} si ha, per
ogni \(\delta_1,\delta_2>0\),
\[
  D(x_0,\delta_1)\cap D(x_0,\delta_2)
  =D\bigl(x_0,\min\{\delta_1,\delta_2\}\bigr).
\]
Se \(x_0=+\infty\), per la \cref{def:lez07-intorni-infinito},
per ogni \(k_1,k_2\in\R\),
\[
  (k_1,+\infty)\cap(k_2,+\infty)
  =(\max\{k_1,k_2\},+\infty).
\]
Analogamente, se \(x_0=-\infty\),
\[
  (-\infty,k_1)\cap(-\infty,k_2)
  =(-\infty,\min\{k_1,k_2\}).
\]
In tutti i casi, l'intersezione è un intorno del medesimo
punto. Applicando ripetutamente la stessa proprietà, si
ottiene il risultato per l'intersezione di un numero finito di intorni di \(x_0\).

\begin{figure}[htbp]
  \centering
  \begin{tikzpicture}[unipd/diagram]
  \draw[->,unipd-muted] (-.6,-1.05) -- (8,-1.05) node[right] {\(x\)};
  \draw[very thick,unipd-primary,->] (1.45,.85) -- (7.55,.85);
  \draw[unipd-primary,fill=white,thick] (1.45,.85) circle (2.4pt);
  \draw[very thick,unipd-accent,->] (3.45,.10) -- (7.55,.10);
  \draw[unipd-accent,fill=white,thick] (3.45,.10) circle (2.4pt);
  \draw[very thick,unipd-primary,->] (3.45,-.60) -- (7.55,-.60);
  \draw[unipd-primary,fill=white,thick] (3.45,-.60) circle (2.4pt);
  \draw[dashed,unipd-muted] (1.45,-1.13) -- (1.45,1.08);
  \draw[dashed,unipd-muted] (3.45,-1.13) -- (3.45,1.08);
  \node[below=6pt] at (1.45,-1.05) {\(k_1\)};
  \node[below=6pt] at (3.45,-1.05) {\(k_2\)};
  \node[above=4pt,text=unipd-primary] at (5.7,.85) {\((k_1,+\infty)\)};
  \node[above=4pt,text=unipd-accent] at (5.7,.10) {\((k_2,+\infty)\)};
  \node[above=4pt,text=unipd-primary] at (5.7,-.60) {\((\max\{k_1,k_2\},+\infty)\)};
\end{tikzpicture}

  \caption{Intersezione di due intorni di \(+\infty\):
  la semiretta comune inizia dalla maggiore delle due soglie,
  qui con \(k_1<k_2\).}
  \label{fig:lez11-intersezione-intorni-infinito}
\end{figure}

L'affermazione non si estende, in generale, alle intersezioni
infinite, come mostra l'esempio seguente.

\begin{example}[Intersezione di infiniti intorni]
\label{ex:lez11-intersezione-infinita}
Per ogni \(n\in\N\) con \(n\ge1\), si consideri
\[
  A_n:=D\!\left(0,\frac1n\right)
       =\left(-\frac1n,\frac1n\right).
\]
Allora
\[
  B:=\bigcap_{n=1}^{\infty}A_n=\{0\}.
\]
L'insieme \(B\) non è quindi un intorno di \(0\), benché
ciascun \(A_n\) lo sia.
\end{example}

\begin{proof}
Poiché \(0\in A_n\) per ogni \(n\ge1\), si ha
\(0\in B\).

Si consideri l'insieme
\[
  E:=\left\{\frac1n:n\in\N,\ n\ge1\right\}.
\]
Il numero \(0\) è un minorante di \(E\). Inoltre, per
ogni \(\varepsilon>0\), per il \cref{thm:proprieta-archimedea}
esiste \(n\in\N\) tale che \(n>1/\varepsilon\);
pertanto \(1/n<\varepsilon\). Per il
\cref{thm:caratterizzazione-inf}, segue
\[
  \inf E=0.
\]
Dall'\cref{ex:estremi-insieme-opposto} si ricava anche
\[
  \sup(-E)=-\inf E=0.
\]

Sia ora \(x\in B\). Per definizione di intersezione,
\[
  -\frac1n<x<\frac1n
  \qquad\forall n\in\N,\ n\ge1.
\]
Il numero \(x\) è dunque un minorante di \(E\) e un
maggiorante di \(-E\). Ne segue che
\[
  x\le\inf E=0,
  \qquad
  x\ge\sup(-E)=0.
\]
Pertanto \(x=0\), e quindi \(B\subseteq\{0\}\).
Insieme all'inclusione opposta, si conclude
\[
  B=\{0\}.
\]
\end{proof}

\section{Unicità del limite}
\label{sec:lez11-unicita-limite}

La \cref{def:lez08-limite} non permette a una funzione
di ammettere due limiti distinti nello stesso punto di
accumulazione.

\begin{theorem}[Unicità del limite]
\label{thm:lez11-unicita-limite}
Sia \(f\colon A\longrightarrow\R\), con \(A\subseteq\R\)
non vuoto, e sia \(x_0\in\overline{\R}\) un punto di
accumulazione per \(A\). Se
\[
  \lim_{x\to x_0}f(x)=\ell_1,
  \qquad
  \lim_{x\to x_0}f(x)=\ell_2,
\]
con \(\ell_1,\ell_2\in\overline{\R}\), allora
\[
  \ell_1=\ell_2.
\]
\end{theorem}

\begin{proof}
Si procede per assurdo, supponendo
\(\ell_1\ne\ell_2\).

Esistono due intorni \(V_{\ell_1}\) e \(V_{\ell_2}\) dei
rispettivi punti tali che
\[
  V_{\ell_1}\cap V_{\ell_2}=\varnothing.
\]
Infatti, se \(\ell_1,\ell_2\in\R\), si può supporre,
senza perdita di generalità, \(\ell_1<\ell_2\) e porre
\[
  \delta:=\frac{\ell_2-\ell_1}{2}>0,
  \qquad
  V_{\ell_1}:=D(\ell_1,\delta),
  \qquad
  V_{\ell_2}:=D(\ell_2,\delta).
\]
I due intorni sono disgiunti, come illustrato nella
\cref{fig:lez11-intorni-disgiunti}. Se almeno uno dei
due limiti è infinito, si possono scegliere intorni
disgiunti utilizzando la \cref{def:lez07-intorni-infinito}.

Per la \cref{def:lez08-limite}, esistono intorni
\(U_{x_0}\) e \(U'_{x_0}\) di \(x_0\) tali che
\[
  f(x)\in V_{\ell_1}
  \quad\forall x\in(U_{x_0}\cap A)\setminus\{x_0\},
\]
\[
  f(x)\in V_{\ell_2}
  \quad\forall x\in(U'_{x_0}\cap A)\setminus\{x_0\}.
\]
Per la \cref{sec:lez11-intersezione-intorni},
\(W_{x_0}:=U_{x_0}\cap U'_{x_0}\) è ancora un intorno
di \(x_0\). Poiché \(x_0\) è di accumulazione per
\(A\), esiste
\[
  x\in(W_{x_0}\cap A)\setminus\{x_0\}.
\]
Per tale \(x\), entrambe le condizioni sul limite
devono essere verificate, quindi
\[
  f(x)\in V_{\ell_1}\cap V_{\ell_2}=\varnothing,
\]
una contraddizione. Pertanto \(\ell_1=\ell_2\).
\end{proof}

\begin{figure}[htbp]
  \centering
  \begin{tikzpicture}[unipd/diagram,scale=1.1]
  \draw[->,unipd-muted] (-.45,0) -- (8.1,0) node[right] {\(x\)};
  \draw[very thick,unipd-primary] (.7,0) -- (3.5,0);
  \draw[very thick,unipd-accent] (3.5,0) -- (6.3,0);
  \foreach \x in {.7,3.5} {
    \draw[unipd-primary,fill=white,thick] (\x,0) circle (2.4pt);
  }
  \foreach \x in {3.5,6.3} {
    \draw[unipd-accent,fill=white,thick] (\x,0) circle (2.0pt);
  }
  \fill[unipd-primary] (2.1,0) circle (2.1pt);
  \fill[unipd-accent] (4.9,0) circle (2.1pt);
  \node[below=7pt] at (2.1,0) {\(\ell_1\)};
  \node[below=7pt] at (4.9,0) {\(\ell_2\)};
  \node[above=8pt,text=unipd-primary] at (2.1,0) {\(V_{\ell_1}\)};
  \node[above=8pt,text=unipd-accent] at (4.9,0) {\(V_{\ell_2}\)};
  \draw[<->,unipd-muted] (2.1,.7) -- (3.5,.7) node[midway,above] {\(\delta\)};
  \draw[<->,unipd-muted] (3.5,.7) -- (4.9,.7) node[midway,above] {\(\delta\)};
  \draw[dashed,unipd-muted] (3.5,.12) -- (3.5,.85);
\end{tikzpicture}

  \caption{Due intorni circolari disgiunti di
  \(\ell_1<\ell_2\), con raggio
  \(\delta=(\ell_2-\ell_1)/2\).}
  \label{fig:lez11-intorni-disgiunti}
\end{figure}

\begin{exercise}[Intorni disgiunti nella retta reale estesa]
\label{ex:lez11-intorni-disgiunti}
Verificare che si possono scegliere due intorni disgiunti
anche quando \(\ell_1\ne\ell_2\) e almeno uno dei due
appartiene a \(\{-\infty,+\infty\}\), distinguendo
il caso in cui l'altro limite è reale da quello
in cui i due limiti sono infiniti.
\end{exercise}

\begin{proof}[Soluzione dell'esercizio]
Si distinguono tre casi.

\emph{Primo caso: \(\ell_1=+\infty\), \(\ell_2\in\R\).}
Si scelgano gli intorni
\[
  V_{\ell_1}:=(\ell_2+1,+\infty),
  \qquad
  V_{\ell_2}:=D(\ell_2,1)
  =(\ell_2-1,\ell_2+1).
\]
Poiché gli intervalli non hanno punti in comune,
si ha \(V_{\ell_1}\cap V_{\ell_2}=\varnothing\).

\emph{Secondo caso: \(\ell_1\in\R\), \(\ell_2=-\infty\).}
Si scelgano gli intorni
\[
  V_{\ell_1}:=D(\ell_1,1)
  =(\ell_1-1,\ell_1+1),
  \qquad
  V_{\ell_2}:=(-\infty,\ell_1-1).
\]
Anche in questo caso
\(V_{\ell_1}\cap V_{\ell_2}=\varnothing\).

\emph{Terzo caso: \(\ell_1=-\infty\), \(\ell_2=+\infty\).}
Si scelgano gli intorni
\[
  V_{\ell_1}:=(-\infty,0),
  \qquad
  V_{\ell_2}:=(0,+\infty).
\]
Si ha nuovamente
\(V_{\ell_1}\cap V_{\ell_2}=\varnothing\).

Scambiando gli indici si ottengono i casi rimanenti.
Pertanto, per ogni coppia di punti distinti della
retta reale estesa, esistono due intorni disgiunti.
\end{proof}

\section{Permanenza del segno}
\label{sec:lez11-permanenza-segno}

\begin{theorem}[Permanenza del segno]
\label{thm:lez11-permanenza-segno}
Sia \(f\colon A\longrightarrow\R\), con \(A\subseteq\R\)
non vuoto, e sia \(x_0\in\overline{\R}\) un punto di
accumulazione per \(A\).
\begin{enumerate}
  \item Se
  \[
    \lim_{x\to x_0}f(x)=\ell
    \in(0,+\infty)\cup\{+\infty\},
  \]
  allora esistono \(k>0\) e un intorno \(U_{x_0}\)
  di \(x_0\) tali che
  \[
    f(x)>k
    \qquad\forall x\in(U_{x_0}\cap A)\setminus\{x_0\}.
  \]
  \item Se
  \[
    \lim_{x\to x_0}f(x)=\ell
    \in(-\infty,0)\cup\{-\infty\},
  \]
  allora esistono \(k>0\) e un intorno \(U_{x_0}\)
  di \(x_0\) tali che
  \[
    f(x)<-k
    \qquad\forall x\in(U_{x_0}\cap A)\setminus\{x_0\}.
  \]
\end{enumerate}
\end{theorem}

\begin{proof}[Dimostrazione del primo punto]
Si distinguono i casi \(\ell=+\infty\) e
\(\ell\in(0,+\infty)\).

\emph{Caso \(\ell=+\infty\).}
Per la \cref{def:lez08-limite}, per ogni \(K\in\R\)
esiste un intorno \(U_{x_0}\) tale che
\[
  f(x)>K
  \qquad\forall x\in(U_{x_0}\cap A)\setminus\{x_0\}.
\]
Scegliendo \(K=1\), si ottiene la tesi con \(k=1\).

\emph{Caso \(\ell\in(0,+\infty)\).}
Per la \cref{def:lez08-limite}, per ogni
\(\varepsilon>0\) esiste un intorno \(U_{x_0}\)
tale che
\[
  |f(x)-\ell|<\varepsilon
  \qquad\forall x\in(U_{x_0}\cap A)\setminus\{x_0\}.
\]
Scegliendo \(\varepsilon=\ell/2>0\), segue
\[
  -\frac\ell2<f(x)-\ell<\frac\ell2,
\]
e quindi
\[
  f(x)>\frac\ell2>0
  \qquad\forall x\in(U_{x_0}\cap A)\setminus\{x_0\}.
\]
La tesi è verificata ponendo \(k=\ell/2\).
\end{proof}

Il secondo punto si dimostra analogamente, considerando
\(\ell=-\infty\) e \(\ell\in(-\infty,0)\); la verifica
non viene sviluppata in questa sede.

\section{Limitatezza locale}
\label{sec:lez11-limitatezza-locale}

\begin{theorem}[Limitatezza locale]
\label{thm:lez11-limitatezza-locale}
Sia \(f\colon A\longrightarrow\R\), con \(A\subseteq\R\)
non vuoto, e sia \(x_0\in\overline{\R}\) un punto di
accumulazione per \(A\). Se
\[
  \lim_{x\to x_0}f(x)=\ell\in\R,
\]
allora esiste un intorno \(U_{x_0}\) di \(x_0\)
tale che l'insieme
\[
  \{f(x):x\in U_{x_0}\cap A\}
\]
è limitato.
\end{theorem}

La dimostrazione del \cref{thm:lez11-limitatezza-locale}
non viene sviluppata in questa sede.

\chapter[Lezione 12 --- 22 ottobre 2026]{Lezione 12: teoremi del confronto e teorema delle tre funzioni}
\label{chap:lezione-12}

\section{Primo teorema del confronto}
\label{sec:lez12-confronto-primo}

\begin{theorem}[Primo teorema del confronto]
\label{thm:lez12-confronto-primo}
Siano \(f,g\colon A\longrightarrow\R\), con
\(A\subseteq\R\) non vuoto, e sia
\(x_0\in\overline{\R}\) un punto di accumulazione per \(A\).
Valgono le seguenti affermazioni.
\begin{enumerate}
  \item Se
  \[
    \lim_{x\to x_0}f(x)=\ell_1\in\R,
    \qquad
    \lim_{x\to x_0}g(x)=\ell_2\in\R,
    \qquad \ell_1<\ell_2,
  \]
  allora esiste un intorno \(U_{x_0}\) di \(x_0\) tale che
  \[
    f(x)<g(x)
    \qquad\forall x\in(U_{x_0}\cap A)\setminus\{x_0\}.
  \]
  \item Se
  \[
    \lim_{x\to x_0}f(x)=-\infty,
    \qquad
    \lim_{x\to x_0}g(x)=\ell\in\R\cup\{+\infty\},
  \]
  allora esiste un intorno \(U_{x_0}\) di \(x_0\) tale che
  \[
    f(x)<g(x)
    \qquad\forall x\in(U_{x_0}\cap A)\setminus\{x_0\}.
  \]
  \item Se
  \[
    \lim_{x\to x_0}f(x)=\ell\in\R,
    \qquad
    \lim_{x\to x_0}g(x)=+\infty,
  \]
  allora esiste un intorno \(U_{x_0}\) di \(x_0\) tale che
  \[
    f(x)<g(x)
    \qquad\forall x\in(U_{x_0}\cap A)\setminus\{x_0\}.
  \]
\end{enumerate}
\end{theorem}

\begin{proof}[Dimostrazione del primo punto]
Poiché \(\ell_1<\ell_2\), si ponga
\[
  \varepsilon:=\frac{\ell_2-\ell_1}{2}>0.
\]
Per la \cref{def:lez08-limite}, esistono intorni
\(U'_{x_0}\) e \(U''_{x_0}\) di \(x_0\) tali che
\[
  |f(x)-\ell_1|<\varepsilon
  \qquad\forall x\in(U'_{x_0}\cap A)\setminus\{x_0\},
\]
\[
  |g(x)-\ell_2|<\varepsilon
  \qquad\forall x\in(U''_{x_0}\cap A)\setminus\{x_0\}.
\]
Per la \cref{sec:lez11-intersezione-intorni},
\[
  U_{x_0}:=U'_{x_0}\cap U''_{x_0}
\]
è ancora un intorno di \(x_0\). Per ogni
\(x\in(U_{x_0}\cap A)\setminus\{x_0\}\), si ottiene
\[
  f(x)<\ell_1+\varepsilon
  =\frac{\ell_1+\ell_2}{2}
  =\ell_2-\varepsilon<g(x).
\]
\end{proof}

\clearpage

\begin{exercise}[Casi con limiti infiniti]
\label{ex:lez12-confronto-primo-infiniti}
Dimostrare il secondo e il terzo punto del
\cref{thm:lez12-confronto-primo} mediante la
\cref{def:lez08-limite}.
\end{exercise}

\begin{proof}[Soluzione dell'esercizio]
\emph{Secondo punto.}
Se \(\ell\in\R\), dalla \cref{def:lez08-limite},
scegliendo \(\varepsilon=1\), esiste un intorno
\(U'_{x_0}\) nel quale \(g(x)>\ell-1\).
Poiché \(f(x)\to-\infty\), esiste un intorno
\(U''_{x_0}\) nel quale \(f(x)<\ell-1\).
Per ogni \(x\in(U'_{x_0}\cap U''_{x_0}\cap A)
\setminus\{x_0\}\), segue
\[
  f(x)<\ell-1<g(x).
\]
Se \(\ell=+\infty\), esistono intorni
\(U'_{x_0}\) e \(U''_{x_0}\) nei quali valgono
rispettivamente \(f(x)<0\) e \(g(x)>0\).
Per ogni \(x\in(U'_{x_0}\cap U''_{x_0}\cap A)\setminus\{x_0\}\),
si ha
\[
  f(x)<0<g(x).
\]

\emph{Terzo punto.}
Poiché \(f(x)\to\ell\in\R\), scegliendo
\(\varepsilon=1\) esiste un intorno \(U'_{x_0}\)
nel quale \(f(x)<\ell+1\).
Poiché \(g(x)\to+\infty\), esiste un intorno
\(U''_{x_0}\) nel quale \(g(x)>\ell+1\).
Per ogni \(x\in(U'_{x_0}\cap U''_{x_0}\cap A)
\setminus\{x_0\}\), segue
\[
  f(x)<\ell+1<g(x).
\]
In ciascun caso, l'intersezione dei due intorni è un
intorno di \(x_0\) per la \cref{sec:lez11-intersezione-intorni}.
\end{proof}

\begin{example}[Limiti uguali e ordinamento variabile]
\label{ex:lez12-limiti-uguali-ordine}
Si considerino le funzioni su \(\R\)
\[
  f(x)=x^3,
  \qquad g(x)=-x^3.
\]
Entrambe hanno limite \(0\) per \(x\to0\):
per ogni \(\varepsilon>0\), scegliendo
\(\delta=\varepsilon^{1/3}\), da \(0<|x|<\delta\)
segue \(|x^3|<\varepsilon\).
Tuttavia,
\[
  f(x)<g(x)\quad\text{se }x<0,
  \qquad
  f(x)>g(x)\quad\text{se }x>0.
\]
Ogni intorno bucato di \(0\) contiene punti nei quali
si verifica ciascuna delle due disuguaglianze.
\end{example}

\begin{figure}[htbp]
  \centering
  \begin{tikzpicture}[unipd/diagram,x=2.25cm,y=1.28cm]
  \draw[->,unipd-muted] (-1.55,0) -- (1.6,0) node[right] {\(x\)};
  \draw[->,unipd-muted] (0,-2.75) -- (0,2.8) node[above] {\(y\)};
  \draw[unipd-primary,line width=1.05pt]
    plot[domain=-1.33:1.33,samples=100,smooth]
      (\x,{\x*\x*\x});
  \draw[unipd-accent,line width=1.05pt]
    plot[domain=-1.33:1.33,samples=100,smooth]
      (\x,{-\x*\x*\x});
  \fill[unipd-muted] (0,0) circle (1.7pt);
  \node[anchor=west,text=unipd-primary] at (1.03,1.78) {\(f(x)=x^3\)};
  \node[anchor=west,text=unipd-accent] at (1.03,-1.78) {\(g(x)=-x^3\)};
  \node[below left] at (0,0) {\(0\)};
\end{tikzpicture}

  \caption{Le funzioni \(x^3\) e \(-x^3\) hanno lo
  stesso limite per \(x\to0\), ma il loro ordine si
  inverte attraversando l'origine.}
  \label{fig:lez12-limiti-uguali-ordine}
\end{figure}

\section{Secondo teorema del confronto}
\label{sec:lez12-confronto-secondo}

\begin{theorem}[Secondo teorema del confronto]
\label{thm:lez12-confronto-secondo}
Siano \(f,g\colon A\longrightarrow\R\), con
\(A\subseteq\R\) non vuoto, e sia
\(x_0\in\overline{\R}\) un punto di accumulazione per \(A\).
Si supponga che esista un intorno \(U_{x_0}\) di
\(x_0\) tale che
\[
  f(x)\le g(x)
  \qquad\forall x\in(U_{x_0}\cap A)\setminus\{x_0\}.
\]
Allora:
\begin{enumerate}
  \item se
  \[
    \lim_{x\to x_0}f(x)=\ell_1\in\R,
    \qquad
    \lim_{x\to x_0}g(x)=\ell_2\in\R,
  \]
  si ha \(\ell_1\le\ell_2\);
  \item se \(\displaystyle\lim_{x\to x_0}g(x)=-\infty\),
  allora \(\displaystyle\lim_{x\to x_0}f(x)=-\infty\);
  \item se \(\displaystyle\lim_{x\to x_0}f(x)=+\infty\),
  allora \(\displaystyle\lim_{x\to x_0}g(x)=+\infty\).
\end{enumerate}
\end{theorem}

La dimostrazione del \cref{thm:lez12-confronto-secondo}
non viene sviluppata in questa sede.

La disuguaglianza stretta \(f(x)<g(x)\)
non garantisce \(\ell_1<\ell_2\).

\begin{example}[Disuguaglianza stretta e limiti uguali]
\label{ex:lez12-confronto-stretto-limiti-uguali}
Si considerino le funzioni \(f,g\colon\R\longrightarrow\R\)
definite da
\[
  f(x):=0,
  \qquad
  g(x):=
  \begin{cases}
    \dfrac1x, & x>0,\\[4pt]
    0, & x=0,\\[4pt]
    -\dfrac1x, & x<0.
  \end{cases}
\]
Si ha \(f(x)<g(x)\) per ogni \(x\ne0\). Inoltre,
per ogni \(\varepsilon>0\), scegliendo \(b=1/\varepsilon\),
per \(x>b\) si ottiene \(0<g(x)=1/x<\varepsilon\).
Pertanto
\[
  \lim_{x\to+\infty}f(x)
  =\lim_{x\to+\infty}g(x)=0.
\]
\end{example}

\begin{figure}[htbp]
  \centering
  \begin{tikzpicture}[unipd/diagram,x=.86cm,y=1.05cm]
  \draw[->,unipd-muted] (-5.1,0) -- (5.35,0) node[right] {\(x\)};
  \draw[->,unipd-muted] (0,-.3) -- (0,3.45) node[above] {\(y\)};
  \draw[unipd-accent,line width=1.1pt] (-4.95,0) -- (5.1,0);
  \draw[unipd-primary,line width=1pt]
    plot[domain=-4.9:-.36,samples=120,smooth]
      (\x,{-1/\x});
  \draw[unipd-primary,line width=1pt]
    plot[domain=.36:4.9,samples=120,smooth]
      (\x,{1/\x});
  \fill[unipd-primary] (0,0) circle (1.8pt);
  \node[text=unipd-primary,anchor=west] at (1.25,1.45)
    {\(g(x)=1/|x|\quad(x\ne0)\)};
  \node[text=unipd-accent,anchor=west] at (2.45,-.23)
    {\(f(x)=0\)};
  \node[below left] at (0,0) {\(0\)};
\end{tikzpicture}

  \caption{La funzione \(g(x)>f(x)=0\) per ogni
  \(x\ne0\), ma entrambe hanno limite \(0\) per
  \(x\to+\infty\).}
  \label{fig:lez12-stretto-limiti-uguali}
\end{figure}

\clearpage

\begin{example}[Un limite infinito mediante confronto]
\label{ex:lez12-limite-x-seno}
Si consideri la funzione
\[
  g\colon\R\longrightarrow\R,
  \qquad g(x)=x+\sin x.
\]
Poiché \(\sin x\ge-1\) per ogni \(x\in\R\), si ha
\[
  x-1\le x+\sin x
  \qquad\forall x\in\R.
\]
Per \(f(x):=x-1\) si ha
\[
  \lim_{x\to+\infty}f(x)=+\infty,
\]
poiché, per ogni \(K\in\R\), la condizione \(x>K+1\)
implica \(f(x)>K\). Per il terzo punto del
\cref{thm:lez12-confronto-secondo}, segue
\[
  \lim_{x\to+\infty}(x+\sin x)=+\infty.
\]
\end{example}

\begin{figure}[htbp]
  \centering
  \begin{tikzpicture}[unipd/diagram,x=.92cm,y=.62cm]
  \draw[->,unipd-muted] (-.35,0) -- (7.2,0) node[right] {\(x\)};
  \draw[->,unipd-muted] (0,-1.25) -- (0,8.25) node[above] {\(y\)};
  \draw[unipd-muted,densely dashed,line width=.75pt]
    plot[domain=0:6.25,samples=2] (\x,{\x+1});
  \draw[unipd-accent,line width=1pt]
    plot[domain=0:6.25,samples=2] (\x,{\x-1});
  \draw[unipd-primary,line width=1.15pt]
    plot[domain=0:6.25,samples=180,smooth]
      (\x,{\x+sin(deg(\x))});
  \node[font=\sffamily\footnotesize,text=unipd-muted,anchor=west]
    at (6.3,7.25) {\(x+1\)};
  \node[font=\sffamily\footnotesize,text=unipd-primary,anchor=west]
    at (6.3,6.22) {\(x+\sin x\)};
  \node[font=\sffamily\footnotesize,text=unipd-accent,anchor=west]
    at (6.3,5.25) {\(x-1\)};
  \node[below left] at (0,0) {\(0\)};
\end{tikzpicture}

  \caption{La funzione \(x+\sin x\) è compresa tra
  \(x-1\) e \(x+1\). In particolare, il minorante
  \(x-1\) tende a \(+\infty\).}
  \label{fig:lez12-confronto-seno}
\end{figure}

\section{Teorema delle tre funzioni}
\label{sec:lez12-tre-funzioni}

\begin{theorem}[Teorema delle tre funzioni]
\label{thm:lez12-tre-funzioni}
Siano \(f,g,h\colon A\longrightarrow\R\), con
\(A\subseteq\R\) non vuoto, e sia
\(x_0\in\overline{\R}\) un punto di accumulazione per \(A\).
Si supponga che
\[
  \lim_{x\to x_0}f(x)=\ell\in\R,
  \qquad
  \lim_{x\to x_0}g(x)=\ell,
\]
e che esista un intorno \(U_{x_0}\) di \(x_0\) tale che
\[
  f(x)\le h(x)\le g(x)
  \qquad\forall x\in(U_{x_0}\cap A)\setminus\{x_0\}.
\]
Allora
\[
  \lim_{x\to x_0}h(x)=\ell.
\]
\end{theorem}

\begin{proof}
Sia \(\varepsilon>0\). Per la \cref{def:lez08-limite},
esistono intorni \(U'_{x_0}\) e \(U''_{x_0}\) di
\(x_0\) tali che
\[
  |f(x)-\ell|<\varepsilon
  \qquad\forall x\in(U'_{x_0}\cap A)\setminus\{x_0\},
\]
\[
  |g(x)-\ell|<\varepsilon
  \qquad\forall x\in(U''_{x_0}\cap A)\setminus\{x_0\}.
\]
Per la \cref{sec:lez11-intersezione-intorni},
\[
  W_{x_0}:=U'_{x_0}\cap U''_{x_0}\cap U_{x_0}
\]
è ancora un intorno di \(x_0\). Per ogni
\(x\in(W_{x_0}\cap A)\setminus\{x_0\}\) valgono
simultaneamente le disuguaglianze
\[
  \ell-\varepsilon<f(x)\le h(x)
  \le g(x)<\ell+\varepsilon.
\]
Quindi
\[
  |h(x)-\ell|<\varepsilon
  \qquad\forall x\in(W_{x_0}\cap A)\setminus\{x_0\},
\]
e, per la \cref{def:lez08-limite}, si conclude
\(\lim_{x\to x_0}h(x)=\ell\).
\end{proof}

\begin{exercise}[Limite di \(x\sin(1/x)\)]
\label{ex:lez12-limite-seno-inverso}
Verificare mediante il \cref{thm:lez12-tre-funzioni} che
\[
  \lim_{x\to0}\left|x\sin\frac1x\right|=0,
  \qquad
  \lim_{x\to0}x\sin\frac1x=0.
\]
\end{exercise}

\begin{figure}[htbp]
  \centering
  \begin{tikzpicture}[unipd/diagram,x=4.5cm,y=2.9cm]
  \draw[->,unipd-muted] (-1.12,0) -- (1.15,0) node[right] {\(x\)};
  \draw[->,unipd-muted] (0,-1.08) -- (0,1.09) node[above] {\(y\)};
  \draw[unipd-accent,densely dashed,line width=.85pt]
    (-.95,.95) -- (0,0) -- (.95,.95);
  \draw[unipd-accent,densely dashed,line width=.85pt]
    (-.95,-.95) -- (0,0) -- (.95,-.95);
  \draw[unipd-primary,line width=1pt]
    plot[domain=-.95:-.032,samples=350,smooth]
      (\x,{\x*sin(deg(1/\x))});
  \draw[unipd-primary,line width=1pt]
    plot[domain=.032:.95,samples=350,smooth]
      (\x,{\x*sin(deg(1/\x))});
  \fill[white] (0,0) circle (2.2pt);
  \draw[unipd-primary,line width=.85pt] (0,0) circle (2.2pt);
  \node[font=\sffamily\footnotesize,text=unipd-accent,anchor=west]
    at (.61,.86) {\(|x|\)};
  \node[font=\sffamily\footnotesize,text=unipd-accent,anchor=west]
    at (.62,-.88) {\(-|x|\)};
\end{tikzpicture}

  \caption{La funzione \(x\sin(1/x)\), non definita in
  \(x=0\), rimane compresa tra \(-|x|\) e \(|x|\),
  entrambe convergenti a \(0\) per \(x\to0\).}
  \label{fig:lez12-seno-inverso}
\end{figure}

\begin{proof}[Soluzione dell'esercizio]
\emph{Dominio e punto di accumulazione.}
La funzione
\[
  h(x):=x\sin\frac1x
\]
ha dominio \(A=\R\setminus\{0\}\). Il punto \(0\) è di
accumulazione per \(A\): per ogni \(\delta>0\), il
numero \(a=\delta/2\) appartiene ad \(A\) e soddisfa
\[
  0<|a|=\frac\delta2<\delta.
\]

\emph{Applicazione del teorema.}
Per ogni \(x\in A\), poiché
\(\lvert\sin(1/x)\rvert\le1\), si ha
\[
  0\le\left|x\sin\frac1x\right|
  =|x|\left|\sin\frac1x\right|\le|x|.
\]
Per ogni \(\varepsilon>0\), scegliendo
\(\delta=\varepsilon\), si verifica
\(|x|<\varepsilon\) quando \(0<|x|<\delta\).
Pertanto
\[
  \lim_{x\to0}|x|=0,
  \qquad
  \lim_{x\to0}(-|x|)=0.
\]
Poiché anche la funzione costante nulla ha limite
\(0\), applicando il \cref{thm:lez12-tre-funzioni}
alla prima disuguaglianza si ottiene
\[
  \lim_{x\to0}\left|x\sin\frac1x\right|=0.
\]
Dalla medesima stima segue inoltre
\[
  -|x|\le x\sin\frac1x\le|x|
  \qquad\forall x\in A.
\]
Applicando nuovamente il \cref{thm:lez12-tre-funzioni},
si conclude
\[
  \lim_{x\to0}x\sin\frac1x=0.
\]
\end{proof}


\chapter[Lezione 13 --- 23 ottobre 2026]{Lezione 13: funzione di Dirichlet e proprietà dei limiti}
\label{chap:lezione-13}

\section{Funzione di Dirichlet}
\label{sec:lez13-dirichlet}

\begin{example}[Funzione di Dirichlet]
\label{ex:lez13-dirichlet}
Si consideri la funzione \(f\colon\R\longrightarrow\R\) definita da
\[
  f(x):=
  \begin{cases}
    1, & x\in\R\setminus\Q,\\
    -1, & x\in\Q.
  \end{cases}
\]
Si ha \(f(0)=-1\), ma \(f\) non ammette limite per \(x\to0\).
\end{example}

\begin{proof}
Il dominio è \(A=\R\). Per ogni \(\delta>0\), il punto
\(a=\delta/2\) appartiene ad \(A\) e soddisfa
\[
  0<|a|=\frac{\delta}{2}<\delta.
\]
Pertanto \(0\) è punto di accumulazione per \(A\).
Si supponga anzitutto che
\[
  \lim_{x\to0}f(x)=\ell\in\R.
\]
Per la \cref{def:lez08-limite}, scegliendo
\(\varepsilon=1/2\), esiste \(\delta>0\) tale che
\[
  |f(x)-\ell|<\frac12
  \qquad\forall x\in\R:\ 0<|x|<\delta.
\]
Per il \cref{thm:razionali-irrazionali-tra-reali}, esistono
\(r\in\Q\) e \(s\in\R\setminus\Q\) tali che
\(0<r,s<\delta\). Poiché \(f(r)=-1\) e \(f(s)=1\), si ottiene
\[
  |-1-\ell|<\frac12,
  \qquad
  |1-\ell|<\frac12.
\]
Per la disuguaglianza triangolare, ne segue
\[
  2=|(1-\ell)-(-1-\ell)|
  \le |1-\ell|+|-1-\ell|<1,
\]
una contraddizione.

Neppure \(+\infty\) e \(-\infty\) possono essere limiti:
poiché \(-1\le f(x)\le1\) per ogni \(x\in\R\),
la \cref{def:lez08-limite} non è soddisfatta scegliendo
rispettivamente la soglia \(K=2\) o \(K=-2\).
Pertanto \(f\) non ammette limite per \(x\to0\).
\end{proof}

\section{Proprietà elementari dei limiti}
\label{sec:lez13-proprieta-limiti}

\begin{proposition}[Proprietà elementari dei limiti]
\label{prop:lez13-proprieta-limiti}
Sia \(f\colon A\longrightarrow\R\), con \(A\subseteq\R\)
non vuoto, e sia \(x_0\in\overline{\R}\) un punto di
accumulazione per \(A\). Allora:
\begin{enumerate}
  \item per ogni \(\ell\in\R\),
  \[
    \lim_{x\to x_0}f(x)=\ell
    \quad\Longleftrightarrow\quad
    \lim_{x\to x_0}(f(x)-\ell)=0;
  \]
  \item
  \[
    \lim_{x\to x_0}f(x)=0
    \quad\Longleftrightarrow\quad
    \lim_{x\to x_0}|f(x)|=0;
  \]
  \item per ogni \(\ell\in\R\setminus\{0\}\),
  \[
    \lim_{x\to x_0}f(x)=\ell
    \quad\Longrightarrow\quad
    \lim_{x\to x_0}|f(x)|=|\ell|;
  \]
  \item
  \[
    \lim_{x\to x_0}f(x)=0
    \quad\Longleftrightarrow\quad
    \lim_{x\to x_0}f(x)^2=0.
  \]
\end{enumerate}
\end{proposition}

\begin{proof}
\emph{Primo punto.}
Per ogni \(x\in A\) e \(\ell\in\R\),
\[
  |(f(x)-\ell)-0|=|f(x)-\ell|.
\]
Le due condizioni della \cref{def:lez08-limite}
coincidono, quindi vale l'equivalenza.

\emph{Secondo punto.}
Per ogni \(x\in A\),
\[
  \bigl||f(x)|-0\bigr|=|f(x)-0|.
\]
Anche in questo caso le condizioni della
\cref{def:lez08-limite} coincidono.

\emph{Terzo punto.}
Si supponga \(f(x)\to\ell\in\R\setminus\{0\}\)
per \(x\to x_0\). Per la disuguaglianza triangolare inversa,
\[
  \bigl||f(x)|-|\ell|\bigr|\le|f(x)-\ell|.
\]
Dato \(\varepsilon>0\), per la \cref{def:lez08-limite}
esiste un intorno \(U_{x_0}\) tale che
\[
  |f(x)-\ell|<\varepsilon
  \qquad\forall x\in(U_{x_0}\cap A)\setminus\{x_0\}.
\]
Per gli stessi \(x\) si ha
\[
  \bigl||f(x)|-|\ell|\bigr|<\varepsilon,
\]
e quindi \(\lim_{x\to x_0}|f(x)|=|\ell|\).

\emph{Quarto punto.}
Si dimostrano separatamente le due implicazioni.

\emph{Prima implicazione \((\Rightarrow)\).}
Si supponga \(f(x)\to0\) per \(x\to x_0\).
Per ogni \(\varepsilon>0\), applicando la
\cref{def:lez08-limite} con \(\sqrt{\varepsilon}>0\),
esiste un intorno \(U_{x_0}\) tale che
\[
  |f(x)|<\sqrt{\varepsilon}
  \qquad\forall x\in(U_{x_0}\cap A)\setminus\{x_0\}.
\]
Ne segue
\[
  |f(x)^2|=|f(x)|^2<\varepsilon
  \qquad\forall x\in(U_{x_0}\cap A)\setminus\{x_0\}.
\]
Pertanto \(\lim_{x\to x_0}f(x)^2=0\).

\emph{Seconda implicazione \((\Leftarrow)\).}
Si supponga \(f(x)^2\to0\) per \(x\to x_0\).
Per ogni \(\varepsilon>0\), applicando la
\cref{def:lez08-limite} con \(\varepsilon^2>0\),
esiste un intorno \(U_{x_0}\) tale che
\[
  |f(x)^2|<\varepsilon^2
  \qquad\forall x\in(U_{x_0}\cap A)\setminus\{x_0\}.
\]
Poiché \(|f(x)^2|=|f(x)|^2\), si ottiene
\[
  |f(x)|<\varepsilon
  \qquad\forall x\in(U_{x_0}\cap A)\setminus\{x_0\}.
\]
Pertanto \(\lim_{x\to x_0}f(x)=0\).
\end{proof}


\courseappendices
\chapter[Tutorato 1 --- 8 ottobre 2026]{Tutorato 1: fattoriale e dimostrazioni per induzione}
\label[appendix]{chap:tutorato-01}

\begin{exercise}[Semplificazione di un rapporto di fattoriali]
\label{ex:tutorato01-rapporto-fattoriali}
Per \(n\in\N\), semplificare l'espressione
\[
  \frac{(n+2)!+(n+1)!}{(n+3)!}.
\]
\end{exercise}

\begin{proof}[Soluzione dell'esercizio]
Per la \cref{def:fattoriale},
\[
  (n+2)!=(n+2)(n+1)!,
  \qquad
  (n+3)!=(n+3)(n+2)(n+1)!.
\]
Raccogliendo \((n+1)!\) al numeratore, si ottiene
\begin{align*}
  \frac{(n+2)!+(n+1)!}{(n+3)!}
    &=\frac{\bigl((n+2)+1\bigr)(n+1)!}
      {(n+3)(n+2)(n+1)!}\\
    &=\frac{(n+3)(n+1)!}{(n+3)(n+2)(n+1)!}\\
    &=\frac{1}{n+2}.
\end{align*}
Il denominatore iniziale e i fattori semplificati sono strettamente
positivi per ogni \(n\in\N\); tutte le divisioni sono quindi lecite,
anche nel caso \(n=0\).
\end{proof}

\begin{exercise}[Somma dei primi \(n\) numeri dispari positivi]
\label{ex:tutorato01-somma-dispari}
Dimostrare che, per ogni \(n\in\N\) con \(n\ge1\), vale
\[
  \sum_{k=1}^{n}(2k-1)=n^2.
\]
\end{exercise}

\begin{proof}[Soluzione dell'esercizio]
Si applica il \cref{thm:principio-induzione} con valore iniziale \(p=1\).

\emph{Passo base.}
Per \(n=1\),
\[
  \sum_{k=1}^{1}(2k-1)=1=1^2.
\]

\emph{Passo induttivo.}
Sia \(n\ge1\) e si supponga che valga l'ipotesi induttiva
\[
  \sum_{k=1}^{n}(2k-1)=n^2.
\]
Separando l'ultimo addendo della somma fino a \(n+1\) e applicando
l'ipotesi induttiva, si ottiene
\begin{align*}
  \sum_{k=1}^{n+1}(2k-1)
    &=\sum_{k=1}^{n}(2k-1)+\bigl(2(n+1)-1\bigr)\\
    &=n^2+2n+1\\
    &=(n+1)^2.
\end{align*}
È dunque verificata la formula per \(n+1\). Per il principio di
induzione, l'identità vale per ogni \(n\ge1\).
\end{proof}

\begin{exercise}[Somma dei quadrati dei primi \(n\) naturali positivi]
\label{ex:tutorato01-somma-quadrati}
Dimostrare che, per ogni \(n\in\N\) con \(n\ge1\), vale
\[
  \sum_{k=1}^{n}k^2=\frac{n(n+1)(2n+1)}6.
\]
\end{exercise}

\begin{proof}[Soluzione dell'esercizio]
Si applica il \cref{thm:principio-induzione} con valore iniziale \(p=1\).

\emph{Passo base.}
Per \(n=1\),
\[
  \sum_{k=1}^{1}k^2=1=\frac{1\cdot2\cdot3}{6}.
\]

\emph{Passo induttivo.}
Sia \(n\ge1\) e si supponga
\[
  \sum_{k=1}^{n}k^2=\frac{n(n+1)(2n+1)}6.
\]
Separando l'ultimo addendo e applicando l'ipotesi induttiva,
\begin{align*}
  \sum_{k=1}^{n+1}k^2
    &=\sum_{k=1}^{n}k^2+(n+1)^2\\
    &=\frac{n(n+1)(2n+1)}6+(n+1)^2\\
    &=\frac{n+1}{6}\bigl(n(2n+1)+6(n+1)\bigr)\\
    &=\frac{n+1}{6}(2n^2+7n+6)\\
    &=\frac{(n+1)(n+2)(2n+3)}6.
\end{align*}
Nell'ultimo passaggio è stata utilizzata la fattorizzazione
\(2n^2+7n+6=(n+2)(2n+3)\). Poiché \(2n+3=2(n+1)+1\),
l'espressione ottenuta coincide con quella richiesta per \(n+1\).
Il passo induttivo è quindi verificato, e la formula vale per ogni
\(n\ge1\).
\end{proof}

\begin{exercise}[Dimostrazione del binomio di Newton]
\label{ex:tutorato01-binomio-newton}
Dimostrare per induzione il \cref{thm:binomio-newton}: per ogni
\(a,b\in\R\) e per ogni \(n\in\N\),
\[
  (a+b)^n=\sum_{k=0}^{n}\binom{n}{k}a^k b^{n-k}.
\]
\end{exercise}

\begin{proof}[Soluzione dell'esercizio]
Fissati \(a,b\in\R\), si applica il \cref{thm:principio-induzione}
con valore iniziale \(p=0\). Si mantiene la convenzione \(t^0=1\),
anche per \(t=0\), adottata nel \cref{thm:binomio-newton}.

\emph{Passo base.}
Per la \cref{def:coefficiente-binomiale}, \(\binom{0}{0}=1\).
Quindi, per \(n=0\),
\[
  (a+b)^0=1=\binom{0}{0}a^0b^0.
\]

\emph{Passo induttivo.}
Sia \(n\in\N\) e si supponga che valga l'ipotesi induttiva
\[
  (a+b)^n=\sum_{k=0}^{n}\binom{n}{k}a^k b^{n-k}.
\]
Moltiplicando per \(a+b\) e distribuendo il prodotto rispetto alla
somma finita, si ottiene
\begin{align*}
  (a+b)^{n+1}
    &=(a+b)^n(a+b)\\
    &=\left(\sum_{k=0}^{n}\binom{n}{k}a^k b^{n-k}\right)(a+b)\\
    &=\sum_{k=0}^{n}\binom{n}{k}a^{k+1}b^{n-k}
      +\sum_{k=0}^{n}\binom{n}{k}a^k b^{n+1-k}.
\end{align*}
Nella prima somma si pone \(j=k+1\). Quando \(k\) varia da \(0\)
a \(n\), il nuovo indice \(j\) varia da \(1\) a \(n+1\);
inoltre \(k=j-1\) e \(n-k=n+1-j\). Pertanto
\[
  \sum_{k=0}^{n}\binom{n}{k}a^{k+1}b^{n-k}
  =\sum_{j=1}^{n+1}\binom{n}{j-1}a^j b^{n+1-j}.
\]
Rinominando \(j\) come \(k\), l'uguaglianza precedente dà
\[
  (a+b)^{n+1}
  =\sum_{k=1}^{n+1}\binom{n}{k-1}a^k b^{n+1-k}
    +\sum_{k=0}^{n}\binom{n}{k}a^k b^{n+1-k}.
\]
Le due somme hanno ora gli stessi fattori \(a^k b^{n+1-k}\),
ma intervalli di sommazione diversi: gli indici comuni sono
\(1,\ldots,n\). Il termine con \(k=n+1\) compare soltanto nella
prima somma e vale \(a^{n+1}\), mentre quello con \(k=0\) compare
soltanto nella seconda e vale \(b^{n+1}\), poiché
\(\binom{n}{n}=\binom{n}{0}=1\).
Separando questi termini e combinando le somme sugli indici comuni,
si ottiene
\[
  (a+b)^{n+1}
  =b^{n+1}
    +\sum_{k=1}^{n}
      \left(\binom{n}{k-1}+\binom{n}{k}\right)a^k b^{n+1-k}
    +a^{n+1}.
\]
Per \(n=0\) non vi sono indici comuni: la somma da \(1\) a \(0\)
è vuota e vale \(0\), e l'uguaglianza resta valida.

Per ogni \(1\le k\le n\), l'identità di Pascal, dimostrata nella
\cref{prop:identita-pascal}, si applica con indice superiore \(n+1\)
e fornisce
\[
  \binom{n}{k-1}+\binom{n}{k}=\binom{n+1}{k}.
\]
Infatti, \(n+1\ge1\) e \(1\le k\le n+1\), come richiesto
dalle ipotesi dell'identità. Per la
\cref{def:coefficiente-binomiale}, anche i termini estremi hanno
coefficienti \(\binom{n+1}{0}=\binom{n+1}{n+1}=1\). Ne segue
\begin{align*}
  (a+b)^{n+1}
    &=b^{n+1}+\sum_{k=1}^{n}\binom{n+1}{k}a^k b^{n+1-k}
      +a^{n+1}\\
    &=\sum_{k=0}^{n+1}\binom{n+1}{k}a^k b^{n+1-k}.
\end{align*}
È quindi verificata la formula per \(n+1\). Il principio di induzione
assicura che essa vale per ogni \(n\in\N\). Poiché \(a,b\in\R\)
erano arbitrari, segue il \cref{thm:binomio-newton}.
\end{proof}

\chapter[Tutorato 2 --- 22 ottobre 2026]{Tutorato 2: estremi di insiemi e verifica di un limite}
\label[appendix]{chap:tutorato-02}

\begin{exercise}[Estremi di un insieme definito da una disequazione]
\label{ex:disequazione-razionale}
Sia
\[
  A:=\left\{x\in\R:\frac{2x-1}{x+2}<1\right\}.
\]
Determinare l'insieme \(A\), i suoi estremi superiore e inferiore
ed eventuali massimo e minimo.
\end{exercise}

\begin{proof}[Soluzione dell'esercizio]
Per \(x\neq-2\), la disequazione equivale a
\[
  \frac{2x-1}{x+2}<1
  \quad\Longleftrightarrow\quad
  \frac{x-3}{x+2}<0.
\]
Il rapporto è negativo se numeratore e denominatore hanno segni
opposti, ossia per \(-2<x<3\). Pertanto
\[
  A=(-2,3).
\]
Per ogni \(\varepsilon>0\), posto
\[
  t:=\min\left\{\frac{\varepsilon}{2},1\right\}>0,
\]
i punti \(-2+t\) e \(3-t\) appartengono ad \(A\) e verificano
\[
  -2<-2+t<-2+\varepsilon,
  \qquad
  3-\varepsilon<3-t<3.
\]
Per i \cref{thm:caratterizzazione-inf,thm:caratterizzazione-sup},
si ha
\[
  \inf A=-2,\qquad \sup A=3.
\]
Poiché \(-2,3\notin A\), per la \cref{def:massimo-minimo}
non esistono né \(\min A\) né \(\max A\).
\end{proof}

\begin{figure}[htbp]
  \centering
  \begin{tikzpicture}[unipd/diagram]
  \draw[->,unipd-muted] (-3.25,0) -- (4.2,0) node[right] {\(\R\)};
  \draw[very thick,unipd-primary] (-2,0) -- (3,0);
  \draw[thick,unipd-primary,fill=white] (-2,0) circle (2.5pt);
  \draw[thick,unipd-primary,fill=white] (3,0) circle (2.5pt);
  \node[below=7pt] at (-2,0) {\(-2\)};
  \node[below=7pt] at (3,0) {\(3\)};
  \node[above=9pt] at (.5,0) {\(A=(-2,3)\)};
\end{tikzpicture}

  \caption{L'insieme \(A=(-2,3)\): gli estremi \(-2\) e \(3\)
  non appartengono all'intervallo.}
  \label{fig:intervallo-aperto}
\end{figure}

\begin{exercise}[Estremi di un insieme definito da termini indicizzati]
\label{ex:estremi-termini}
Sia
\[
  B:=\left\{1-\frac{(-1)^n}{n}:n\in\N,\ n\ge1\right\}.
\]
Determinare gli estremi superiore e inferiore di \(B\) e
stabilire se sono rispettivamente massimo e minimo.
\end{exercise}

\begin{proof}[Soluzione dell'esercizio]
Posto
\[
  b_n:=1-\frac{(-1)^n}{n},\qquad n\ge1,
\]
i primi quattro termini sono
\[
  b_1=2,\quad b_2=\frac12,\quad
  b_3=\frac43,\quad b_4=\frac34.
\]
Si distinguono i due casi, secondo che \(n\) sia pari o dispari:
\[
  b_n=\begin{cases}
    1+\dfrac1n,&n\text{ dispari},\\[4pt]
    1-\dfrac1n,&n\text{ pari}.
  \end{cases}
\]

\emph{Indici dispari.}
Per ogni \(n\ge1\) dispari,
\[
  b_{n+2}=1+\frac1{n+2}<1+\frac1n=b_n.
\]
I termini dispari sono quindi decrescenti e soddisfano
\[
  1<b_n\le b_1=2.
\]

\emph{Indici pari.}
Per ogni \(n\ge2\) pari,
\[
  b_{n+2}=1-\frac1{n+2}>1-\frac1n=b_n.
\]
I termini pari sono quindi crescenti e soddisfano
\[
  \frac12=b_2\le b_n<1.
\]
In entrambi i casi, \(\lvert b_n-1\rvert=1/n\), quindi i termini si
avvicinano a \(1\) al crescere di \(n\). Inoltre,
\[
  \frac12\le b_n\le2\qquad\forall n\ge1.
\]
Poiché \(b_1=2\) e \(b_2=\frac12\), per la
\cref{def:massimo-minimo} si conclude
\[
  \sup B=\max B=2,\qquad
  \inf B=\min B=\frac12.
\]
\end{proof}

\begin{figure}[htbp]
  \centering
  \begin{tikzpicture}[unipd/diagram,x=.75cm,y=2.6cm]
  \draw[->,unipd-muted] (0,.92) -- (10.4,.92) node[right] {\(n\)};
  \draw[->,unipd-muted] (.25,.92) -- (.25,2.18) node[above] {\(b_n\)};
  \draw[dashed,unipd-muted] (.25,1) -- (10,1);
  \node[left=5pt] at (.25,1) {\(1\)};
  \foreach \n in {1,3,5,7,9} {
    \pgfmathsetmacro\v{1+1/\n}
    \fill[unipd-primary] (\n,\v) circle (2pt);
    \draw[unipd-muted] (\n,.9) -- (\n,.94);
    \node[below=5pt] at (\n,.92) {\(\n\)};
  }
  \node[above right=2pt,text=unipd-primary] at (1,2) {\(b_1=2\)};
  \node[above right=2pt,text=unipd-primary] at (3,1.33333) {\(b_3=\frac43\)};
\end{tikzpicture}

  \caption{Termini di indice dispari di \(b_n\), decrescenti
  a partire da \(b_1=2\).}
  \label{fig:termini-dispari}
\end{figure}

\begin{figure}[htbp]
  \centering
  \begin{tikzpicture}[unipd/diagram,x=.75cm,y=3.5cm]
  \draw[->,unipd-muted] (0,.37) -- (11.2,.37) node[right] {\(n\)};
  \draw[->,unipd-muted] (.25,.37) -- (.25,1.13) node[above] {\(b_n\)};
  \draw[dashed,unipd-muted] (.25,1) -- (10.7,1);
  \node[left=5pt] at (.25,1) {\(1\)};
  \foreach \n in {2,4,6,8,10} {
    \pgfmathsetmacro\v{1-1/\n}
    \fill[unipd-primary] (\n,\v) circle (2pt);
    \draw[unipd-muted] (\n,.355) -- (\n,.385);
    \node[below=5pt] at (\n,.37) {\(\n\)};
  }
  \node[above right=2pt,text=unipd-primary] at (2,.5) {\(b_2=\frac12\)};
  \node[below right=2pt,text=unipd-primary] at (4,.75) {\(b_4=\frac34\)};
\end{tikzpicture}

  \caption{Termini di indice pari di \(b_n\), crescenti
  a partire da \(b_2=\frac12\).}
  \label{fig:termini-pari}
\end{figure}

\clearpage

\begin{exercise}[Verifica di un limite finito]
\label{ex:limite-lineare}
Verificare mediante la \cref{def:lez08-limite} che
\[
  \lim_{x\to2}(4x-3)=5.
\]
\end{exercise}

\begin{proof}[Soluzione dell'esercizio]
\emph{Dominio e punto di accumulazione.}
La funzione \(f\colon\R\longrightarrow\R\), definita da
\[
  f(x):=4x-3,
\]
ha dominio \(A=\R\).

Il punto \(2\) è di accumulazione per \(A\). Infatti, per ogni
\(\delta>0\), il numero \(a=2+\delta/2\) appartiene ad \(A\)
e soddisfa
\[
  0<|a-2|=\frac{\delta}{2}<\delta.
\]

\emph{Verifica della definizione.}
Fissato \(\varepsilon>0\), occorre trovare \(\delta_\varepsilon>0\)
tale che
\[
  |f(x)-5|<\varepsilon
  \qquad\forall x\in A:
  0<|x-2|<\delta_\varepsilon.
\]
Per ogni \(x\in A\), si ha
\[
  |f(x)-5|=|4x-8|=4|x-2|.
\]
È pertanto sufficiente scegliere
\[
  \delta_\varepsilon:=\frac{\varepsilon}{4}>0.
\]
Infatti, per ogni \(x\in A\) tale che
\(0<|x-2|<\delta_\varepsilon\), si ottiene
\[
  |f(x)-5|=4|x-2|
  <4\delta_\varepsilon=\varepsilon.
\]
Poiché \(\varepsilon>0\) era arbitrario, per la
\cref{def:lez08-limite} si conclude
\[
  \lim_{x\to2}(4x-3)=5.
\]
\end{proof}
\begin{figure}[htbp]
  \centering
  \begin{tikzpicture}[unipd/diagram,x=2.25cm,y=.57cm]
  \fill[unipd-soft] (1.5,3) rectangle (2.5,7);
  \draw[->,unipd-muted] (.35,0) -- (3.15,0) node[right] {\(x\)};
  \draw[->,unipd-muted] (.5,0) -- (.5,8.6) node[above] {\(y\)};
  \foreach \xx in {1.5,2,2.5} {
    \draw[dashed,unipd-muted] (\xx,0) -- (\xx,7.25);
  }
  \foreach \yy in {3,5,7} {
    \draw[densely dotted,unipd-muted] (.5,\yy) -- (2.82,\yy);
  }
  \draw[very thick,unipd-primary,domain=1.18:2.83,samples=2]
    plot (\x,{4*\x-3});
  \fill[unipd-primary] (2,5) circle (2pt);
  \node[below=4pt] at (1.5,0) {\(2-\delta\)};
  \node[below=4pt] at (2,0) {\(2\)};
  \node[below=4pt] at (2.5,0) {\(2+\delta\)};
  \node[left=5pt] at (.5,3) {\(5-\varepsilon\)};
  \node[left=5pt] at (.5,5) {\(5\)};
  \node[left=5pt] at (.5,7) {\(5+\varepsilon\)};
  \node[above left=3pt,text=unipd-primary] at (2.72,7.88) {\(f(x)=4x-3\)};
\end{tikzpicture}

  \caption{Verifica di \(\lim_{x\to2}(4x-3)=5\):
  per \(\delta=\varepsilon/4\), la condizione
  \(0<|x-2|<\delta\) implica \(|f(x)-5|<\varepsilon\).}
  \label{fig:limite-lineare}
\end{figure}


\glsaddall

\end{document}
